\chapter{积分作为扩张机制：极限、密度与函数空间}\label{ch:03}

上一章已经能构造和传输测度，却还没有回答最朴素的计算问题：一个可测函数究竟怎样积？
对非负简单函数，答案只是有限个“高度乘质量”的和；一旦函数由无限层逼近，有限可加已经不够，
极限与积分能否交换便成为全部理论的关节。单调收敛、Fatou 引理和支配收敛定理正是在这一压力下
出现的，而不是三条彼此无关的极限口诀。

建立标量积分以后，我们会连续遇到四次同一种失败。第一，重复积分若只按形式交换次序，条件收敛
会给出不同答案；Tonelli 与 Fubini 必须分开处理非负性和绝对可积性。第二，全局测度未必由点值
公式给出；Radon--Nikodym 定理要从集合上的绝对连续性恢复密度，并把剩余质量分离到零测骨架上。
第三，范数 Cauchy 的函数列不一定逐点 Cauchy；$L^p$ 的完备性要从一条可和的快速子列重建极限。
第四，向量值函数逐坐标可测仍可能没有简单函数逼近；Pettis 定理和 Bochner 积分必须把可分值域
与范数可积性同时纳入。

本章的主线因而是“先在有限层上计算，再用受控极限扩张”。它从简单函数一路延伸到乘积积分、
换元、核、$L^p$ 几何、Hahn--Banach、对偶、一致可积以及 Banach 值积分。每次扩张都会明确说明
什么条件阻止质量集中、逃逸或相消。条件期望由 Radon--Nikodym 表示自然引出；正则条件分布与
Banach 值 Radon--Nikodym 性质则标明从标量密度扩张到核或向量值密度时所需的额外结构。

\section{积分的构造与极限定理}\label{sec:03-01}

\subsection{从简单函数到正积分：Daniell 视角}\label{subsec:3-1-1}
\index{简单函数}\index{非负积分}\index{Daniell--Young 积分}

测度已经回答了集合有多大，却还没有告诉我们函数图像有多少 ``质量''。在有限原子空间里，
答案没有悬念。若
\[
  \mu=\sum_{j=1}^N w_j\delta_{x_j},\qquad w_j\geq0,
\]
那么任何合理的积分都应复现加权和
\[
  \int f\,\dd\mu=\sum_{j=1}^Nw_jf(x_j).
  \tag{3.1.1}\label{eq:3-1-finite-observation}
\]
真正的问题是：当函数有无穷多个取值、空间也不再由有限个原子组成时，怎样从有限和走过去？

最先能够逐项计算的函数不是连续函数，而是只取有限多个值的函数。由集合到函数的扩张因而从这里开始。
这条路线与 Daniell--Young 的观点平行：先把 ``正线性'' 和 ``从下逼近'' 当成不可让步的规则，
再让极限定理说明所得对象确实稳定。这里从已有测度出发构造积分；从正泛函反向恢复测度则是
Daniell 表示定理所处理的逆问题。

\begin{definition}[非负简单函数的积分]\label{def:3-1-1-1}
设 $(X,\mathcal A,\mu)$ 为测度空间。沿用定义~\ref{def:2-2-1-simple-function}，把非负简单函数写成
\[
  s=\sum_{k=1}^m a_k\1_{E_k},
\]
其中 $a_k\in[0,\infty)$，$E_k\in\mathcal A$ 两两不交，并定义
\[
  \int_Xs\,\dd\mu:=\sum_{k=1}^ma_k\mu(E_k).
\]
这里以及下文均约定 $0\cdot\infty=0$。
\end{definition}

同一简单函数可以有许多写法。例如，值为零的集合可以任意拆分，两个取同一系数的集合也可以合并。
因此首先要确认上式算的是函数，而不是写法。

\begin{lemma}[简单函数积分表示无关]\label{lemma:3-1-1-1}
简单函数的积分不依赖于所取的不交值层表示。对非负简单函数 $s,t$ 与 $a,b\geq0$，还有
\[
  \int(as+bt)\,\dd\mu=a\int s\,\dd\mu+b\int t\,\dd\mu.
\]
\end{lemma}

\begin{proof}
设
\[
 s=\sum_{i=1}^m a_i\1_{E_i}=\sum_{j=1}^n b_j\1_{F_j}
\]
是两个两两不交表示。加入余集并令其系数为零后，可设两族各自分割 $X$。集合
$E_i\cap F_j$ 构成公共有限分割。若 $E_i\cap F_j\ne\varnothing$，在其中任取一点便有
$a_i=s(x)=b_j$；若交集为空，其测度为零。因此
\[
 \sum_i a_i\mu(E_i)
 =\sum_{i,j}a_i\mu(E_i\cap F_j)
 =\sum_{i,j}b_j\mu(E_i\cap F_j)
 =\sum_jb_j\mu(F_j).
\]
等式只涉及非负扩展实数，因而没有无穷减无穷的问题。对 $s,t$ 的分割取公共细分，
在每个小块上直接合并系数，即得所述正线性。
\end{proof}

\begin{example}[三个有限观测模型]\label{ex:3-1-1-1}
若 $\mu$ 是计数测度，则对非负简单函数 $s$，
\[
 \int s\,\dd\mu=\sum_{x\in X}s(x),
\]
其中右侧定义为所有有限子集和的上确界；由于 $s$ 只有有限多个正值层，这与上面的定义一致。
若 $\mu=\delta_x$，则 $\int s\,\dd\delta_x=s(x)$。更一般地，有限原子测度立即给出
\eqref{eq:3-1-finite-observation}。这些计算提示我们：对非负量，添加更多观测只能使答案增加，
所以下确界逼近不如从下逼近自然。
\end{example}

\begin{example}[Dirac 测度]\label{ex:3-1-1-2}
上例的单原子情形值得单独记住。对任意非负可测函数 $f$，稍后的定义将给出
\[
 \int f\,\dd\delta_x=f(x),
\]
包括 $f(x)=+\infty$ 的情形。换言之，Dirac 积分就是一次点值观测。
\end{example}

\begin{definition}[非负积分]\label{def:3-1-1-2}
对可测函数 $f:X\to[0,+\infty]$，定义
\[
 \int_X f\,\dd\mu
 :=\sup\left\{\int_Xs\,\dd\mu:
      s\text{ 为非负简单函数且 }0\leq s\leq f\right\}.
 \tag{3.1.2}\label{eq:3-1-positive-integral}
\]
对 $E\in\mathcal A$，记
\[
 \int_Ef\,\dd\mu:=\int_Xf\1_E\,\dd\mu.
\]
\end{definition}

定义中的集合非空，因为零函数总是候选者；其上确界允许为 $+\infty$。经验测度
$n^{-1}\sum_{j=1}^n\delta_{x_j}$ 下，$\int_Ef$ 最终将等于只对满足 $x_j\in E$ 的样本求平均，
所以集合下标只是对有限观测作筛选，并未偷偷引入另一种积分。

\begin{definition}[几乎处处与共同零集]\label{def:3-1-1-ae}
若存在可测零集 $N$，使性质 $P(x)$ 对每个 $x\in X\setminus N$ 成立，就称 $P$ 在
$\mu$-几乎处处意义下成立，简记为 $\mu$-a.e. 成立。数列 $f_n\to f$ a.e. 时，
要求有一个对所有 $n$ 共同有效的可测零集。可数多个零集之并仍为零集，正是数列可以共同化异常集的原因。
\end{definition}

定义~\ref{def:3-1-1-2} 看似要在一个巨大的函数族上取上确界。事实上，有一列完全明确的阶梯就够了。

\begin{lemma}[规范层级逼近]\label{lemma:3-1-1-layer}
对非负可测函数 $f$，令
\[
 s_n(f)=2^{-n}\bigl\lfloor 2^n(f\wedge n)\bigr\rfloor,
 \qquad n\geq1,
 \tag{3.1.3}\label{eq:3-1-canonical-simple}
\]
并约定 $\lfloor+\infty\rfloor=+\infty$，但由于先取了 $f\wedge n$，右侧始终有限。
则 $s_n(f)$ 是有限值非负简单函数，
\[
  s_n(f)\uparrow f,
  \qquad
  \int f\,\dd\mu=\sup_n\int s_n(f)\,\dd\mu.
  \tag{3.1.4}\label{eq:3-1-canonical-sup}
\]
此外，若简单函数 $t_n\uparrow t$，且 $t$ 也是有限值简单函数，则
$\int t_n\uparrow\int t$。
\end{lemma}

\begin{proof}
函数 $s_n(f)$ 的值属于有限集合
$\{0,2^{-n},\ldots,n\}$，其各个值层由 $f$ 的可测逆像构成，故它是简单函数。
对任意 $u\in[0,+\infty]$，先逐情形核对单调性。若 $u\le n$，则
\[
 \lfloor2^{n+1}u\rfloor=\lfloor2(2^nu)\rfloor
 \ge 2\lfloor2^nu\rfloor,
\]
从而
\[
 2^{-(n+1)}\lfloor2^{n+1}(u\wedge(n+1))\rfloor
 \ge 2^{-n}\lfloor2^n(u\wedge n)\rfloor.
\]
若 $u>n$，则
\[
 2^{-(n+1)}\lfloor2^{n+1}(u\wedge(n+1))\rfloor
 \ge 2^{-(n+1)}\lfloor2^{n+1}n\rfloor=n
 =2^{-n}\lfloor2^n(u\wedge n)\rfloor,
\]
所以同一不等式仍成立。再若 $u<\infty$，当 $n>u$ 时
\[
 0\le u-2^{-n}\lfloor2^nu\rfloor<2^{-n},
\]
故该数列趋于 $u$；若 $u=+\infty$，则其第 $n$ 项正好为 $n$。逐点应用便得
$s_n(f)\uparrow f$。

由定义，$\sup_n\int s_n(f)\leq\int f$。反过来，任取非负简单函数 $s\leq f$ 和
$0<c<1$。若 $a>0$ 遍历 $s$ 的有限多个正值，取 $n$ 足够大，使
$n\geq\max a$ 且 $2^{-n}\leq(1-c)\min a$。在 $s(x)=a$ 的值层上，
\[
 s_n(f)(x)\geq2^{-n}\lfloor2^na\rfloor\geq a-2^{-n}\geq ca,
\]
而在零值层上不等式也成立。因此 $cs\leq s_n(f)$，从而
\[
 c\int s\,\dd\mu\leq\sup_n\int s_n(f)\,\dd\mu.
\]
先令 $c\uparrow1$，再对 $s$ 取上确界，得到 \eqref{eq:3-1-canonical-sup}。
这里没有使用任何积分极限定理。

最后设 $t=\sum_{j=1}^r a_j\1_{E_j}$，其中 $a_j>0$ 且 $E_j$ 两两不交；零值层可略去。
固定 $0<c<1$，令 $E_{j,n}=E_j\cap\{t_n\geq ca_j\}$。由 $t_n\uparrow t$，
$E_{j,n}\uparrow E_j$。测度从下连续给出
\[
 \begin{aligned}
 \liminf_n\int t_n\,\dd\mu
 &\geq \lim_n\int c\sum_j a_j\1_{E_{j,n}}\,\dd\mu\\
 &=c\sum_j a_j\lim_n\mu(E_{j,n})
 =c\int t\,\dd\mu.
 \end{aligned}
\]
令 $c\uparrow1$；再结合 $t_n\leq t$ 所给的反向不等式，结论成立。
这一段只用了测度的从下连续性，而该性质已由可数可加性证明。
\end{proof}

\begin{example}[层级阶梯]\label{ex:3-1-1-3}
若 $f(x)=x$ 定义在 $[0,1]$ 上，则 $s_n(f)$ 在每个
$[k2^{-n},(k+1)2^{-n})$ 上取值 $k2^{-n}$。因此
\[
 \int_0^1s_n(f)\,\dd x
 =2^{-2n}\sum_{k=0}^{2^n-1}k
 =\frac12-2^{-n-1}.
\]
由 \eqref{eq:3-1-canonical-sup}，$\int_0^1x\,\dd x=1/2$。同一计算也说明，
定义中的上确界不是一件不可计算的装置：它把熟悉的下矩形和直接纳入其中。
\end{example}

\begin{definition}[正线性泛函]\label{def:3-1-1-3}
设 $\mathcal C$ 是一个非负函数锥。映射 $I:\mathcal C\to[0,+\infty]$ 若保持次序，且
\[
 I(af+bg)=aI(f)+bI(g)\qquad(a,b\geq0),
\]
则称为正线性泛函。扩展值等式仍采用 $0\cdot(+\infty)=0$，并且不作
$+\infty-(+\infty)$ 运算。
\end{definition}

下一命题说明，从下方简单函数取上确界确实保留了有限和的代数。正向不等式直接来自定义，
反向不等式则需要同时细分两个简单函数逼近。

\begin{proposition}[非负积分的代数性质]\label{proposition:3-1-1-2}
若 $f,g:X\to[0,+\infty]$ 可测且 $a,b\geq0$，则
\[
 \int(af+bg)\,\dd\mu
 =a\int f\,\dd\mu+b\int g\,\dd\mu.
 \tag{3.1.5}\label{eq:3-1-positive-linearity}
\]
若 $f\leq g$，则 $\int f\leq\int g$。
\end{proposition}

\begin{proof}
单调性直接来自 \eqref{eq:3-1-positive-integral}：$f$ 下方的每个简单函数也是 $g$ 下方的候选者。
若 $a>0$，映射 $s\mapsto as$ 把 $f$ 下方的简单函数与 $af$ 下方的简单函数一一对应，
所以 $\int af=a\int f$；$a=0$ 由约定成立。因而只须证明可加性。

若 $p\leq f$、$q\leq g$ 为非负简单函数，则 $p+q\leq f+g$，由引理
\ref{lemma:3-1-1-1}，
\[
 \int(f+g)\geq\int(p+q)=\int p+\int q.
\]
对非负扩展实数族有
$\sup_{p,q}(A_p+B_q)=\sup_pA_p+\sup_qB_q$：有限情形可分别逼近两个上确界，
任一上确界为无穷时则先把对应项取得任意大。因而分别对 $p,q$ 取上确界，得到
$\int(f+g)\geq\int f+\int g$。

为证另一方向，固定任意有限值简单函数 $s\leq f+g$，并令
\[
 u=s\wedge f,
 \qquad v=s-u.
\]
由于 $0\leq u\leq s$，$u,v$ 都是有限值可测函数，且 $u+v=s$、$u\leq f$。
若 $f\geq s$，则 $v=0$；若 $f<s$，则 $v=s-f\leq g$。故始终有 $v\leq g$。
令 $u_n=s_n(u)$、$v_n=s_n(v)$。它们是简单函数，且
$u_n+v_n\uparrow u+v=s$。引理~\ref{lemma:3-1-1-layer} 的最后一项以及简单函数正线性给出
\[
 \int s
 =\lim_n\int(u_n+v_n)
 =\lim_n\left(\int u_n+\int v_n\right)
 \leq\int f+\int g.
\]
对所有 $s\leq f+g$ 取上确界，便得反向不等式。最后把可加性和齐次性合并即得
\eqref{eq:3-1-positive-linearity}。
\end{proof}

\begin{proposition}[零积分判据]\label{proposition:3-1-1-3}
对非负可测函数 $f$，
\[
 \int f\,\dd\mu=0
 \quad\Longleftrightarrow\quad
 f=0\quad\mu\text{-a.e.}
 \tag{3.1.6}\label{eq:3-1-zero-integral}
\]
因而非负可测函数若 a.e. 相等，其积分相同。
\end{proposition}

\begin{proof}
若 $\int f=0$，令 $E_n=\{f>1/n\}$。因为
$(1/n)\1_{E_n}\leq f$，单调性给出
\[
 0\leq n^{-1}\mu(E_n)\leq\int f=0,
\]
故 $\mu(E_n)=0$。而 $\{f>0\}=\bigcup_nE_n$，所以 $f=0$ a.e.

反过来，设 $f=0$ 于 $X\setminus N$，其中 $\mu(N)=0$。若简单函数 $0\leq s\leq f$，
则 $s$ 的每个正值层都包含于 $N$，故 $\int s=0$。对 $s$ 取上确界即得 $\int f=0$。

最后若 $f=g$ 于 $X\setminus N$，则 $f\1_N$ 与 $g\1_N$ 均由前半部分积分为零；
命题~\ref{proposition:3-1-1-2} 给出
\[
 \int f=\int f\1_{X\setminus N}=\int g\1_{X\setminus N}=\int g.
\]
即使把零函数在一个零测点改成 $+\infty$，这个结论仍然适用。
\end{proof}

现在回算前面的原子模型。对任意 $f\geq0$，规范简单函数满足
$s_n(f)(x)\uparrow f(x)$，故
\[
 \int f\,\dd\delta_x=\sup_ns_n(f)(x)=f(x).
\]
若 $\mu=\sum_{j=1}^Nw_j\delta_{x_j}$ 且各 $x_j$ 互异，则简单函数情形由定义给出
$\int s\,\dd\mu=\sum_jw_js(x_j)$；对 $s_n(f)$ 使用这一等式并令 $n\to\infty$，
有限个非负项可逐项取极限，得到
$\int f\,\dd\mu=\sum_jw_jf(x_j)$。
若 $\mu$ 是计数测度，则简单函数计算和上确界定义给出
\[
 \int f\,\dd\mu
 =\sup\left\{\sum_{x\in F}f(x):F\subset X\text{ 有限}\right\}.
\]
确切地，对有限 $F$ 与 $M>0$，简单函数
$\sum_{x\in F}(f(x)\wedge M)\1_{\{x\}}$ 位于 $f$ 下方；先令 $M\to\infty$，再对 $F$ 取上确界，
给出 ``$\geq$''。
为核对反向不等式，任一 $s\leq f$ 的正值层若有一个无限集，$\int s$ 与右侧上确界都为无穷；
否则 $s$ 的正支撑有限，$\int s=\sum_xs(x)\leq\sup_F\sum_{x\in F}f(x)$。
这就是非负无序和；此处没有任何重排造成的条件抵消。

\begin{theorem}[Radon 开集公式]\label{theorem:3-1-1-4}
设 $X$ 为局部紧 Hausdorff 空间，$\mu$ 为 $X$ 上的 Radon 测度，$U\subset X$ 开。则
\[
 \mu(U)=\sup\left\{\int_X f\,\dd\mu:
 f\in C_c(X),\ 0\leq f\leq1,\ \supp f\subset U\right\}.
 \tag{3.1.7}\label{eq:3-1-radon-open}
\]
\end{theorem}

\begin{proof}
若 $f$ 属于右侧测试类，则 $f\leq\1_U$，所以
$\int f\leq\mu(U)$。反过来，任取紧集 $K\subset U$。局部紧 Hausdorff 空间上的紧支撑截断定理
\ref{theorem:1-3-1-2} 给出 $f\in C_c(X)$，满足
$0\leq f\leq1$、$f=1$ 于 $K$ 且 $\supp f\subset U$。于是
\[
 \mu(K)=\int\1_K\,\dd\mu\leq\int f\,\dd\mu.
\]
Radon 测度在开集上的内正则性给出
$\mu(U)=\sup_{K\subset U,\,K\text{ 紧}}\mu(K)$，故右侧上确界不小于 $\mu(U)$。
\end{proof}

\begin{example}[区间上的帽函数]\label{ex:3-1-1-radon-hat}
取 $X=\R$、$U=(a,b)$，并设 $0<\varepsilon<(b-a)/4$。定义
\[
 f_\varepsilon(x)=
 \begin{cases}
 0,&x\leq a+\varepsilon\ \text{或}\ x\geq b-\varepsilon,\\
 (x-a-\varepsilon)/\varepsilon,&a+\varepsilon<x<a+2\varepsilon,\\
 1,&a+2\varepsilon\leq x\leq b-2\varepsilon,\\
 (b-\varepsilon-x)/\varepsilon,&b-2\varepsilon<x<b-\varepsilon.
 \end{cases}
\]
这是定理~\ref{theorem:3-1-1-4} 中的测试函数。在线性斜坡上用 $N$ 个等长小区间作下阶梯和上阶梯，
相应简单函数的积分分别为
\[
 \frac{\varepsilon}{N}\sum_{j=0}^{N-1}\frac jN
 =\frac{\varepsilon(N-1)}{2N},
 \qquad
 \frac{\varepsilon}{N}\sum_{j=1}^{N}\frac jN
 =\frac{\varepsilon(N+1)}{2N}.
\]
它们分别从下、从上夹住斜坡，故令 $N\to\infty$ 得每条斜坡的积分为 $\varepsilon/2$。因此
\[
 \int f_\varepsilon\,\dd x
 =(b-a-4\varepsilon)+2\cdot\frac{\varepsilon}{2}
 =b-a-3\varepsilon.
\]
为核对所用的单调逼近，令 $d=x-a$。在区间左半部，帽函数可写为
\[
 \max\left\{0,\min\left\{1,\frac d\varepsilon-1\right\}\right\},
\]
该量随 $\varepsilon\downarrow0$ 单调不减，并在 $d>0$ 时趋于 $1$；右半部对
$b-x$ 有同样的表示，而 $x=a,b$ 时始终取零。因此，例如对 $k\geq1$ 取
$\varepsilon_k=(b-a)2^{-k-2}$，就有
\[
 f_{\varepsilon_k}(x)\uparrow\1_{(a,b)}(x)
 \quad(x\in\R),
 \qquad
 \int f_{\varepsilon_k}\,\dd x
 =b-a-3(b-a)2^{-k-2}\longrightarrow b-a.
\]
连续紧支撑测试确实恢复了开区间的长度。
\end{example}

\begin{remark}[拓扑与正则性确实参与了公式]
定理~\ref{theorem:3-1-1-4} 不是任意测度空间上的积分恒等式：没有拓扑时，$C_c(X)$ 根本没有
定义；即使有拓扑，若缺少开集内正则性或把紧集截在开集内的连续帽函数，也无法从
$\mu(K)$ 反向逼近 $\mu(U)$。因此右侧测试类之所以能恢复开集质量，恰好同时使用了
局部紧 Hausdorff 结构与 Radon 正则性。
\end{remark}

\begin{proposition}[非负推前积分公式]\label{proposition:3-1-1-5}
设 $T:(X,\mathcal A)\to(Y,\mathcal B)$ 可测，$\mu$ 为 $X$ 上的测度。对每个非负
$\mathcal B$-可测函数 $f$，
\[
 \int_Y f\,\dd(T_*\mu)=\int_Xf\circ T\,\dd\mu.
 \tag{3.1.8}\label{eq:3-1-pushforward-positive}
\]
\end{proposition}

\begin{proof}
对 $f=\1_B$，两边分别为
$(T_*\mu)(B)$ 与 $\mu(T^{-1}B)$，由推前测度的定义相等。有限正线性继而给出简单函数情形。
对一般 $f$，令 $s_n=s_n(f)$。由于截断、乘法和取整都是逐点操作，
\[
 s_n(f)\circ T=s_n(f\circ T).
\]
引理~\ref{lemma:3-1-1-layer} 分别应用于两侧，再用简单函数情形，得到
\[
 \int_Yf\,\dd(T_*\mu)
 =\sup_n\int_Ys_n(f)\,\dd(T_*\mu)
 =\sup_n\int_Xs_n(f\circ T)\,\dd\mu
 =\int_Xf\circ T\,\dd\mu.
\]
证明并未假定 $f\circ T$ 下方的每个简单函数都来自某个复合函数。
\end{proof}

对经验测度，上式两边都等于 $n^{-1}\sum_{j=1}^nf(Tx_j)$；抽象公式正是有限样本恒等式的延伸。

\begin{figure}[htbp]
\centering
\begin{tikzpicture}[x=1.05cm,y=1.05cm]
  \draw[->,book line] (-0.2,0)--(6.3,0) node[right]{$x$};
  \draw[->,book line] (0,-0.2)--(0,3.7) node[above]{$y$};
  \draw[BookAccent,very thick,smooth] plot coordinates
    {(0,0.2) (0.7,0.65) (1.2,0.55) (1.8,1.45) (2.4,1.2)
     (3.0,2.35) (3.7,2.0) (4.3,3.1) (5.0,2.55) (5.8,3.25)};
  \draw[BookWarm,thick]
    (0,0)--(0.7,0)--(0.7,0.5)--(1.2,0.5)--(1.2,0.5)--(1.8,0.5)
    --(1.8,1.25)--(2.4,1.25)--(2.4,1.0)--(3.0,1.0)--(3.0,2.0)
    --(3.7,2.0)--(3.7,1.75)--(4.3,1.75)--(4.3,2.75)--(5.0,2.75)
    --(5.0,2.5)--(5.8,2.5)--(5.8,0);
  \fill[BookWarm!12] (0,0)--(0.7,0)--(0.7,0.5)--(1.8,0.5)--(1.8,1.25)
    --(2.4,1.25)--(2.4,1)--(3,1)--(3,2)--(3.7,2)--(3.7,1.75)--(4.3,1.75)
    --(4.3,2.75)--(5,2.75)--(5,2.5)--(5.8,2.5)--(5.8,0)--cycle;
  \node[book label,BookAccent] at (4.9,3.45) {$f$};
  \node[book label,BookWarm] at (2.0,0.78) {$s_n(f)$};
\end{tikzpicture}
\caption{非负积分由图像下方的有限阶梯填充；网格越细、截断越高，阶梯单调上升。}
\label{fig:3-1-simple-staircase}
\end{figure}

\begin{remark}[从集合路线到 Daniell 语言]
若一开始给定非负函数锥上的正线性泛函 $I$，自然会令
$\mu(E)=I(\1_E)$，在简单函数上要求
$I(\sum a_k\1_{E_k})=\sum a_k\mu(E_k)$，再以简单函数下逼近的上确界扩张。
上面的构造证明了这套代数骨架在已有测度时是一致的。若要反向推出 Daniell 表示定理，还必须对
$I$ 假设适当的序列连续性，使集合函数 $E\mapsto I(\1_E)$ 具有可数可加性。
\end{remark}

\begin{exercise}
证明若 $s,t$ 是非负简单函数，则 $s\wedge t$ 与 $s\vee t$ 仍为简单函数，并直接从共同分割计算其积分。
\end{exercise}

\begin{exercise}
从两种不交值层表示的公共细分出发，独立重证引理~\ref{lemma:3-1-1-1} 的表示无关性；说明为何
证明中不会出现 $+\infty-(+\infty)$。
\end{exercise}

\begin{exercise}
对 $\mu=\sum_{j=1}^Nw_j\delta_{x_j}$，由定义~\ref{def:3-1-1-2} 证明
$\int f\,\dd\mu=\sum_jw_jf(x_j)$，包括某些 $f(x_j)=+\infty$ 的情形。
\end{exercise}

\begin{exercise}
设 $s_n(f)$ 由 \eqref{eq:3-1-canonical-simple} 定义。证明
$0\leq f-s_n(f)<2^{-n}$ 在 $\{f\leq n\}$ 上成立，并据此给出有界函数的逐点误差界。
\end{exercise}

\begin{exercise}
直接比较相邻两级的截断高度与二进网格，证明
$s_n(f)\le s_{n+1}(f)$，并逐点证明 $s_n(f)\uparrow f$，不得调用 MCT。
\end{exercise}

\begin{exercise}
由定义~\ref{def:3-1-1-2} 重新证明：若 $f\geq0$ 且 $\int f=0$，则
$\mu\{f>\varepsilon\}=0$ 对每个 $\varepsilon>0$，并推出 $f=0$ a.e.
\end{exercise}

从下逼近被写进积分的定义之后，递增极限理应是最容易处理的极限。但 ``理应'' 不是证明；
下一小节将从定义本身建立 MCT，再由它长出 Fatou 与 DCT。

\subsection{MCT、Fatou、DCT 与可数性边界}\label{subsec:3-1-2}
\index{单调收敛定理}\index{Fatou 引理}\index{支配收敛定理}

极限与积分能否交换，是积分理论的第一道压力测试。三种看似相近的函数列给出三种不同答案。
若 $E_n\uparrow E$，指标函数 $\1_{E_n}$ 没有丢失质量；若一个单位质量的区间一路向右移动，
每个固定点最终都看不见它；若质量挤进越来越窄的区间，逐点极限同样可能漏掉全部积分。
因此，逐点收敛不是通行证。单调性、下极限和可积支配分别承担不同的控制任务。
还有一处更隐蔽的断裂：把序列换成任意有向网时，即便各项递增，极限也可能由不可数次拼接产生，
而测度只承诺可数可加。序列定理建立后，我们会给出一个积分始终为零、上确界积分却为一的显式网。

\begin{example}[单调示性函数]\label{ex:3-1-2-1}
若 $E_n\uparrow E$，则 $\1_{E_n}\uparrow\1_E$，而测度从下连续已经告诉我们
$\mu(E_n)\uparrow\mu(E)$。这应当成为一般非负函数的模型。若 $E_n\downarrow E$，
则只有在 $\mu(E_1)<\infty$ 时才能从补集得到 $\mu(E_n)\downarrow\mu(E)$；有限性不是排版上的小字，
而是允许从 $\mu(E_1)$ 中作减法的条件。
\end{example}

\begin{example}[质量逃向无穷远]\label{ex:3-1-2-2}
在 $(\R,m)$ 上令 $f_n=\1_{[n,n+1]}$。对每个 $x$，至多一个 $n$ 使 $f_n(x)=1$，故
$f_n(x)\to0$；但 $\int f_n\,\dd m=1$。于是
\[
 \int\liminf_nf_n\,\dd m=0
 <1=\liminf_n\int f_n\,\dd m.
 \tag{3.1.9}\label{eq:3-1-fatou-strict}
\]
下极限只能保证不凭空增加质量，不能追赶已经跑到无穷远的质量。
\end{example}

\begin{example}[尖峰压缩]\label{ex:3-1-2-3}
在 $(0,1)$ 上令 $f_n=n\1_{(0,1/n)}$。对每个 $x>0$，当 $n>1/x$ 时 $f_n(x)=0$，
所以 $f_n\to0$ 处处；然而
\[
 \int_0^1 f_n\,\dd x=n\,m(0,1/n)=1.
\]
若存在与 $n$ 无关且可积的 $g$ 满足 $f_n\leq g$，下面的支配收敛定理将迫使积分趋零。
所以这个例子不只是说某个证明失效，它证明了统一可积支配函数根本不存在。
\end{example}

我们先证明生成其他结论的根部。

\begin{theorem}[单调收敛定理（MCT）]\label{theorem:3-1-2-1}
设 $f_n:X\to[0,+\infty]$ 可测，并且在一个共同零集之外 $f_n\leq f_{n+1}$。
在该零集上把每个 $f_n$ 改为零，并令 $f$ 为所得序列的逐点上确界。则 $f$ 可测，且
\[
 \int f_n\,\dd\mu\uparrow\int f\,\dd\mu,
 \tag{3.1.10}\label{eq:3-1-mct}
\]
等式允许两边为 $+\infty$；把 $f$ 换成任何与它 a.e. 相等的可测函数，积分不变。
\end{theorem}

\begin{proof}
设 $N$ 是共同异常零集。把每个 $f_n$ 和 $f$ 在 $N$ 上改为零；
命题~\ref{proposition:3-1-1-3} 保证积分不变。此后可以假设 $f_n\uparrow f$ 处处。
由单调性，数列 $\int f_n$ 递增且不超过 $\int f$。记其极限为 $L$。

任取非负简单函数 $s\leq f$ 和 $0<c<1$，令
\[
 E_n=\{x:f_n(x)\geq cs(x)\}.
\]
集合 $E_n$ 递增，并且 $\bigcup_nE_n$ 包含 $\{s>0\}$：若 $s(x)>0$，则
$f_n(x)\uparrow f(x)\geq s(x)>cs(x)$，所以某个 $n$ 后 $x\in E_n$。
由 $f_n\geq cs\1_{E_n}$ 与积分单调性，
\[
 \int f_n\,\dd\mu\geq c\int s\1_{E_n}\,\dd\mu.
\]
把 $s$ 写成有限个正值层 $\sum_{j=1}^ra_j\1_{A_j}$。对每个 $j$，
$A_j\cap E_n\uparrow A_j$；测度从下连续给出
\[
 \lim_n\int s\1_{E_n}\,\dd\mu
 =\sum_{j=1}^ra_j\lim_n\mu(A_j\cap E_n)
 =\int s\,\dd\mu.
\]
因此 $L\geq c\int s$。令 $c\uparrow1$，再对全部 $s\leq f$ 取上确界，得到
$L\geq\int f$。结合先前的反向不等式即得 \eqref{eq:3-1-mct}。

证明中唯一把无穷多个阶段合并起来的地方，是对递增的集合列使用测度从下连续；
这也标出了可数性进入证明的位置。
\end{proof}

\begin{corollary}[非负级数可逐项积分]\label{corollary:3-1-2-series}
若 $u_k:X\to[0,+\infty]$ 可测，则
\[
 \int_X\sum_{k=1}^{\infty}u_k\,\dd\mu
 =\sum_{k=1}^{\infty}\int_Xu_k\,\dd\mu,
 \tag{3.1.10a}\label{eq:3-1-series-integral}
\]
两边允许为 $+\infty$。
\end{corollary}

\begin{proof}
部分和 $S_n=\sum_{k=1}^n u_k$ 递增至 $S=\sum_{k=1}^\infty u_k$。
MCT 与非负积分的有限可加性依次给出
\[
 \int S=\lim_n\int S_n
 =\lim_n\sum_{k=1}^n\int u_k
 =\sum_{k=1}^\infty\int u_k.
\]
\end{proof}

\begin{corollary}[集合的单调连续性]\label{corollary:3-1-2-set-continuity}
若 $E_n\uparrow E$，则 $\mu(E_n)\uparrow\mu(E)$。若 $E_n\downarrow E$ 且
$\mu(E_1)<\infty$，则 $\mu(E_n)\downarrow\mu(E)$。
\end{corollary}

\begin{proof}
第一项对 $\1_{E_n}\uparrow\1_E$ 应用 MCT。第二项令
$F_n=E_1\setminus E_n$，则 $F_n\uparrow E_1\setminus E$，故
\[
 \mu(E_1)-\mu(E_n)=\mu(F_n)\uparrow\mu(E_1\setminus E)=\mu(E_1)-\mu(E).
\]
由于 $\mu(E_1)<\infty$，这些减法都发生在有限实数中。
\end{proof}

有限性不能删去。在 $\R$ 上取 $E_n=[n,+\infty)$，则 $E_n\downarrow\varnothing$，
但 $m(E_n)=+\infty$ 对所有 $n$ 成立，因而这些测度不可能下降到 $m(\varnothing)=0$。

回到例~\ref{ex:3-1-2-1}，示性函数计算现在已成为一般定理的直接特例。
下一条结论不要求函数列单调；代价是它只保留下界。

\begin{theorem}[Fatou 引理]\label{theorem:3-1-2-2}
若 $f_n:X\to[0,+\infty]$ 可测，则
\[
 \int_X\liminf_{n\to\infty}f_n\,\dd\mu
 \leq\liminf_{n\to\infty}\int_Xf_n\,\dd\mu.
 \tag{3.1.11}\label{eq:3-1-fatou}
\]
\end{theorem}

\begin{proof}
令 $g_n=\inf_{k\geq n}f_k$。这是可数族可测函数的下确界，因而可测，并且
$g_n\uparrow\liminf_kf_k$。MCT 给出
\[
 \int\liminf_kf_k=\lim_n\int g_n.
\]
对每个 $k\geq n$ 有 $g_n\leq f_k$，故
\[
 \int g_n\leq\inf_{k\geq n}\int f_k.
\]
令 $n\to\infty$ 即得 \eqref{eq:3-1-fatou}。
\end{proof}

例~\ref{ex:3-1-2-2} 说明 Fatou 不等式可以严格。若另有一个可积上界，
把同一不等式同时用于 ``剩余质量''，便能夹出等号。

\begin{theorem}[非负支配收敛定理]\label{theorem:3-1-2-3}
设 $g:X\to[0,\infty]$ 为可测函数，$0\leq f_n\leq g$ a.e.，且
$f_n\to f$ a.e.，其中 $f$ 可测并满足 $\int g\,\dd\mu<\infty$。则
$0\leq f\leq g$ a.e.，并且
\[
 \int|f_n-f|\,\dd\mu\longrightarrow0,
 \qquad
 \int f_n\,\dd\mu\longrightarrow\int f\,\dd\mu.
 \tag{3.1.12}\label{eq:3-1-nonnegative-dct}
\]
\end{theorem}

\begin{proof}
先令 $Z=\{g=+\infty\}$。对每个 $M>0$，有 $M\1_Z\le g$，因而
$M\mu(Z)\le\int g<\infty$；令 $M\to\infty$ 得 $\mu(Z)=0$。把 $Z$ 与收敛、支配失败的
可数多个零集并成一个零集，再在其上把 $g,f_n,f$ 全部改为零。积分不变，此后 $g$ 处处有限，
$0\leq f_n\leq g$ 且 $f_n\to f$ 处处，故 $0\leq f\leq g$。
特别地，$\int f_n$ 与 $\int f$ 都不超过有限数 $\int g$。

Fatou 引理先给出
\[
 \int f\leq\liminf_n\int f_n.
 \tag{3.1.13}\label{eq:3-1-dct-liminf}
\]
再对非负函数 $g-f_n\to g-f$ 应用 Fatou。非负积分的可加性以及 $\int g<\infty$ 给出
\[
 \int g-\int f
 \leq\liminf_n\left(\int g-\int f_n\right)
 =\int g-\limsup_n\int f_n.
\]
在有限实数中消去 $\int g$，得到
$\limsup_n\int f_n\leq\int f$。与 \eqref{eq:3-1-dct-liminf} 合并，便有积分收敛。

还需证明绝对误差趋零。函数
$h_n=2g-|f_n-f|$ 非负且处处趋于 $2g$。Fatou 引理给出
\[
 2\int g
 \leq\liminf_n\int h_n
 =\liminf_n\left(2\int g-\int|f_n-f|\right)
 =2\int g-\limsup_n\int|f_n-f|.
\]
再次在有限实数中消去 $2\int g$，得到
$\limsup_n\int|f_n-f|\leq0$。这些积分非负，所以它们趋于零。
\end{proof}

例~\ref{ex:3-1-2-3} 的尖峰由此排除了任何统一可积支配函数。
请注意，本定理只积分非负函数；允许负值和复值的版本要等到下一小节定义相应积分后才能陈述。

\begin{example}[一次参数积分的回收计算]\label{ex:3-1-2-parameter}
令
\[
 F(t)=\int_0^1e^{-tx}\,\dd x\qquad(t\in\R).
\]
若 $t_n\to t$，则存在 $M$ 使 $|t_n|\leq M$，从而
$0<e^{-t_nx}\leq e^M$ 于 $[0,1]$，而常数 $e^M$ 在该区间可积。
非负 DCT 给出 $F(t_n)\to F(t)$。由于 $\R$ 是第一可数的，序列连续性等价于连续性，故 $F$ 连续；
一元微积分的直接计算还给出
\[
 F(t)=
 \begin{cases}
 (1-e^{-t})/t,&t\ne0,\\
 1,&t=0.
 \end{cases}
\]
这里极限定理把逐点连续提升成了积分后的连续性；同一机制也适用于由核给出的参数积分。
\end{example}

\begin{definition}[一致尾控制]\label{def:3-1-2-3}
设 $\mathcal F$ 是一族实值或复值可测函数，且每个 $f\in\mathcal F$ 都满足
$\int|f|<\infty$。若
\[
 \lim_{K\to\infty}\sup_{f\in\mathcal F}
 \int_{\{|f|>K\}}|f|\,\dd\mu=0,
 \tag{3.1.14}\label{eq:3-1-uniform-tail}
\]
则称 $\mathcal F$ 具有一致尾控制。
\end{definition}

式~\eqref{eq:3-1-uniform-tail} 只涉及非负函数 $|f|$ 的积分，因而定义已经完备。
尖峰族 $f_n=n\1_{(0,1/n)}$ 不满足该条件：取 $n>K$ 时，
$\int_{|f_n|>K}|f_n|=1$。不过，在无限测度空间上，一致尾控制连一致 $L^1$ 有界性也不能单独保证。
例如 $f_n=\1_{[0,n]}$ 在 $K>1$ 时使 \eqref{eq:3-1-uniform-tail} 的上确界为零，
但 $\int f_n=n$。即使另加一致 $L^1$ 界，高度尾控制仍不阻止空间逃逸：
$g_n=\1_{[n,n+1]}$ 满足 $\|g_n\|_1=1$，且在 $K>1$ 时尾积分恒为零，但对每个紧集
$C\subset\R$ 都有 $\int_{\R\setminus C}g_n=1$ 对所有充分大的 $n$ 成立。
因此一致尾控制只限制函数的高度尾部；在无限测度空间上，还必须分别检查
$L^1$ 有界性和防止质量逃向无穷远的紧性条件。

序列极限定理背后是测度的可数可加性。拓扑中常用网描述极限，那么 MCT 能否把 $n$ 换成任意有向指标？
下面的例子把答案钉死。

\begin{example}[不可数网陷阱]\label{ex:3-1-2-4}
设 $X$ 是不可数集合，$\mathcal A$ 是可数--余可数 $\sigma$-代数，并定义
\[
 \mu(A)=
 \begin{cases}
 0,&A\text{ 可数},\\
 1,&X\setminus A\text{ 可数}.
 \end{cases}
\]
这是概率测度。以 $X$ 的所有有限子集组成的集合 $I$ 为有向集，次序取包含关系，并令
$f_F=\1_F$。则 $F\subset G$ 蕴含 $f_F\leq f_G$，而逐点上
$\sup_{F\in I}f_F=\1_X$。然而每个 $F$ 有 $\int f_F=\mu(F)=0$，而
$\int\1_X=1$。因此任意网版本的 MCT 是假的。
\end{example}

这个有向集没有可数共尾子集：若 $(F_n)$ 共尾，则可数集 $\bigcup_nF_n$ 应包含每个单点，
从而等于不可数的 $X$，矛盾。安全的推广恰好止于能够降回序列的网。

不可数方向还可能更早出故障。取一个非 Borel 集 $A\subset\R$，并令有限子集
$F\subset A$ 按包含关系成网。每个 $\1_F$ 都是 Borel 可测的，但
\[
 \sup_{F\subset A,\,F\text{ 有限}}\1_F=\1_A
\]
不是 Borel 可测函数。可数个可测函数的上确界可测，并不蕴含任意族也有这一性质。

\begin{definition}[可数共尾网]\label{def:3-1-2-2}
有向集 $I$ 若含有可数共尾子集，即存在可数 $C\subset I$，使对每个 $i\in I$ 都有
$c\in C$ 满足 $i\leq c$，则称 $I$ 可数共尾。以这样的 $I$ 为指标的网简称可数共尾网。
\end{definition}

第一可数性与这里的概念不同。第一可数性说某个拓扑空间在每一点有可数邻域基；
可数共尾性说索引集内部有一条可数的最终路径。值域即使是 $\R$，
例~\ref{ex:3-1-2-4} 的索引集仍然没有可数共尾子集。

\begin{proposition}[可数共尾网版本]\label{proposition:3-1-2-4}
设 $I$ 为可数共尾有向集。
\begin{enumerate}
 \item 若 $(f_i)_{i\in I}$ 是递增的非负可测网，且 $f=\sup_{i\in I}f_i$ 可测，则
 \[
   \sup_{i\in I}\int f_i\,\dd\mu=\int f\,\dd\mu.
 \]
 \item 若 $0\leq f_i\leq g$ a.e.，$\int g<\infty$，且 $f_i\to f$ a.e.，其中 $f$ 可测，
 则 $\int|f_i-f|\to0$。
\end{enumerate}
这里的 a.e. 陈述均使用一个对全网共同有效的零集。一般网不在本命题中享有无条件的 Fatou 引理。
\end{proposition}

\begin{proof}
枚举一个可数共尾集为 $(c_n)$。利用有向性递归选取 $d_n$，使
$d_1\geq c_1$ 且 $d_n\geq d_{n-1},c_n$。于是 $(d_n)$ 是递增共尾序列。
第一项中，$f_{d_n}\uparrow f$：它们的上确界不超过 $f$；反过来，对任意 $i$，
共尾性给出某个 $c_n\geq i$，继而 $d_n\geq c_n$，所以 $f_i\leq f_{d_n}$。
序列 MCT 因而给出
\[
 \int f=\lim_n\int f_{d_n}=\sup_{i\in I}\int f_i.
\]

对第二项，若结论失败，则存在 $\varepsilon>0$，使对每个 $i_0\in I$ 都能找到
$i\geq i_0$ 满足 $\int|f_i-f|\geq\varepsilon$。先选 $i_1\ge d_1$ 保持这一坏性质；已选 $i_{n-1}$ 后，
再递归选取 $i_n\geq d_n,i_{n-1}$ 且保持坏性质。于是 $(i_n)$ 递增，且因 $i_n\ge d_n$ 而共尾；
映射 $n\mapsto i_n$ 保序且共尾，所以它按子网的定义确实给出原网的一个子网。因此
$f_{i_n}\to f$ a.e.，且仍由 $g$ 支配。序列的非负支配收敛定理给出
$\int|f_{i_n}-f|\to0$，与其始终不小于 $\varepsilon$ 矛盾。

从可数共尾集构造递增链以及逐步选取坏指标，是可数次递归选择；按照第一章的选择约定，
这些步骤使用可数依赖选择。一般网的尾下确界可能是不可数族的下确界，甚至未必可测，
所以没有额外的可数尾决定条件时，不能把 Fatou 的序列证明照抄过去。
\end{proof}

\begin{figure}[htbp]
\centering
\begin{tikzpicture}[node distance=10mm and 13mm]
  \node[book strong node] (mct) {MCT\\递增极限};
  \node[book node,above left=of mct] (sets) {集合从下连续\\从上连续需有限性};
  \node[book node,above right=of mct] (fatou) {Fatou\\尾下确界};
  \node[book warm node,right=of fatou] (dct) {非负 DCT\\可积支配};
  \node[book warning node,below right=of mct] (net) {任意网断裂\\不可数拼接};
  \draw[book accent arrow] (mct)--(sets);
  \draw[book accent arrow] (mct)--(fatou);
  \draw[book accent arrow] (fatou)--(dct);
  \draw[book dashed arrow] (mct)--(net);
  \node[book label,below=2mm of dct] {$2g-|f_n-f|$ 控制误差};
\end{tikzpicture}
\caption{极限定理的生成关系。虚线分支不是另一个定理，而是不可数网反例标出的边界。}
\label{fig:3-1-convergence-tree}
\end{figure}

\begin{remark}[三个直接用途]
推论~\ref{corollary:3-1-2-series} 已经给出级数与积分交换的非负版本，
例~\ref{ex:3-1-2-parameter} 则展示 DCT 如何产生参数积分连续性。
第三种用途是把 Cauchy 列抽成绝对增量可求和的快速子列：若
$\sum_k\|f_{n_{k+1}}-f_{n_k}\|_p<\infty$，则
\eqref{eq:3-1-series-integral} 能把范数估计转化为几乎处处的绝对收敛。这正是
Riesz--Fischer 完备性证明的核心机制。
\end{remark}

\begin{exercise}
利用推论~\ref{corollary:3-1-2-series} 计算
$\int_0^1\sum_{k=0}^\infty x^k\,\dd x$，并解释为什么答案为 $+\infty$。
\end{exercise}

\begin{exercise}
设 $0\leq f_n\leq g$ 且 $\int g<\infty$。用 Fatou 引理分别作用于 $f_n$ 与 $g-f_n$，
证明
\[
 \limsup_n\int f_n\leq\int\limsup_nf_n.
\]
说明这里为什么需要逐点不等式 $f_n\leq g$。
\end{exercise}

\begin{exercise}
构造三个函数列，分别说明非负 DCT 中 ``a.e. 收敛''、``同一个支配函数'' 和
``支配函数可积'' 不能任意删去。
\end{exercise}

\begin{exercise}
完整证明例~\ref{ex:3-1-2-4} 中的 $\mu$ 是可数可加测度，并证明有限子集有向集没有可数共尾子集。
\end{exercise}

非负极限定理已经允许我们稳定地处理大小，却仍不允许抵消。
$+\infty-(+\infty)$ 不是一个等待聪明记号来修补的空格；下一小节将精确规定何时减法合法。

\subsection{实值与复值积分：抵消、绝对值与局部化}\label{subsec:3-1-3}
\index{正部与负部}\index{$L^1$ 空间@$L^1$ 空间}\index{局部可积}

在 $\R$ 上看函数 $f(x)=x$。每个对称截断 $[-R,R]$ 都给出零，但在整个实线上，
正半轴和负半轴各自携带无穷质量。若把这两个无穷硬写成 ``相减为零''，换一种截断方式就会得到别的答案。
对称主值是一种附加的求和规则，不是 Lebesgue 积分。我们需要的原则是：先分别测量正负质量，
只有在减法有意义时才允许抵消。

\begin{definition}[正负部分]\label{def:3-1-3-1}
对扩展实值可测函数 $f$，定义
\[
 f^+=\max\{f,0\},\qquad f^-=\max\{-f,0\}.
\]
于是 $f=f^+-f^-$ 且 $|f|=f^++f^-$。若 $\int f^+$ 与 $\int f^-$ 不同时为
$+\infty$，定义
\[
 \int f\,\dd\mu:=\int f^+\,\dd\mu-\int f^-\,\dd\mu.
 \tag{3.1.15}\label{eq:3-1-signed-integral}
\]
若两者都有限，则称 $f$ 可积。
\end{definition}

\begin{definition}[$L^1$ 可积]\label{def:3-1-3-2}
实值或复值可测函数 $f$ 若满足
\[
 \int_X|f|\,\dd\mu<\infty,
\]
就称为可积。所有这样的点值函数组成 $\mathcal L^1(\mu)$；按 a.e. 相等取商所得空间记为
$L^1(\mu)$。对复值 $f$ 定义
\[
 \int f\,\dd\mu:=\int\Re f\,\dd\mu+i\int\Im f\,\dd\mu,
 \qquad
 \|[f]\|_1:=\int|f|\,\dd\mu.
 \tag{3.1.16}\label{eq:3-1-complex-integral}
\]
不致混淆时省略等价类的方括号。
\end{definition}

因为 $|\Re f|,|\Im f|\leq|f|$，复积分右侧确实由两个有限实数构成。
命题~\ref{proposition:3-1-1-3} 保证 $\|[f]\|_1=0$ 当且仅当 $[f]=0$；
所以它在等价类空间上才是范数，在点值函数空间上只是半范数。

\begin{example}[条件抵消不可定义]\label{ex:3-1-3-1}
对 $f(x)=x$，把 $[0,R]$ 等分成 $N$ 段。以左、右端点高度构成的简单函数分别夹住 $x$，其积分为
\[
 \frac{R^2(N-1)}{2N}
 \leq\int_0^R x\,\dd x
 \leq\frac{R^2(N+1)}{2N}.
\]
令 $N\to\infty$ 得 $\int_0^Rx\,\dd x=R^2/2$。
因此 $\int_\R f^+=\int_0^\infty x\,\dd x=+\infty$；同样
$\int_\R f^-=+\infty$。故 $\int_\R x\,\dd x$ 没有定义。
另一方面，每个 $[-R,R]$ 上的积分为零：把两侧分别用等长阶梯计算，正负部分的有限积分相等。
这说明 ``所有对称截断都为零'' 并不能替代定义~\ref{def:3-1-3-1}。
\end{example}

\begin{definition}[局部可积]\label{def:3-1-3-3}
设 $X$ 为局部紧 Hausdorff 空间，$\mu$ 为 Radon 测度。若可测函数 $f$ 对每个紧集
$K\subset X$ 都满足
\[
 \int_K|f|\,\dd\mu<\infty,
\]
则称 $f$ 局部可积，记作 $f\in L^1_{\mathrm{loc}}(X,\mu)$。
\end{definition}

若 $X$ 还第二可数，第一章给出的紧耗尽满足
$K_n\subset\operatorname{int}K_{n+1}$ 且 $\bigcup_nK_n=X$。此时只须检查每个 $K_n$：
任意紧集 $K$ 被开覆盖 $\{\operatorname{int}K_n\}$ 覆盖，有限子覆盖与递增性给出
$K\subset K_N$。

\begin{example}[局部可积但不可积]\label{ex:3-1-3-2}
相对于 $\R^n$ 上的 Lebesgue 测度，常数函数 $1$ 在每个紧集 $K$ 上满足
$\int_K1=m_n(K)<\infty$，故属于 $L^1_{\mathrm{loc}}(\R^n)$。但
\[
 \int_{\R^n}1\,\dd x=m_n(\R^n)=+\infty,
\]
所以它不属于 $L^1(\R^n)$。局部化允许我们在每个有限窗口内工作，并没有把无穷空间伪装成有限空间。
\end{example}

\begin{example}[幂奇点的局部可积阈值]\label{ex:3-1-3-power}
在 $\R^n$ 上令 $f(x)=|x|^{-\alpha}$，其中 $\alpha\geq0$，并令 $f(0)=+\infty$。
只需检查原点附近。置
\[
 A_k=\{x:2^{-k-1}<|x|\leq2^{-k}\},\qquad k\geq0.
\]
若 $v_n=m_n(B(0,1))$，第~2~章的线性体积缩放给出
\[
 m_n(A_k)=v_n(1-2^{-n})2^{-kn}.
\]
在 $A_k$ 上有
$2^{k\alpha}\leq f(x)<2^{(k+1)\alpha}$。对不交环带的有限并使用非负可加性，
再令环带数目增加并用 MCT，得到
\[
 C_1\sum_{k=0}^N2^{-k(n-\alpha)}
 \leq\int_{\bigcup_{k=0}^NA_k}|x|^{-\alpha}\,\dd x
 \leq C_2\sum_{k=0}^N2^{-k(n-\alpha)},
 \tag{3.1.17}\label{eq:3-1-power-annuli}
\]
其中 $C_1=v_n(1-2^{-n})$、$C_2=2^\alpha C_1$。
几何级数收敛当且仅当 $\alpha<n$。原点是零测单点，远离原点时函数在紧集上有界，故
\[
 |x|^{-\alpha}\in L^1_{\mathrm{loc}}(\R^n)
 \quad\Longleftrightarrow\quad \alpha<n.
\]
极坐标公式也会给出同一阈值；上面的环带论证则只使用体积缩放与几何级数。
\end{example}

\begin{proposition}[限制测度与集合上积分]\label{proposition:3-1-3-restriction}
对 $E\in\mathcal A$ 定义限制测度
$(\mu|_E)(A)=\mu(A\cap E)$。若 $f$ 非负可测或可积，则
\[
 \int_Ef\,\dd\mu
 =\int_Xf\1_E\,\dd\mu
 =\int_Xf\,\dd(\mu|_E).
 \tag{3.1.18}\label{eq:3-1-restriction}
\]
\end{proposition}

\begin{proof}
第一式是记号的定义。若 $f=\1_A$，则后两式都是 $\mu(A\cap E)$；有限正线性给出简单函数情形。
对一般非负 $f$，$s_n(f)$ 在两种测度下都是规范简单逼近，故由
引理~\ref{lemma:3-1-1-layer}，
\[
 \int f\,\dd(\mu|_E)
 =\sup_n\int s_n(f)\,\dd(\mu|_E)
 =\sup_n\int s_n(f)\1_E\,\dd\mu
 =\int f\1_E\,\dd\mu.
\]
实值可积函数拆成正负部分，复值函数拆成实虚部，即得一般情形。
\end{proof}

经验测度下，\eqref{eq:3-1-restriction} 的三式都等于
$n^{-1}\sum_{j:x_j\in E}f(x_j)$，限制测度就是把不符合筛选条件的样本权重置零。

下面建立 $L^1$ 的代数。证明不能先假定 ``积分差由差的积分控制''，因为这句话本身需要积分线性。

\begin{theorem}[$L^1$ 上积分的线性]\label{theorem:3-1-3-2}
$L^1(\mu)$ 是实或复向量空间，并且
\[
 I:L^1(\mu)\to\R\ \text{或}\ \C,
 \qquad I(f)=\int f\,\dd\mu,
\]
是线性泛函。a.e. 相等的可积函数有相同积分。
\end{theorem}

\begin{proof}
若 $f,g\in\mathcal L^1$，则点态三角不等式和非负积分的可加性给出
\[
 \int|f+g|\leq\int(|f|+|g|)=\int|f|+\int|g|<\infty;
\]
对任意标量 $a$，$\int|af|=|a|\int|f|<\infty$。因此可积函数形成向量空间，a.e. 等价与
向量运算相容。

先设 $f,g$ 实值。点态恒等式
\[
 (f+g)^+ +f^-+g^-
 =(f+g)^-+f^++g^+
 \tag{3.1.19}\label{eq:3-1-linearity-identity}
\]
两边都是非负可积函数。对两边积分并使用命题~\ref{proposition:3-1-1-2}，
所有项均为有限实数，故可以移项，得到
\[
 \int(f+g)=\int f+\int g.
\]
当 $a\geq0$ 时，$(af)^+=af^+$、$(af)^-=af^-$；当 $a<0$ 时，
正负部分交换并乘以 $|a|$。所以 $\int af=a\int f$。

对复值函数，实虚部分别应用上述结论，先得加法性。又
\[
 \Re(if)=-\Im f,
 \qquad \Im(if)=\Re f,
\]
故 $\int(if)=i\int f$；结合实齐次性便得复线性。

若 $f=g$ a.e.，则 $|f-g|=0$ a.e.，命题~\ref{proposition:3-1-1-3} 给出
$\int|f-g|=0$。在实值情形，$0\le(f-g)^\pm\le|f-g|$，故两个正负部分的积分都为零，
从定义得到 $\int(f-g)=0$；复值情形对实部、虚部重复此论证，因为二者的绝对值都不超过
$|f-g|$。由已经在点值可积函数上证明的线性，$\int f-\int g=\int(f-g)=0$。因此积分确实
下降为 $L^1$ 等价类上的线性泛函。
\end{proof}

\begin{example}[积分平均与中心化]\label{ex:3-1-3-expectation}
设 $(\Omega,\mathcal F,\Prob)$ 是总质量为一的测度空间。对可积可测函数 $X$，记其积分平均为
$\operatorname{Av}_{\Prob}(X):=\int_\Omega X\,\dd\Prob$；在概率论中，这一数值称为期望。常数函数
$c$ 的积分是
$c\Prob(\Omega)=c$，故积分线性给出
\[
 \operatorname{Av}_{\Prob}\!\left(X-\operatorname{Av}_{\Prob}(X)\right)
 =\operatorname{Av}_{\Prob}(X)
  -\operatorname{Av}_{\Prob}(X)\Prob(\Omega)=0.
\]
因此 ``中心化'' 不是形式上的减号，而是由总质量为一和 $L^1$ 线性共同保证的零平均操作。
\end{example}

\begin{proposition}[绝对值估计]\label{proposition:3-1-3-1}
对每个 $f\in L^1(\mu)$，
\[
 \left|\int f\,\dd\mu\right|\leq\int|f|\,\dd\mu.
 \tag{3.1.20}\label{eq:3-1-integral-absolute}
\]
\end{proposition}

\begin{proof}
若 $f$ 实值，则 $-|f|\leq f\leq|f|$。积分线性与非负积分单调性给出
$-\int|f|\leq\int f\leq\int|f|$。

若 $f$ 复值且 $z=\int f\ne0$，取单位复数 $\lambda=\overline z/|z|$，使
$\lambda z=|z|$。由已经证明的复线性和实值情形，
\[
 \left|\int f\right|
 =\Re\left(\lambda\int f\right)
 =\int\Re(\lambda f)
 \leq\int|\lambda f|
 =\int|f|.
\]
当 $z=0$ 时结论成立。
\end{proof}

\begin{example}[复相位只改变抵消]\label{ex:3-1-3-3}
设 $g$ 可测，$\theta:X\to\R$ 可测，并令 $f=e^{i\theta}g$。连续复合保证
$e^{i\theta}$ 可测，而且 $|f|=|g|$，所以
\[
 \int|f|=\int|g|
\]
允许共同为 $+\infty$。若这一数值有限，则 $f,g\in L^1$，但一般
$\int f\ne\int g$：相位可以制造抵消，却不会改变控制量。命题~\ref{proposition:3-1-3-1}
说明绝对值恰是所有相位共同服从的上界。
\end{example}

\begin{corollary}[实值与复值推前积分公式]\label{corollary:3-1-3-pushforward}
设 $T:X\to Y$ 可测。若 $f:Y\to\R$ 或 $\C$ 对 $T_*\mu$ 可积，则
$f\circ T$ 对 $\mu$ 可积，且
\[
 \int_Yf\,\dd(T_*\mu)=\int_Xf\circ T\,\dd\mu.
 \tag{3.1.21}\label{eq:3-1-pushforward-signed}
\]
\end{corollary}

\begin{proof}
对非负函数 $|f|$ 使用命题~\ref{proposition:3-1-1-5}，得到
\[
 \int_X|f\circ T|\,\dd\mu
 =\int_Y|f|\,\dd(T_*\mu)<\infty.
\]
故复合函数可积。实值时分别对 $f^+,f^-$ 使用非负推前公式；复值时再对实部、虚部使用实值结论。
\end{proof}

\begin{theorem}[实值与复值支配收敛]\label{theorem:3-1-3-5}
设 $f_n:X\to\R$ 或 $\C$ 可测，$f_n\to f$ a.e.，其中 $f$ 可测。若存在非负可测函数
$g$ 满足 $|f_n|\leq g$ a.e. 且 $\int g<\infty$，则 $f,f_n\in L^1(\mu)$，并且
\[
 \|f_n-f\|_1\longrightarrow0,
 \qquad
 \int f_n\,\dd\mu\longrightarrow\int f\,\dd\mu.
 \tag{3.1.22}\label{eq:3-1-full-dct}
\]
\end{theorem}

\begin{proof}
共同化异常零集并在其上把所有函数改为零后，可设收敛与支配处处成立。
取极限得 $|f|\leq g$，故 $f,f_n$ 可积；并且
$|f_n-f|\leq2g$ 且 $|f_n-f|\to0$。对非负函数列 $|f_n-f|$ 应用
定理~\ref{theorem:3-1-2-3}，得到 $\|f_n-f\|_1\to0$。最后由积分线性和绝对值估计，
\[
 \left|\int f_n-\int f\right|
 =\left|\int(f_n-f)\right|
 \leq\int|f_n-f|\longrightarrow0.
\]
\end{proof}

这就是实值或复值可积函数的支配收敛定理（DCT）。

\begin{proposition}[限制与拼接]\label{proposition:3-1-3-3}
若 $(E_k)_{k\geq1}$ 是两两不交的可测集且 $f\in L^1(\mu)$，则
\[
 \int_{\bigcup_{k=1}^\infty E_k}f\,\dd\mu
 =\sum_{k=1}^\infty\int_{E_k}f\,\dd\mu,
 \tag{3.1.23}\label{eq:3-1-patching}
\]
右侧级数绝对收敛。
\end{proposition}

\begin{proof}
对非负函数 $|f|$，有限部分和
$\sum_{k=1}^n|f|\1_{E_k}$ 递增至 $|f|\1_{\cup_kE_k}$。MCT 与有限可加性给出
\[
 \sum_{k=1}^\infty\int_{E_k}|f|
 =\int_{\cup_kE_k}|f|
 \leq\int|f|<\infty.
\]
由命题~\ref{proposition:3-1-3-1}，
$|\int_{E_k}f|\leq\int_{E_k}|f|$，故右侧绝对收敛。
对每个 $n$，积分线性给出
\[
 \sum_{k=1}^n\int_{E_k}f=\int_{\cup_{k=1}^nE_k}f.
\]
两边与目标值之差的绝对值不超过
$\int_{\cup_{k>n}E_k}|f|$；该量由上面的收敛级数趋于零。令 $n\to\infty$ 即得公式。
\end{proof}

有限复测度的积分还要解决另一个问题：积分变量本身携带相位。第~2~章由分割定义的全变差
$|\nu|$ 能直接控制简单函数积分，标量 $L^1(|\nu|)$ 的完备性继而给出延拓。

\begin{theorem}[复测度积分估计]\label{theorem:3-1-3-4}
设 $\nu$ 是 $(X,\mathcal A)$ 上的有限复测度。若
$s=\sum_{k=1}^ma_k\1_{E_k}$ 是复值简单函数，其中 $E_k$ 两两不交，定义
\[
 \int_Xs\,\dd\nu:=\sum_{k=1}^ma_k\nu(E_k).
\]
这个定义与表示无关，并且可唯一延伸到每个 $f\in L^1(|\nu|)$，使
\[
 \left|\int_Xf\,\dd\nu\right|
 \leq\int_X|f|\,\dd|\nu|.
 \tag{3.1.24}\label{eq:3-1-complex-measure-bound}
\]
延伸后的积分在 $L^1(|\nu|)$ 上复线性。
\end{theorem}

\begin{proof}
表示无关的证明与引理~\ref{lemma:3-1-1-1} 相同：取公共有限分割，并使用复测度的有限可加性。
由全变差的定义和 $|\nu(E)|\leq|\nu|(E)$，
\[
 \left|\int s\,\dd\nu\right|
 \leq\sum_k|a_k|\,|\nu(E_k)|
 \leq\sum_k|a_k|\,|\nu|(E_k)
 =\int|s|\,\dd|\nu|.
 \tag{3.1.25}\label{eq:3-1-simple-complex-measure-bound}
\]

给定 $f\in L^1(|\nu|)$，分别以规范层级逼近 $\Re f$、$\Im f$ 的正负部分：
\[
 s_n=\bigl(s_n((\Re f)^+)-s_n((\Re f)^-)\bigr)
   +i\bigl(s_n((\Im f)^+)-s_n((\Im f)^-)\bigr).
\]
这些是复值简单函数。MCT 给出
$\int s_n(h)\uparrow\int h$ 对每个非负可积的 $h$；由非负可加性，
\[
 \int(h-s_n(h))\,\dd|\nu|
 =\int h\,\dd|\nu|-\int s_n(h)\,\dd|\nu|\longrightarrow0.
\]
四项相加并用点态三角不等式，得到
\[
 \int|s_n-f|\,\dd|\nu|\longrightarrow0.
 \tag{3.1.26}\label{eq:3-1-simple-l1-approx}
\]
由 \eqref{eq:3-1-simple-complex-measure-bound}，
\[
 \left|\int s_n\,\dd\nu-\int s_m\,\dd\nu\right|
 \leq\int|s_n-s_m|\,\dd|\nu|,
\]
所以 $\bigl(\int s_n\,\dd\nu\bigr)$ 是 $\C$ 中的 Cauchy 列。利用 $\C$ 的完备性，定义
\[
 \int f\,\dd\nu:=\lim_n\int s_n\,\dd\nu.
\]
若 $(t_n)$ 是另一列满足 $\int|t_n-f|\,\dd|\nu|\to0$ 的简单函数，
则
\[
 \left|\int s_n\,\dd\nu-\int t_n\,\dd\nu\right|
 \leq\int|s_n-f|\,\dd|\nu|+\int|t_n-f|\,\dd|\nu|\to0,
\]
故极限与逼近选择无关。对简单逼近逐项使用线性并取极限，得到复线性。

最后，$\bigl||s_n|-|f|\bigr|\leq|s_n-f|$，所以
$\int|s_n|\,\dd|\nu|\to\int|f|\,\dd|\nu|$。在
\eqref{eq:3-1-simple-complex-measure-bound} 中令 $n\to\infty$，便得
\eqref{eq:3-1-complex-measure-bound}。若另一个延伸在线性且受同一 $L^1$ 误差控制，
它在简单函数上相同，并由 \eqref{eq:3-1-simple-l1-approx} 在每个 $f$ 上相同，故延伸唯一。
整个构造只使用了标量域的完备性，没有调用 Hahn--Banach、$L^p$ 完备性或复测度的极分解。
\end{proof}

\begin{example}[两个复原子]\label{ex:3-1-3-complex-atoms}
令
\[
 \nu=i\delta_0-(1+i)\delta_1.
\]
分割定义给出 $|\nu|=\delta_0+\sqrt2\,\delta_1$。因此对任意函数 $f:\{0,1\}\to\C$，
\[
 \int f\,\dd\nu=if(0)-(1+i)f(1),
\]
且
\[
 \left|if(0)-(1+i)f(1)\right|
 \leq |f(0)|+\sqrt2|f(1)|
 =\int|f|\,\dd|\nu|.
\]
这里相位与控制量都能逐原子看见。
\end{example}

\begin{figure}[htbp]
\centering
\begin{tikzpicture}[x=1cm,y=1cm]
  \draw[->,book line] (-3.8,0)--(3.8,0) node[right]{$x$};
  \draw[->,book line] (0,-2.4)--(0,2.7) node[above]{$y$};
  \draw[BookAccent,very thick,smooth] plot coordinates
    {(-3,0.2) (-2.4,1.25) (-1.7,1.75) (-1.0,0.55) (-0.3,0)
     (0.35,-0.8) (1.1,-1.65) (1.8,-0.7) (2.5,0) (3.1,1.0)};
  \fill[BookAccent!12] (-3,0)--(-3,0.2)--(-2.4,1.25)--(-1.7,1.75)--(-1,0.55)--(-0.3,0)--cycle;
  \fill[BookWarning!12] (0.35,0)--(0.35,-0.8)--(1.1,-1.65)--(1.8,-0.7)--(2.5,0)--cycle;
  \draw[BookWarning,dashed,thick,smooth] plot coordinates
    {(0.35,0.8) (1.1,1.65) (1.8,0.7) (2.5,0)};
  \node[book label,BookAccent] at (-1.8,2.15) {$f^+$};
  \node[book label,BookWarning] at (1.2,-2.0) {$f^-$};
  \node[book label,BookWarning] at (1.3,2.1) {$|f|$ 将负区翻上};
\end{tikzpicture}
\caption{带符号积分是正区质量减去负区质量；绝对值把两区都变成非负控制量。}
\label{fig:3-1-positive-negative-parts}
\end{figure}

\begin{remark}[复测度的极表示]
定理~\ref{theorem:3-1-3-4} 直接从全变差定义了 $\int f\,\dd\nu$。稍后的
定理~\ref{theorem:3-3-1-3} 还将给出 $\nu=u|\nu|$，其中 $|u|=1$ a.e.；由此可把同一积分写成
\[
  \int f\,\dd\nu=\int fu\,\dd|\nu|.
\]
前一种定义强调全变差控制，后一种表示则把复测度积分化为正测度积分。
\end{remark}

\begin{remark}[等价类没有点态支撑]
$L^1$ 元素是 a.e. 等价类。在零测稠密集上修改一个代表，可能彻底改变其点态非零集的闭包，
却不改变 $L^1$ 元素、范数或积分。因此，若要谈函数的拓扑支撑，必须指定代表；
若只谈 $L^1$ 结论，则应使用关于零集修改稳定的本质支撑概念。
\end{remark}

\begin{exercise}
证明 $\|f\|_1=\int|f|$ 在 $L^1(\mu)$ 上满足三角不等式，并说明为什么在
$\mathcal L^1(\mu)$ 上可能只有半范数。
\end{exercise}

\begin{exercise}
分别给出 $\R$ 上一个属于 $L^1_{\mathrm{loc}}\setminus L^1$ 的有界函数和一个无界函数，
并按定义验证。
\end{exercise}

\begin{exercise}
若 $f\in L^1$ 且 $E_n\downarrow\varnothing$，证明
$\int_{E_n}|f|\to0$。提示：对 $|f|\1_{E_n}$ 使用 DCT。
\end{exercise}

\begin{exercise}
证明定理~\ref{theorem:3-1-3-4} 中简单函数积分的表示无关；再证明若
$f_n\to f$ 于 $L^1(|\nu|)$，则 $\int f_n\,\dd\nu\to\int f\,\dd\nu$。
\end{exercise}

我们现在拥有两套互补规则：Tonelli 型问题应先保持非负，Fubini 型问题则要用绝对可积性控制抵消。
下一节把它们放到积空间中，并检验两种积分次序何时能够交换。

\section{积空间、换元与参数积分}\label{sec:03-02}

\subsection{Tonelli 与 Fubini：从矩形到迭代积分}\label{subsec:3-2-1}
\index{Tonelli 定理}\index{Fubini 定理}\index{截面}

设 $(X,\mathcal A,\mu)$ 与 $(Y,\mathcal B,\nu)$ 是两个测度空间。若 $E\subset X\times Y$，固定
$x\in X$ 或 $y\in Y$ 后得到的集合
\[
 E_x=\{y:(x,y)\in E\},\qquad E^y=\{x:(x,y)\in E\}
\]
称为 $E$ 的竖直截面与水平截面。对矩形 $E=A\times B$，截面计算没有任何悬念：
\[
 \nu(E_x)=\1_A(x)\nu(B),\qquad
 \int_X\nu(E_x)\,\dd\mu(x)=\mu(A)\nu(B).
 \tag{3.2.1}\label{eq:3-2-rectangle-slice}
\]
问题在于，乘积 $\sigma$-代数中的集合经过可数次补集与并运算后不再是矩形；即使每个截面都有
测度，也还要证明截面测度随参数可测，外层积分才有意义。

\begin{example}[矩形示性函数]\label{ex:3-2-1-1}
令 $f=\1_{A\times B}$。由 \eqref{eq:3-2-rectangle-slice}，
\[
 \int_X\left(\int_Y f(x,y)\,\dd\nu(y)\right)\dd\mu(x)
 =\mu(A)\nu(B)
 =\int_Y\left(\int_X f(x,y)\,\dd\mu(x)\right)\dd\nu(y).
\]
这项逐矩形计算是迭代积分公式的起点，但尚不能替代从矩形到任意乘积可测函数的扩张；其中各项
可以取 $+\infty$。
\end{example}

\begin{example}[对角线的两种答案]\label{ex:3-2-1-2}
在 $\R^2$ 中令 $D=\{(x,y):x=y\}$。对每个 $x$，$D_x=\{x\}$，故
$m_1(D_x)=0$。第~\ref{subsec:2-3-2} 小节已用盒覆盖独立证明 $m_2(D)=0$；本小节建立
Tonelli 定理后，截面计算会再次给出同一答案。盒覆盖论证与截面论证分别从外测度和乘积积分
解释同一个零测现象。
\end{example}

首先处理可测性。注意以下引理不给 $X$ 配测度；这一点以后构造密度核时很重要。

\begin{lemma}[非负参数积分的可测性]\label{proposition:3-2-1-3}
设 $(X,\mathcal A)$ 为可测空间，$(Y,\mathcal B,\nu)$ 为 $\sigma$-有限测度空间。若
$f:X\times Y\to[0,\infty]$ 对 $\mathcal A\otimes\mathcal B$ 可测，则每个截面
$f(x,\cdot)$ 可测，并且
\[
 x\longmapsto \int_Y f(x,y)\,\dd\nu(y)
\]
是 $\mathcal A$-可测函数。
\end{lemma}

\begin{proof}
固定 $x$。所有满足 $E_x\in\mathcal B$ 的集合 $E\subset X\times Y$ 构成一个
$\sigma$-代数，并包含每个可测矩形，因而包含 $\mathcal A\otimes\mathcal B$。应用于
$E=\{f>t\}$ 可知 $f(x,\cdot)$ 可测。

先设 $\nu(Y)<\infty$。令 $\mathcal D$ 为所有满足
$x\mapsto\nu(E_x)$ 可测的 $E\in\mathcal A\otimes\mathcal B$。若 $E=A\times B$，则
$\nu(E_x)=\1_A(x)\nu(B)$，所以可测矩形属于 $\mathcal D$。此外，
\[
 \nu((E^c)_x)=\nu(Y)-\nu(E_x),
\]
故 $\mathcal D$ 对相对于 $X\times Y$ 的补集封闭。若 $E_j\in\mathcal D$ 两两不交，则
\[
 \nu\!\left(\left(\bigcup_{j=1}^{\infty}E_j\right)_x\right)
 =\sum_{j=1}^{\infty}\nu((E_j)_x),
\]
右侧是可测非负函数部分和的点态极限。因此 $\mathcal D$ 是一个 $\lambda$-系。可测矩形构成
生成 $\mathcal A\otimes\mathcal B$ 的 $\pi$-系，$\pi$--$\lambda$ 定理给出
$\mathcal D=\mathcal A\otimes\mathcal B$。

一般地，取 $Y_k\uparrow Y$，其中 $Y_k\in\mathcal B$ 且 $\nu(Y_k)<\infty$。对任意
$E\in\mathcal A\otimes\mathcal B$，有限测度情形说明
$x\mapsto\nu(E_x\cap Y_k)$ 可测，而从下连续性给出
\[
 \nu(E_x)=\lim_{k\to\infty}\nu(E_x\cap Y_k).
\]
故集合示性函数的结论成立。非负简单函数的结论由有限线性组合得到。最后，若
$0\le s_n\uparrow f$ 是非负简单函数列，则对每个 $x$，单调收敛定理给出
\[
 \int_Ys_n(x,y)\,\dd\nu(y)\uparrow
 \int_Yf(x,y)\,\dd\nu(y).
\]
可测函数的递增极限可测，证明完毕。
\end{proof}

现在可以给外层积分赋予含义。对非负函数允许答案为 $+\infty$；这正是 Tonelli 定理与
Fubini 定理的分界。

\begin{definition}[乘积空间上的绝对可积]\label{def:3-2-1-2}
设 $f:X\times Y\to\R$ 或 $\C$ 可测。若
\[
 \int_{X\times Y}|f|\,\dd(\mu\otimes\nu)<\infty,
\]
则称 $f$ 在乘积空间上\emph{绝对可积}；等价地，$f\in L^1(\mu\otimes\nu)$。
\end{definition}

\begin{theorem}[Tonelli 定理]\label{theorem:3-2-1-1}
设 $\mu$ 与 $\nu$ 均为 $\sigma$-有限测度，$\mu\otimes\nu$ 是
$\mathcal A\otimes\mathcal B$ 上的乘积测度。若
$f:X\times Y\to[0,\infty]$ 可测，则两个截面积函数均可测，并且
\[
 \int_{X\times Y}f\,\dd(\mu\otimes\nu)
 =\int_X\!\left(\int_Yf(x,y)\,\dd\nu(y)\right)\dd\mu(x)
 =\int_Y\!\left(\int_Xf(x,y)\,\dd\mu(x)\right)\dd\nu(y).
 \tag{3.2.2}\label{eq:tonelli}
\]
三个量允许同时等于 $+\infty$。
\end{theorem}

\begin{proof}
引理~\ref{proposition:3-2-1-3} 已给出截面积函数的可测性。对
$E\in\mathcal A\otimes\mathcal B$ 定义
\[
 \lambda(E)=\int_X\nu(E_x)\,\dd\mu(x).
\]
若 $(E_j)$ 两两不交，则对每个 $x$，截面 $(E_j)_x$ 两两不交，因而
\[
 \nu\!\left(\left(\bigcup_jE_j\right)_x\right)
 =\lim_{N\to\infty}\sum_{j=1}^N\nu((E_j)_x).
\]
对右侧的递增部分和应用单调收敛定理，得到
$\lambda(\bigcup_jE_j)=\sum_j\lambda(E_j)$；又有 $\lambda(\varnothing)=0$，所以
$\lambda$ 是测度。对可测矩形，
\[
 \lambda(A\times B)=\int_X\1_A(x)\nu(B)\,\dd\mu(x)=\mu(A)\nu(B).
\]
取有限测度集的可数覆盖 $X=\bigcup_iX_i$、$Y=\bigcup_jY_j$。矩形
$X_i\times Y_j$ 可数地覆盖 $X\times Y$，且
$\lambda(X_i\times Y_j)=\mu(X_i)\nu(Y_j)<\infty$，所以 $\lambda$ 也是 $\sigma$-有限的。
乘积测度唯一性定理
\ref{theorem:2-3-2-2} 说明 $\lambda=\mu\otimes\nu$。因此
\[
 (\mu\otimes\nu)(E)=\int_X\nu(E_x)\,\dd\mu(x).
 \tag{3.2.3}\label{eq:cavalieri-set-pre}
\]

令 $f$ 为非负简单函数，把 \eqref{eq:cavalieri-set-pre} 对其有限个示性函数逐项相加，便得到
\eqref{eq:tonelli} 的第一个等号。一般 $f\ge0$ 可取非负简单函数
$s_n\uparrow f$。固定 $x$ 后，单调收敛定理给出
\[
 \int_Ys_n(x,y)\,\dd\nu(y)\uparrow \int_Yf(x,y)\,\dd\nu(y);
\]
再对 $x$ 应用单调收敛定理，并同时对积空间应用一次，得到
\[
 \int_X\int_Y f\,\dd\nu\,\dd\mu
 =\lim_n\int_X\int_Ys_n\,\dd\nu\,\dd\mu
 =\lim_n\int_{X\times Y}s_n\,\dd(\mu\otimes\nu)
 =\int_{X\times Y}f\,\dd(\mu\otimes\nu).
\]
交换 $X,Y$ 重复同一论证，得到另一个次序的等式。
\end{proof}

\begin{corollary}[Cavalieri 原理]\label{corollary:3-2-1-4}
对每个 $E\in\mathcal A\otimes\mathcal B$，
\[
 (\mu\otimes\nu)(E)=\int_X\nu(E_x)\,\dd\mu(x)
 =\int_Y\mu(E^y)\,\dd\nu(y).
\]
\end{corollary}

\begin{proof}
在 Tonelli 定理中取 $f=\1_E$。
\end{proof}

\begin{theorem}[Fubini 定理]\label{theorem:3-2-1-2}
在 Tonelli 定理的假设下，设 $f:X\times Y\to\R$ 或 $\C$ 在乘积空间上绝对可积。则
\[
 \int_Y|f(x,y)|\,\dd\nu(y)<\infty
\]
对 $\mu$-几乎处处的 $x$ 成立；交换 $X,Y$ 后的陈述也成立。把异常参数处的截面积分置为零，
所得函数可积，并且
\[
 \int_{X\times Y}f\,\dd(\mu\otimes\nu)
 =\int_X\!\left(\int_Y f(x,y)\,\dd\nu(y)\right)\dd\mu(x)
 =\int_Y\!\left(\int_X f(x,y)\,\dd\mu(x)\right)\dd\nu(y).
 \tag{3.2.4}\label{eq:fubini}
\]
\end{theorem}

\begin{proof}
对 $|f|$ 应用 Tonelli 定理。函数
\[
 h(x)=\int_Y|f(x,y)|\,\dd\nu(y)
\]
非负可测，并满足
\[
 \int_Xh\,\dd\mu=\int_{X\times Y}|f|\,\dd(\mu\otimes\nu)<\infty.
\]
若 $A=\{h=\infty\}$ 有正测度，则对每个 $M>0$ 有
$\int h\,\dd\mu\ge M\mu(A)$；当 $\mu(A)>0$ 时令 $M\to\infty$ 会与积分有限矛盾。
故 $h<\infty$ 几乎处处。交换变量得到另一个方向的截面绝对可积性。

若 $f$ 为实值，对 $f^+$ 与 $f^-$ 分别应用 Tonelli。由于
$0\le f^\pm\le|f|$，相应迭代积分都是有限数，于是可以相减，得到
\eqref{eq:fubini}。若 $f$ 为复值，则对 $\Re f$ 与 $\Im f$ 重复实值论证；二者的绝对值都不超过
$|f|$，所以全部相减均在有限数之间进行。异常集上的置零只改变零集上的函数，不改变外层积分。
\end{proof}

\begin{corollary}[可积参数积分的可测性]\label{corollary:3-2-1-integrable-parameter}
在 Fubini 定理的假设下，把不可积截面处的值置零后，
\[
 x\longmapsto\int_Yf(x,y)\,\dd\nu(y)
\]
可测且属于 $L^1(\mu)$。
\end{corollary}

\begin{proof}
对实值函数，把引理~\ref{proposition:3-2-1-3} 应用于 $(\Re f)^\pm$；对复值函数再处理
$(\Im f)^\pm$。这些非负参数积分均可测，在几乎处处截面可积的集合上作差即得所需函数。
置零集合由上一证明中的可测函数 $h$ 的无穷值集给出，故仍可测。此外
\[
 \left|\int_Yf(x,y)\,\dd\nu(y)\right|
 \le \int_Y|f(x,y)|\,\dd\nu(y)=h(x)
\]
几乎处处成立，而 $h\in L^1(\mu)$。
\end{proof}

\begin{remark}[完成化上的代表]
若 $f$ 只对乘积测度的完成化可测，应先选取一个
$\mathcal A\otimes\mathcal B$-可测代表 $\widetilde f$。若两个这样的代表只在乘积零集 $N$ 上
不同，取乘积可测零集 $Z\supset N$；Tonelli 应用于 $\1_Z$ 给出
$\nu(Z_x)=0$ 对几乎处处 $x$ 成立。交换 $X,Y$ 后再次应用 Tonelli，又得到
$\mu(Z^y)=0$ 对几乎处处 $y$ 成立。因此两代表给出相同的几乎处处截面积分。
任意完成化可测代表的每一个截面却未必可测；Fubini 定理并不作这个更强的断言。
\end{remark}

现在回看几个计算。对非负可测 $u:X\to[0,\infty]$，令
\[
 G_u=\{(x,t)\in X\times(0,\infty):0<t<u(x)\}.
\]
集合恒等式
\[
 G_u=\bigcup_{q\in\Q,\ q>0}\{u>q\}\times(0,q)
\]
证明它乘积可测；其竖直截面长度为 $u(x)$。故 Cavalieri 原理给出
\[
 (\mu\otimes m_1)(G_u)=\int_Xu\,\dd\mu.
 \tag{3.2.5}\label{eq:subgraph-area}
\]
简单函数时，右侧是有限个矩形面积之和；一般情形则是这些阶梯下图的单调极限。

\begin{figure}[htbp]
\centering
\begin{tikzpicture}[x=1.05cm,y=.75cm]
  \draw[->,book line] (-.2,0)--(6.2,0) node[right]{$x$};
  \draw[->,book line] (0,-.15)--(0,4.4) node[above]{$t$};
  \fill[BookAccent!10] (0,0)--(0,1.2)--(1.15,1.2)--(1.15,2.4)--(2.3,2.4)--
    (2.3,3.5)--(3.45,3.5)--(3.45,2.75)--(4.6,2.75)--(4.6,1.65)--(5.75,1.65)--(5.75,0)--cycle;
  \draw[BookAccent,semithick] (0,1.2)--(1.15,1.2)--(1.15,2.4)--(2.3,2.4)--
    (2.3,3.5)--(3.45,3.5)--(3.45,2.75)--(4.6,2.75)--(4.6,1.65)--(5.75,1.65);
  \foreach \x in {1.15,2.3,3.45,4.6,5.75}{\draw[BookInk!35,dashed] (\x,0)--(\x,3.8);}
  \draw[BookWarm,very thick] (2.85,0)--(2.85,3.5);
  \node[book label,right] at (2.85,1.7) {截面长度 $u(x)$};
  \draw[BookAccent!75,very thick,dashed] (0,2.1)--(4.6,2.1);
  \node[book label,above] at (4.9,2.1) {交换次序后的水平截面};
\end{tikzpicture}
\caption{竖直截面给 $\int_X\int_Y$；把同一区域改用水平截面读取，就得到相反的积分次序。}
\label{fig:3-2-1-slices}
\end{figure}

\begin{example}[换序为何会改变答案]\label{ex:3-2-1-3}
在 $\N^2$ 上取计数测度，令
\[
 a_{mn}=\begin{cases}
 1,&n=m,\\
 -1,&n=m+1,\\
 0,&\text{其他}.
 \end{cases}
\]
固定 $m$ 时，一行只有 $1,-1$ 两项，故
$\sum_n a_{mn}=0$，从而 $\sum_m(\sum_na_{mn})=0$。固定 $n$ 时，$n=1$ 的列和为 $1$；
若 $n\ge2$，则 $m=n$ 给 $1$，$m=n-1$ 给 $-1$，列和为零。因此
\[
 \sum_n\left(\sum_ma_{mn}\right)=1.
\]
另一方面，每行的绝对值之和为 $2$，故
$\sum_{m,n}|a_{mn}|=\infty$。这里没有违反 Fubini：失败的正是绝对可积假设。Tonelli 也不能
用于 $a_{mn}$ 本身，因为它有正有负；分别用于正负部分只会得到 $+\infty$ 与 $+\infty$，二者
不能相减。
\end{example}

\begin{example}[对角线的截面计算]\label{ex:3-2-1-diagonal-recheck}
对 $D=\{(x,y):x=y\}$ 应用 Cavalieri 原理。每个 $D_x$ 是单点，故
\[
 m_2(D)=\int_\R m_1(\{x\})\,\dd x=0.
\]
这与第~\ref{subsec:2-3-2} 小节的覆盖计算一致，并从截面长度恒为零解释了二维零测性。
\end{example}

最后考虑尚未完成的高斯计算。记
$I=\int_\R e^{-x^2}\,\dd x$。非负性与 Tonelli 给出
\[
 I^2=\int_{\R^2}e^{-(x^2+y^2)}\,\dd(x,y),
 \tag{3.2.6}\label{eq:gaussian-tonelli-gap}
\]
但矩形切片并未利用被积函数的圆对称性，所以此处还不能求出 $I$。下一小节将把圆切开并拉直，
从而补上这一步。

同一机制还支撑卷积与 Fourier 积分的换序：只有当被积函数非负，或其绝对值在乘积空间可积时，
换序才由本小节的定理保证。把联合推前测度恰好等于两个边缘推前测度之积时，乘积型积分也直接来自
Tonelli；第~5~章会把这一条件纳入随机变量独立性的正式定义。事件示性函数的非负级数同样可在此
逐项积分；这正是第一 Borel--Cantelli 估计的积分核心。向量值积分的换序还需要强可测性与
范数可积，不能仅由标量 Fubini 推出。

\begin{exercise}
用 Tonelli 定理证明：若 $u,v\ge0$ 可测，则
$\int_{X\times Y}u(x)v(y)\,\dd(\mu\otimes\nu)=\int_Xu\,\dd\mu\int_Yv\,\dd\nu$，
其中乘积按非负扩展实数解释。
\end{exercise}
\begin{exercise}
设 $(S,\Sigma,\lambda)$ 是总质量为一的测度空间，$\xi,\eta:S\to[0,\infty]$ 可测，并且联合映射
$(\xi,\eta)$ 的推前测度满足
\[
  (\xi,\eta)_*\lambda=(\xi_*\lambda)\otimes(\eta_*\lambda).
\]
由推前积分公式与 Tonelli 定理推出
\[
  \int_S\xi\eta\,\dd\lambda
  =\left(\int_S\xi\,\dd\lambda\right)
   \left(\int_S\eta\,\dd\lambda\right).
\]
等式允许取值 $+\infty$。第~5~章才会把假设和结论改写为独立随机变量的语言。
\end{exercise}
\begin{exercise}
逐项核算例~\ref{ex:3-2-1-3} 的两个迭代和，并说明任意有限矩形部分和为什么不产生矛盾。
\end{exercise}
\begin{exercise}
对对角线 $D=\{(x,y):x=y\}$ 直接应用 Tonelli 定理，证明 $m_2(D)=0$；明确写出每个竖直截面，
并说明为何不需要把“不可数个零测截面的并仍为零测”当作一般原则。
\end{exercise}
\begin{exercise}
从 Tonelli 定理取 $f=\1_E$ 推导 Cavalieri 公式，再对非负简单函数的下图集逐层核对
\eqref{eq:subgraph-area}。
\end{exercise}
\begin{exercise}
设 $(\Omega,\mathcal F,\Prob)$ 是总质量为一的测度空间，$E_n\in\mathcal F$，且
$\sum_n\Prob(E_n)<\infty$。对
$\sum_n\1_{E_n}$ 使用 Tonelli，证明几乎每个点只属于有限多个 $E_n$；这就是第一
Borel--Cantelli 引理的测度论形式。
\end{exercise}

式 \eqref{eq:gaussian-tonelli-gap} 已把一维问题提升成二维问题。提升维数只有在我们能够使用圆对称时
才有收益；这要求一个严格的非线性换元公式。

\subsection{\texorpdfstring{$\R^n$}{R n} 中的换元公式}\label{subsec:3-2-2}
\index{换元公式}\index{Jacobian}

初等微积分常把极坐标写成
\[
 x=r\cos\theta,\qquad y=r\sin\theta,
\]
然后在面积元旁添上一个 $r$。这个因子并不是记忆口诀：半径为 $r$、径向厚度为
$\dd r$、角宽为 $\dd\theta$ 的小格近似一个边长为 $\dd r$ 与 $r\dd\theta$ 的矩形。
在一般坐标中，同一局部线性化由行列式记录。

局部线性化却不会自动消除全局重叠。映射
$\Psi(x,t)=(e^x\cos t,e^x\sin t)$ 的 Jacobian 处处非零，但 $t$ 与 $t+2\pi$ 被送到同一点；
若不先切开角变量，同一块面积会被反复计算。下面的主定理因此从真正的全局微分同胚开始，定理后
再把这个周期例完整算一遍，并说明何时可以按单射分支求和。

\begin{definition}[Jacobian]\label{def:3-2-2-1}
设 $U\subset\R^n$ 开，$\Phi:U\to\R^n$ 为 $C^1$ 映射。定义
\[
 J_\Phi(x)=|\det D\Phi(x)|,
\]
称为 $\Phi$ 在 $x$ 处的 Jacobian。绝对值不可省略：测度记录无向体积；只有讨论有向体积时才保留
$\det D\Phi$ 的符号。
\end{definition}

若 $\Phi:U\to V$ 是 $C^1$ 微分同胚，我们希望证明的测度关系可以先写成
\[
 \Phi_*\bigl(J_\Phi m_n|_U\bigr)=m_n|_V.
 \tag{3.2.7}\label{eq:cov-pushforward-target}
\]
这不是新定义，而是下文集合公式的等价目标。

\begin{remark}[目标推前关系]\label{def:3-2-2-2}
式 \eqref{eq:cov-pushforward-target} 的集合版本是
$m_n(\Phi(E))=\int_EJ_\Phi$。先证明集合版本，再按示性函数、简单函数和单调极限扩张，
可以避免把形式上的“$\dd y=J_\Phi(x)\dd x$”当成证明。
\end{remark}

第~\ref{subsec:2-4-2} 小节已经证明线性映射 $A$ 满足
$m_n(AE)=|\det A|m_n(E)$。下一引理说明，当 $D\Phi$ 在一个小立方体上接近 $A$ 时，
$\Phi$ 对每个可测子集的体积作用也一致接近 $A$。这个“对每个子集”比只计算整个立方体强，
它使后面的有限分割可以逐块相加。

\begin{lemma}[局部体积估计]\label{lemma:3-2-2-1}
给定可逆矩阵 $A$ 与 $0<\varepsilon<1$，存在 $\delta>0$，使得下述结论成立。设 $Q$ 是闭立方体，
$\Phi$ 在 $Q$ 的一个邻域上为 $C^1$，在 $Q$ 上单射，且
\[
 \sup_{x\in Q}\|D\Phi(x)-A\|<\delta.
\]
则对每个 Lebesgue 可测集 $E\subset Q$，$\Phi(E)$ 可测，并且
\[
 (1-\varepsilon)|\det A|m_n(E)
 \le m_n(\Phi(E))
 \le(1+\varepsilon)|\det A|m_n(E).
 \tag{3.2.8}\label{eq:local-volume-estimate}
\]
\end{lemma}

\begin{proof}
令 $\Psi=A^{-1}\Phi$。若
$\eta:=\sup_Q\|D\Psi-I\|<1$，则立方体的凸性和线段积分给出
\[
 \begin{aligned}
 |\Psi(x)-\Psi(y)-(x-y)|
 &=\left|\int_0^1(D\Psi(y+t(x-y))-I)(x-y)\,\dd t\right|\\
 &\le\eta|x-y|,
 \end{aligned}
\]
所以
\[
 (1-\eta)|x-y|\le|\Psi(x)-\Psi(y)|\le(1+\eta)|x-y|.
 \tag{3.2.9}\label{eq:near-id-bilip}
\]
特别地，$\Psi$ 在 $Q$ 上为双 Lipschitz 映射到其像。

我们需要一个当 $\eta\downarrow0$ 时趋于 $1$ 的体积界，因而在此把盒覆盖估计写得更精细。
写 $\Psi(z)=z+R(z)$，则 $R$ 是 $\eta$-Lipschitz。若 $C$ 是边长 $a$ 的坐标立方体，
对 $C\cap Q$ 中任意两点 $x,y$，每个坐标满足
\[
 |\Psi_i(x)-\Psi_i(y)|
 \le |x_i-y_i|+|R(x)-R(y)|
 \le a+\eta\sqrt n\,a.
\]
故 $\Psi(C\cap Q)$ 包含在各边长均不超过 $(1+\sqrt n\eta)a$ 的闭坐标盒中；该闭盒
Lebesgue 可测，且测度不超过 $(1+\sqrt n\eta)^na^n$。更精确地，给定
$\varepsilon_0>0$；因 $E\subset Q$，有 $m_n(E)\le m_n(Q)<\infty$。按外测度定义取立方体覆盖
$E\subset\bigcup_r C_r$，记 $C_r$ 边长为 $a_r$，使
\[
 \sum_r a_r^n<m_n(E)+\varepsilon_0.
\]
对每个 $r$，上述坐标振幅估计把 $\Psi(C_r\cap Q)$ 放入一个各边长不超过
$(1+\sqrt n\eta)a_r$ 的闭坐标盒，因而由外测度的单调性与闭盒体积公式，
\[
 m_n^*(\Psi(E))
 \le (1+\sqrt n\eta)^n\sum_r a_r^n
 <(1+\sqrt n\eta)^n\bigl(m_n(E)+\varepsilon_0\bigr).
\]
令 $\varepsilon_0\downarrow0$ 得到
\[
 m_n^*(\Psi(E))\le(1+\sqrt n\eta)^n m_n(E).
 \tag{3.2.10}\label{eq:near-id-upper}
\]
由 \eqref{eq:near-id-bilip}，对 $u=\Psi(x),v=\Psi(y)$，
\[
 \begin{aligned}
 |\Psi^{-1}(u)-\Psi^{-1}(v)-(u-v)|
 &=|R(x)-R(y)|\\
 &\le \frac{\eta}{1-\eta}|u-v|.
 \end{aligned}
\]
这一估计不要求 $\Psi(Q)$ 为凸集。确实，若坐标立方体 $C$ 的边长为 $a$，则对
$u,v\in C\cap\Psi(Q)$，上式说明每个坐标都满足
\[
 \bigl|(\Psi^{-1}(u))_i-(\Psi^{-1}(v))_i\bigr|
 \le a+\frac{\eta}{1-\eta}\sqrt n\,a.
\]
因此 $\Psi^{-1}(C\cap\Psi(Q))$ 包含于各边长均不超过
$\bigl(1+\sqrt n\eta/(1-\eta)\bigr)a$ 的闭坐标盒中。用可数立方体覆盖
$\Psi(E)$，若这些立方体边长为 $a_r$，则对对应坐标盒体积求和给出
\[
 m_n(E)
 \le \left(1+\frac{\sqrt n\eta}{1-\eta}\right)^n\sum_r a_r^n.
\]
再对 $\Psi(E)$ 的所有立方体覆盖取下确界，得到外测度估计
\[
 m_n(E)\le
 \left(1+\frac{\sqrt n\eta}{1-\eta}\right)^n m_n^*(\Psi(E)).
 \tag{3.2.11}\label{eq:near-id-lower}
\]
这里的可测性也需说明。$Q$ 上连续单射到 Hausdorff 空间的映射是到其像的同胚。而且
$\Psi(Q)$ 是紧集，因而在 $\R^n$ 中闭；所以相对 Borel 集同时也是 $\R^n$ 中的
Borel 集，特别地 Borel 子集的像是 Borel 集。任意 Lebesgue 可测 $E\subset Q$ 与某个
Borel 集 $B\subset Q$ 只差一个零集。由
\eqref{eq:near-id-upper}，该零集的像外测度为零；又
$\Psi(E)\mathbin\triangle\Psi(B)\subset\Psi(E\mathbin\triangle B)$，故 $\Psi(E)$ 可测。
于是 \eqref{eq:near-id-lower} 右端的外测度就是 $m_n(\Psi(E))$。

由线性体积公式，
$m_n(\Phi(E))=|\det A|m_n(\Psi(E))$。取 $0<\eta<1$ 足够小，使
\[
 (1+\sqrt n\eta)^n\le1+\varepsilon,
 \qquad
 \left(1+\frac{\sqrt n\eta}{1-\eta}\right)^{-n}\ge1-\varepsilon.
\]
再取 $0<\delta\le\eta/\|A^{-1}\|$。由
\[
 \sup_Q\|D\Psi-I\|
 \le\|A^{-1}\|\sup_Q\|D\Phi-A\|<\eta,
\]
式 \eqref{eq:near-id-upper}--\eqref{eq:near-id-lower} 依次给出
\[
 (1-\varepsilon)m_n(E)\le m_n(\Psi(E))
 \le(1+\varepsilon)m_n(E).
\]
乘以 $|\det A|$ 即得 \eqref{eq:local-volume-estimate}。
\end{proof}

\begin{theorem}[集合换元公式]\label{theorem:3-2-2-2}
设 $U,V\subset\R^n$ 为开集，$\Phi:U\to V$ 为 $C^1$ 微分同胚。若
$E\subset U$ Lebesgue 可测，则 $\Phi(E)$ Lebesgue 可测，且
\[
 m_n(\Phi(E))=\int_EJ_\Phi(x)\,\dd x.
 \tag{3.2.12}\label{eq:set-change-variables}
\]
\end{theorem}

\begin{proof}
先处理可测性。若 $B\subset U$ 为 Borel 集，则
$\Phi(B)=(\Phi^{-1})^{-1}(B)$ 是 $V$ 中的 Borel 集。开集 $U$ 可由可数多个闭立方体
$Q_j\Subset U$ 覆盖；$D\Phi$ 在每个 $Q_j$ 上有界，线段积分说明
$\Phi|_{Q_j}$ 为 Lipschitz 映射。因此，若 $Z\subset U$ 是 Lebesgue 零集，定理
\ref{theorem:2-4-3-2} 给出
\[
 \Phi(Z)=\bigcup_j\Phi(Z\cap Q_j)
\]
仍为零集。任意 Lebesgue 可测 $E\subset U$ 可写成
$E\mathbin\triangle B\subset Z$，其中 $B\subset U$ 为 Borel 集、$Z\subset U$ 为 Borel 零集。于是
$\Phi(E)\mathbin\triangle\Phi(B)\subset\Phi(Z)$，故 $\Phi(E)$ Lebesgue 可测。

下面证明等式。先固定紧集 $K\subset U$ 与 $0<\varepsilon<1$。对每个 $a\in K$，矩阵
$A_a=D\Phi(a)$ 可逆。由 $D\Phi$ 与 $J_\Phi$ 的连续性，结合引理
\ref{lemma:3-2-2-1}，可取以 $a$ 为内点的闭立方体 $Q_a\subset U$，使对每个可测
$F\subset Q_a$ 都有
\[
 (1-\varepsilon)\int_FJ_\Phi\,\dd x
 \le m_n(\Phi(F))
 \le(1+\varepsilon)\int_FJ_\Phi\,\dd x.
 \tag{3.2.13}\label{eq:local-integral-sandwich}
\]
为说明这一步的常数选择，先取 $0<\rho<1$，使
\[
 \frac{1-\rho}{1+\rho}\ge1-\varepsilon,
 \qquad
 \frac{1+\rho}{1-\rho}\le1+\varepsilon.
\]
令引理中的相对误差为 $\rho$，再缩小 $Q_a$ 使
$J_\Phi(x)/J_\Phi(a)\in[1-\rho,1+\rho]$。于是对 $F\subset Q_a$，
\[
 \begin{aligned}
 m_n(\Phi(F))
 &\ge(1-\rho)J_\Phi(a)m_n(F)
 \ge\frac{1-\rho}{1+\rho}\int_FJ_\Phi\,\dd x,\\
 m_n(\Phi(F))
 &\le(1+\rho)J_\Phi(a)m_n(F)
 \le\frac{1+\rho}{1-\rho}\int_FJ_\Phi\,\dd x,
 \end{aligned}
\]
这便给出 \eqref{eq:local-integral-sandwich}。

$K$ 的紧性给出有限个 $Q_1,\dots,Q_N$，其内部覆盖 $K$。令
\[
 P_1=K\cap Q_1,
 \qquad
 P_j=(K\cap Q_j)\setminus\bigcup_{i<j}P_i\quad(2\le j\le N).
\]
这些是两两不交的可测集，覆盖 $K$，且 $P_j\subset Q_j$。若 $E\subset K$ 可测，令
$E_j=E\cap P_j$。由于 $\Phi$ 单射，$\Phi(E_j)$ 也两两不交。对
\eqref{eq:local-integral-sandwich} 有限求和，得到
\[
 (1-\varepsilon)\int_EJ_\Phi\,\dd x
 \le m_n(\Phi(E))
 \le(1+\varepsilon)\int_EJ_\Phi\,\dd x.
\]
紧集上 $J_\Phi$ 有界且 $m_n(K)<\infty$，故中间积分有限。令
$\varepsilon\downarrow0$，得到 \eqref{eq:set-change-variables} 对每个 $E\subset K$ 成立。

最后取递增紧耗尽
\[
 K_j=\{x\in U:|x|\le j,\ \dist(x,U^c)\ge j^{-1}\},
\]
当 $U=\R^n$ 时约定 $\dist(x,\varnothing)=\infty$。则 $K_j\uparrow U$。对
$E\cap K_j$ 使用已证等式；由于 $\Phi$ 单射，$\Phi(E\cap K_j)\uparrow\Phi(E)$。
测度的从下连续性与单调收敛定理分别给出
\[
 m_n(\Phi(E))
 =\lim_jm_n(\Phi(E\cap K_j))
 =\lim_j\int_{E\cap K_j}J_\Phi\,\dd x
 =\int_EJ_\Phi\,\dd x.
\]
证明完成。局部单射用于引理，整体单射则保证分块的像不重叠；这两个假设承担的作用不同。
\end{proof}

\begin{theorem}[函数换元公式]\label{theorem:3-2-2-3}
在定理~\ref{theorem:3-2-2-2} 的假设下，若 $g:V\to[0,\infty]$ 可测，则
\[
 \int_Vg(y)\,\dd y
 =\int_Ug(\Phi(x))J_\Phi(x)\,\dd x.
 \tag{3.2.14}\label{eq:function-change-variables}
\]
若 $g$ 为实值或复值且任一侧绝对可积，同一等式仍成立，并且另一侧也绝对可积。
\end{theorem}

\begin{proof}
先补上复合函数的可测性。若 $B\subset V$ Lebesgue 可测，把已经证明的集合换元定理应用于
逆微分同胚 $\Phi^{-1}:V\to U$，便知 $\Phi^{-1}(B)$ Lebesgue 可测。因此，对任意 Lebesgue
可测的扩展实值函数 $g$ 以及任意 $a\in\R$，
\[
 \{x\in U:g(\Phi(x))>a\}
 =\Phi^{-1}(\{y\in V:g(y)>a\})
\]
可测；故 $g\circ\Phi$ 可测。又 $J_\Phi$ 连续，所以在 $g\ge0$ 时
$g\circ\Phi\,J_\Phi$ 是非负可测函数。

当 $g=\1_B$ 且 $B\subset V$ 可测时，集合换元公式应用于
$E=\Phi^{-1}(B)$ 给出
\[
 \int_U\1_B(\Phi(x))J_\Phi(x)\,\dd x
 =m_n(\Phi(\Phi^{-1}B))=m_n(B).
\]
有限非负线性组合给出简单函数情形。对一般 $g\ge0$ 取简单函数
$s_k\uparrow g$，则 $s_k\circ\Phi\,J_\Phi\uparrow g\circ\Phi\,J_\Phi$；两侧分别应用
单调收敛定理得到 \eqref{eq:function-change-variables}。实值可积函数拆成正负部分，复值函数再拆成
实部与虚部。对 $|g|$ 的非负公式同时说明两侧的绝对可积性等价。
\end{proof}

\begin{example}[极坐标]\label{ex:3-2-2-1}
映射
\[
 \Phi:(0,\infty)\times(0,2\pi)\longrightarrow
 \R^2\setminus\{(r,0):r\ge0\},
 \qquad \Phi(r,\theta)=(r\cos\theta,r\sin\theta)
\]
是 $C^1$ 微分同胚，并且
\[
 D\Phi(r,\theta)=
 \begin{pmatrix}
 \cos\theta&-r\sin\theta\\
 \sin\theta&r\cos\theta
 \end{pmatrix},
 \qquad J_\Phi(r,\theta)=r.
\]
被切去的闭射线是 Lebesgue 零集，故对任意非负可测或可积 $g$，
\[
 \int_{\R^2}g(x,y)\,\dd x\dd y
 =\int_0^{2\pi}\int_0^\infty
 g(r\cos\theta,r\sin\theta)r\,\dd r\dd\theta.
 \tag{3.2.15}\label{eq:polar-coordinates}
\]
切开射线不是审美选择：它消除了角变量的周期重叠，使映射真正单射。
\end{example}

\begin{example}[高斯积分的回收]\label{ex:3-2-2-gaussian-recovery}
由式 \eqref{eq:gaussian-tonelli-gap} 与极坐标公式，
\[
 \begin{aligned}
 I^2
 &=\int_0^{2\pi}\int_0^\infty e^{-r^2}r\,\dd r\dd\theta\\
 &=2\pi\left[-\frac12e^{-r^2}\right]_{0}^{\infty}=\pi.
 \end{aligned}
\]
第一行中的平方只用 Tonelli；把圆对称变成一维积分才使用换元。又因
$I\ge\int_{-1}^1e^{-1}\,\dd x>0$，所以
\[
 \int_\R e^{-x^2}\,\dd x=\sqrt\pi.
\]
\end{example}

\begin{example}[球坐标的径向部分]\label{ex:3-2-2-2}
记 $\omega_n=m_n(B(0,1))$，并令 $R(x)=|x|$。线性体积公式给出
$m_n(B(0,r))=\omega_nr^n$。当 $r=0$ 时，球面是零测单点。若 $r>0$，则对
$0<\delta<r$，球面包含在
$B(0,r+\delta)\setminus B(0,r-\delta)$ 中；该环带测度为
$\omega_n((r+\delta)^n-(r-\delta)^n)$，令 $\delta\downarrow0$ 得球面测度为零。
因此半径推前测度 $R_*m_n$ 满足
\[
 (R_*m_n)((a,b])=\omega_n(b^n-a^n)\qquad(0\le a<b).
\]
另一方面，对 $0\le a<b$ 有
\[
 \int_{(a,b]}n\omega_nr^{n-1}\,\dd r
 =\omega_n[b^n-a^n],
\]
且测度 $n\omega_nr^{n-1}\,\dd r$ 与 $R_*m_n$ 赋予 $\{0\}$ 的质量都为零。两者在
$[0,N]$ 上都有限，而 $\{0\}$ 与半开区间 $(a,b]$ 构成生成 $[0,\infty)$ Borel
$\sigma$-代数的 $\pi$-系统；因此测度的 $\sigma$-有限唯一性（也可等价地应用定理
\ref{theorem:2-3-3-2} 的 Lebesgue--Stieltjes 唯一性）给出
\[
 R_*m_n=n\omega_nr^{n-1}\,\dd r.
\]
于是对任意非负
可测函数 $f:[0,\infty)\to[0,\infty]$，以及任一侧绝对可积的实值或复值 $f$，推前积分公式给出
\[
 \int_{\R^n}f(|x|)\,\dd x
 =n\omega_n\int_0^\infty f(r)r^{n-1}\,\dd r.
 \tag{3.2.16}\label{eq:radial-integration}
\]
常数 $|S^{n-1}|:=n\omega_n$ 只是方便记号；这里没有构造曲面测度，也没有使用一般面积公式。
\end{example}

\begin{example}[密度变换]\label{ex:3-2-2-3}
设 $U$ 上的概率测度为 $\lambda=p\,m_n|_U$，并令 $\kappa=\Phi_*\lambda$。对 Borel 集
$B\subset V$，函数换元公式应用于 $\Phi^{-1}:V\to U$，给出
\[
 \begin{aligned}
 \kappa(B)
 &=\int_{\Phi^{-1}(B)}p(x)\,\dd x\\
 &=\int_Bp(\Phi^{-1}(y))J_{\Phi^{-1}}(y)\,\dd y.
 \end{aligned}
\]
逆函数定理给出
$D\Phi^{-1}(y)=(D\Phi(\Phi^{-1}y))^{-1}$，所以
\[
 p_\kappa(y)=
 \begin{cases}
 p(\Phi^{-1}(y))|\det D\Phi^{-1}(y)|,&y\in V,\\
 0,&y\notin V,
 \end{cases}
\]
是 $\kappa$ 在全 $\R^n$ 上的密度。若把 $\lambda$ 看作随机向量 $X$ 的分布，则 $\kappa$ 就是
$\Phi(X)$ 的分布。密度公式只确定 a.e. 等价类，在零集上修改密度不改变推前测度。
\end{example}

\begin{example}[非零 Jacobian 不保证全局单射]\label{ex:3-2-2-4}
令
\[
 \Psi(x,t)=(e^x\cos t,e^x\sin t).
\]
直接计算得 $\det D\Psi(x,t)=e^{2x}>0$，但
$\Psi(x,t+2\pi)=\Psi(x,t)$。因此处处非退化只由逆函数定理保证局部单射，不能推出全局单射。
把 $t$ 限制在任一长度为 $2\pi$ 的开区间，才得到去掉一条射线的平面坐标图；若保留所有
$t$，每个目标点会被无穷多次计算。
\end{example}

对有限或可数个已知单射分支，可以直接计算每个目标点的覆盖重数。

\begin{remark}[一般重数公式的边界]\label{def:3-2-2-3}
对一般 Lipschitz 映射，形式上常写
$N(\Phi,y)=\#\Phi^{-1}\{y\}$。一般重数的可测性、近似微分与面积公式需要几何测度论。
以下命题处理可数个明确的 $C^1$ 微分同胚分支；在这一情形，重数是可测示性函数的可数和。
\end{remark}

\begin{proposition}[分片单射版本]\label{proposition:3-2-2-4}
设 $U\subset\R^n$ 开，$\Phi:U\to\R^n$ 为 $C^1$。设 $N\subset U$ 是 Lebesgue 零集，且
\[
 U\setminus N=\bigsqcup_{j=1}^{\infty}E_j,
\]
其中 $E_j$ Lebesgue 可测，并包含于开集 $U_j\subset U$，而
$\Phi|_{U_j}$ 是到开像的 $C^1$ 微分同胚。则每个 $\Phi(E_j)$ 可测；对非负可测 $g$，
\[
 \int_Ug(\Phi(x))J_\Phi(x)\,\dd x
 =\int_{\R^n}g(y)\sum_{j=1}^{\infty}\1_{\Phi(E_j)}(y)\,\dd y.
 \tag{3.2.17}\label{eq:piecewise-change-variables}
\]
实值或复值情形在任一侧绝对可积时同样成立。若各 $\Phi(E_j)$ 两两几乎不交，右侧就是
$\int_{\Phi(U)}g$。
\end{proposition}

\begin{proof}
集合换元公式的可测性部分应用于 $\Phi|_{U_j}$，说明 $\Phi(E_j)$ 可测。又因 $N$ 零测，
每个 $E_j$ 上的 $g\circ\Phi\,J_\Phi$ 可测。可数个可测分支给出它在 $U\setminus N$ 上的
可测性；而 $N$ 的任意子集都 Lebesgue 可测，所以把它补到 $N$ 上仍是可测函数，且在 $N$
上的积分为零。对每个 $j$，在函数换元公式中取
$g\1_{\Phi(E_j)}$，得到
\[
 \int_{E_j}g(\Phi(x))J_\Phi(x)\,\dd x
 =\int_{\R^n}g(y)\1_{\Phi(E_j)}(y)\,\dd y.
\]
对 $j$ 求和，并在两侧对非负部分和使用单调收敛定理，即得
\eqref{eq:piecewise-change-variables}。带符号或复值情形对绝对值先用非负公式，再拆正负部或实虚部。

还可核对异常像确实不会藏入额外质量。用可数个相对紧凸邻域覆盖 $U$，$D\Phi$ 在每个较小闭块
上有界，故 $\Phi$ 在各块上 Lipschitz。定理~\ref{theorem:2-4-3-2} 与可数并给出
$m_n(\Phi(N))=0$。这一步不声称 $\Phi$ 在整个无界开集上一致 Lipschitz。
\end{proof}

\begin{figure}[htbp]
\centering
\begin{tikzpicture}[x=.8cm,y=.8cm]
  \begin{scope}
    \draw[book line] (0,0) rectangle (4,3);
    \foreach \x in {1,2,3}{\draw[BookInk!45] (\x,0)--(\x,3);}
    \foreach \y in {1,2}{\draw[BookInk!45] (0,\y)--(4,\y);}
    \node[book label] at (2,-.38) {参数域中的小格};
  \end{scope}
  \draw[book accent arrow] (4.45,1.5)--(5.55,1.5) node[midway,above]{$\Phi$};
  \begin{scope}[shift={(6,0)}]
    \draw[book line] plot[smooth] coordinates{(0,.2)(1.1,0)(2.4,.2)(4,.65)};
    \draw[book line] plot[smooth] coordinates{(.1,3)(1.2,2.85)(2.7,3.15)(4.25,3.5)};
    \draw[BookInk!45] plot[smooth] coordinates{(.05,.18)(.15,1.2)(.12,2.2)(.1,3)};
    \draw[BookInk!45] plot[smooth] coordinates{(1.05,.28)(1.15,1.33)(1.12,2.36)(1.1,3.12)};
    \draw[BookInk!45] plot[smooth] coordinates{(2.05,.38)(2.15,1.46)(2.12,2.52)(2.1,3.24)};
    \draw[BookInk!45] plot[smooth] coordinates{(3.05,.48)(3.15,1.59)(3.12,2.68)(3.1,3.36)};
    \draw[BookInk!45] plot[smooth] coordinates{(4.05,.58)(4.15,1.72)(4.12,2.84)(4.1,3.48)};
    \draw[BookInk!45] plot[smooth] coordinates{(.03,1.15)(1.1,1.05)(2.4,1.3)(4.1,1.65)};
    \draw[BookInk!45] plot[smooth] coordinates{(.06,2.1)(1.15,1.98)(2.55,2.25)(4.18,2.58)};
    \node[book label] at (2.1,-.38) {局部面积因子 $|\det D\Phi|$};
  \end{scope}
\end{tikzpicture}
\caption{细网格在每一点近似由线性映射 $D\Phi$ 送成平行多面体；单射保证不同小格的像不重叠。}
\label{fig:3-2-2-curved-grid}
\end{figure}

线性仿射变换的公式立即给出 Gaussian 型密度在缩放下的归一化常数；第~5~章引入随机向量后，
同一计算会成为 Gaussian 随机向量的坐标变换公式。极坐标与径向公式则反复出现在卷积、
Fourier 缩放以及 PDE 基本解的常数计算中。这里证明的是 $C^1$ 微分同胚及其可数分支版本，
因此，映射若出现折叠，应用分片公式前必须先给出满足命题假设的单射分支分解。
\begin{exercise}
从函数换元公式推导一维严格单调 $C^1$ 微分同胚的换元公式，并分别处理递增与递减情形。
\end{exercise}
\begin{exercise}
设 $A$ 可逆、$b\in\R^n$，$X$ 有密度 $p$。计算 $AX+b$ 的密度，并核对其积分为一。
\end{exercise}
\begin{exercise}
由式 \eqref{eq:radial-integration} 计算
$\int_{\R^n}e^{-|x|^2}\,\dd x$，并与一维高斯积分的乘积比较，求出 $\omega_n$。
\end{exercise}
\begin{exercise}
对例~\ref{ex:3-2-2-4} 把 $t$ 分成区间 $(2\pi k,2\pi(k+1))$。写出
\eqref{eq:piecewise-change-variables} 的重数函数，并说明为什么常数函数 $g=1$ 使两侧同为无穷。
\end{exercise}

换元公式处理的是确定的坐标搬运。概率转移和积分算子更进一步：一个起点不再去往一个点，而是产生一整个
目标测度。要让这些测度能够继续积分，参数可测性必须成为定义的一部分。

\subsection{测度核、参数积分与积分号下微分}\label{subsec:3-2-3}
\index{测度核}\index{Markov 核}\index{积分号下微分}

先看一个没有可测性困难的模型。设 $X=\{1,\dots,m\}$、$Y=\{1,\dots,n\}$，并给定非负矩阵
$(k_{ij})$。固定 $i$ 后，第 $i$ 行定义 $Y$ 上的测度
$K(i,\{j\})=k_{ij}$；若每行之和为一，它就是转移概率矩阵。对列向量 $f$，
\[
 Kf(i)=\sum_{j=1}^nk_{ij}f(j).
\]
这里“一个状态 $i$ 被送到一行权重”比确定性映射“一个点被送到一个点”多出了一层平均。无限空间中，
要使这层平均随 $i$ 合法变化，逐点为测度还不够。

\begin{definition}[测度核]\label{def:3-2-3-1}
设 $(X,\mathcal A)$、$(Y,\mathcal B)$ 为可测空间。从 $X$ 到 $Y$ 的\emph{测度核}
$K:X\leadsto Y$ 是一个映射 $K:X\times\mathcal B\to[0,\infty]$，满足：
\begin{enumerate}
  \item 对每个 $x\in X$，$B\mapsto K(x,B)$ 是 $Y$ 上的测度；
  \item 对每个 $B\in\mathcal B$，$x\mapsto K(x,B)$ 是 $\mathcal A$-可测函数。
\end{enumerate}
若 $K(x,Y)<\infty$ 对每个 $x$ 成立，称 $K$ 为有限核；这里不要求这些总质量关于 $x$ 一致有界。
若 $K(x,Y)=1$ 对每个 $x$ 成立，称 $K$ 为概率核或 Markov 核。
\end{definition}

第二项不是第一项的推论。例如，在任意不可测函数 $a:X\to\{0,1\}$ 上令
$K(x,\cdot)=\delta_{a(x)}$，每个 $K(x,\cdot)$ 都是概率测度，但
$x\mapsto K(x,\{1\})=a(x)$ 不可测，因而它不是核。

\begin{definition}[核作用]\label{def:3-2-3-2}
若 $K:X\leadsto Y$ 为核、$f:Y\to[0,\infty]$ 可测，定义
\[
 Kf(x)=\int_Yf(y)K(x,\dd y).
\]
若 $\mu$ 是 $X$ 上的测度，定义 $Y$ 上的集合函数
\[
 (\mu K)(B)=\int_XK(x,B)\,\mu(\dd x),\qquad B\in\mathcal B.
\]
\end{definition}

后一个对象确实是测度。空集的积分为零；若 $B_j$ 两两不交，则对每个 $x$，
\[
 K\!\left(x,\bigcup_jB_j\right)
 =\lim_{N\to\infty}\sum_{j=1}^NK(x,B_j).
\]
对部分和应用单调收敛定理，得到
$(\mu K)(\bigcup_jB_j)=\sum_j(\mu K)(B_j)$。有限状态情形中，$\mu K$ 就是行向量乘矩阵。

\begin{definition}[核复合]\label{def:3-2-3-3}
若 $K:X\leadsto Y$、$L:Y\leadsto Z$，定义
\[
 (KL)(x,C)=\int_YL(y,C)K(x,\dd y),\qquad C\in\mathcal C.
\]
\end{definition}

这个定义是否真的给出核，首先取决于核积分保持可测性。

\begin{theorem}[核积分可测性]\label{theorem:3-2-3-1}
若 $K:X\leadsto Y$ 为测度核，$f:Y\to[0,\infty]$ 可测，则 $Kf$ 可测。
\end{theorem}

\begin{proof}
若 $f=\1_B$，则 $Kf(x)=K(x,B)$，由核定义可测。若
$f=\sum_{j=1}^ra_j\1_{B_j}$ 是非负简单函数，积分的有限线性性给出
$Kf=\sum_ja_jK\1_{B_j}$，故仍可测。对一般 $f\ge0$，取非负简单函数
$s_n\uparrow f$。固定 $x$ 后，对测度 $K(x,\cdot)$ 应用单调收敛定理，得到
\[
 Ks_n(x)\uparrow Kf(x).
\]
右侧是可测函数列的点态极限，因此可测。
\end{proof}

证明特意只在非负锥上进行。若核不是有限核，任意有界带符号 $f$ 仍可能使正负两部分积分都为
$+\infty$，此时 $Kf$ 根本没有定义，不能无条件套用带符号函数的单调类论证。
对实值或复值可测函数 $f$，令
$D_f=\{x:K|f|(x)=\int_Y|f(y)|K(x,\dd y)<\infty\}$。定理~\ref{theorem:3-2-3-1}
说明 $D_f$ 可测；在 $D_f$ 上按第~3.1 节的实部、虚部定义 $Kf(x)$。同一定理分别作用于正负部，
说明 $Kf$ 在 $D_f$ 上可测；在 $X\setminus D_f$ 上置零便得到一个全局可测代表。特别地，若
$K$ 是有限核且 $f$ 有界，则该条件对每个 $x$ 成立，因而 $Kf$ 是处处有定义的可测函数。

\begin{example}[确定性核]\label{ex:3-2-3-1}
可测映射 $T:X\to Y$ 定义
\[
 K_T(x,B)=\delta_{T(x)}(B)=\1_B(T(x)).
\]
它是概率核，并且
$K_Tf=f\circ T$。对 $X$ 上测度 $\mu$，
\[
 (\mu K_T)(B)=\int_X\1_B(T(x))\,\dd\mu(x)
 =\mu(T^{-1}B)=T_*\mu(B).
\]
因此推前测度正是确定性核作用的特例。
\end{example}

\begin{example}[密度核]\label{ex:3-2-3-2}
设 $(Y,\mathcal B,\nu)$ 为 $\sigma$-有限测度空间，
$k:X\times Y\to[0,\infty]$ 联合可测，并且
\[
 \int_Yk(x,y)\,\dd\nu(y)=1\qquad(x\in X).
\]
令
$K(x,B)=\int_Bk(x,y)\,\dd\nu(y)$。固定 $x$ 时这是概率测度；固定 $B$ 时，
引理~\ref{proposition:3-2-1-3} 应用于 $k\1_{X\times B}$，说明
$x\mapsto K(x,B)$ 可测。因此 $K$ 是概率核。密度记号本身不会自动提供这项参数可测性；联合可测与
$\sigma$-有限性正是在这里使用的。
\end{example}

\begin{proposition}[核复合与结合律]\label{proposition:3-2-3-2}
若 $K:X\leadsto Y$、$L:Y\leadsto Z$ 为测度核，则 $KL$ 仍为测度核。对任意第三个核
$M:Z\leadsto W$，
\[
 (KL)M=K(LM).
 \tag{3.2.18}\label{eq:kernel-associativity}
\]
等式允许共同值为 $+\infty$。若各核均为概率核，则复合核仍为概率核。
\end{proposition}

\begin{proof}
固定 $x$ 后，$C\mapsto(KL)(x,C)$ 是测度：证明与定义~\ref{def:3-2-3-2} 后验证
$\mu K$ 为测度的论证相同，只需把 $K(x,\cdot)$ 当作 $Y$ 上的固定测度。固定 $C$ 后，
$y\mapsto L(y,C)$ 非负可测；定理~\ref{theorem:3-2-3-1} 说明
$x\mapsto K(L\1_C)(x)=(KL)(x,C)$ 可测。因此 $KL$ 是核。

先证明一个函数恒等式。若 $f=\1_C$，核复合的定义给出
\[
 (KL)f(x)=(KL)(x,C)=K(Lf)(x).
\]
有限线性组合把等式扩张到非负简单函数。若 $s_n\uparrow f$，分别对
$L(y,\cdot)$、$K(x,\cdot)$ 与 $(KL)(x,\cdot)$ 使用单调收敛定理，得到
\[
 (KL)f=K(Lf)\qquad(f\ge0).
 \tag{3.2.19}\label{eq:kernel-function-assoc}
\]
于是对 $C\subset W$ 可测，取 $f=M\1_C$，便有
\[
 \begin{aligned}
 ((KL)M)(x,C)
 &=(KL)(M\1_C)(x)\\
 &=K(L(M\1_C))(x)
 =(K(LM))(x,C).
 \end{aligned}
\]
这证明结合律。若 $K,L$ 为概率核，则
$(KL)(x,Z)=K(L\1_Z)(x)=K\1_Y(x)=1$。整个证明只有非负积分，没有
$\infty-\infty$，也不需要 $\sigma$-有限性。
\end{proof}

在有限状态空间上，式 \eqref{eq:kernel-associativity} 就是矩阵乘法结合律。第~5~章定义
Markov 链后，会把 $K^n$ 解释为 $n$ 步转移核，并把
$K^{m+n}=K^mK^n$ 称为 Chapman--Kolmogorov 方程。它的解析基础正是这里证明的概率核复合与
结合律。

\begin{example}[热核的平移密度]\label{ex:3-2-3-3}
对 $t>0$ 定义
\[
 p_t(z)=(4\pi t)^{-n/2}\exp\!\left(-\frac{|z|^2}{4t}\right).
\]
由一维高斯积分、Tonelli 以及线性换元 $z=2\sqrt t\,u$，
\[
 \int_{\R^n}e^{-|z|^2/(4t)}\,\dd z=(4\pi t)^{n/2},
\]
故 $p_t$ 非负且积分为一。联合连续函数 $(x,y)\mapsto p_t(y-x)$ 给出密度核
\[
 K_t(x,\dd y)=p_t(y-x)\,\dd y.
\]
它是平移不变 Markov 核。用卷积记号，核作用可以写成热核卷积；上面的归一化计算同时证明了
$K_t(x,\R^n)=1$。
\end{example}

\begin{example}[Green 型有限核]
密度核也包括经典积分方程中的 Green 型核。在区间 $(0,1)$ 上令
\[
 G(x,y)=\min\{x,y\}-xy
 =\begin{cases}
   y(1-x),&0<y\le x,\\
   x(1-y),&x<y<1.
  \end{cases}
\]
函数 $G$ 非负且联合连续，因而
\[
 K_G(x,B)=\int_B G(x,y)\,\dd y
\]
定义一个从 $(0,1)$ 到自身的测度核。它不是概率核，因为直接分段积分得到
\[
\begin{aligned}
 K_G(x,(0,1))
 &=\int_0^x y(1-x)\,\dd y+\int_x^1x(1-y)\,\dd y\\
 &=(1-x)\frac{x^2}{2}+x\frac{(1-x)^2}{2}
 =\frac{x(1-x)}2.
\end{aligned}
\]
相应核作用是 $K_Gf(x)=\int_0^1G(x,y)f(y)\,\dd y$。在常微分方程中，$G$ 是
Dirichlet 边界条件下算子 $-\frac{\dd^2}{\dd x^2}$ 的 Green 函数。这个解释说明了公式的来源，
而上面的分段积分则具体展示了非归一化有限核的总质量。
\end{example}

核定义只保证 $Kf$ 可测。若还希望参数连续，必须把点态连续与一个共同的可积支配放在一起。

\begin{theorem}[参数连续积分]\label{theorem:3-2-3-3}
设拓扑空间 $T$ 在 $t_0$ 处第一可数，$f:T\times X\to\R$ 或 $\C$。假设每个
$f(t,\cdot)$ 可测，并存在 $t_0$ 的邻域 $U$、非负函数 $g\in L^1(\mu)$ 与零集 $N$，使得对
$x\notin N$：
\begin{enumerate}
  \item $t\mapsto f(t,x)$ 在 $t_0$ 连续；
  \item $|f(t,x)|\le g(x)$ 对所有 $t\in U$ 成立。
\end{enumerate}
则 $F(t)=\int_Xf(t,x)\,\dd\mu(x)$ 对 $t\in U$ 有定义，并在 $t_0$ 连续。
\end{theorem}

\begin{proof}
支配条件说明每个 $t\in U$ 的截面 $f(t,\cdot)$ 都绝对可积。若 $(t_n)\subset U$ 且
$t_n\to t_0$，则在共同零集外
$f(t_n,x)\to f(t_0,x)$，且
$|f(t_n,x)-f(t_0,x)|\le2g(x)$。复值支配收敛定理给出
\[
 F(t_n)-F(t_0)=\int_X(f(t_n,x)-f(t_0,x))\,\dd\mu(x)\longrightarrow0.
\]

还需把序列结论升级为拓扑连续性。将 $t_0$ 的一个可数邻域基逐项与支配邻域
$U$ 交，再取前 $n$ 项之交，得到一个满足 $U_n\subset U$ 的递减可数邻域基
$(U_n)$。若 $F$ 不在
$t_0$ 连续，则存在 $\varepsilon>0$，使每个 $U_n$ 中都可取 $t_n$ 满足
$|F(t_n)-F(t_0)|\ge\varepsilon$。递减邻域基保证 $t_n\to t_0$，与上一段矛盾。
第一可数性正用于这一步；对任意拓扑空间中的网，不能把序列版支配收敛定理无声替代为网版本。
\end{proof}

\begin{remark}[Feller 性]
当 $X,Y$ 带拓扑时，可测 Markov 核未必把连续函数送到连续函数。若
$K(C_b(Y))\subset C_b(X)$，称 $K$ 具有 Feller 性；在局部紧理论中也常使用 $C_0$ 版本。
例如，对热核与 $f\in C_b(\R^n)$，若 $x_k\to x$，则
\[
 K_tf(x_k)=\int_{\R^n}f(x_k+z)p_t(z)\,\dd z
 \longrightarrow\int_{\R^n}f(x+z)p_t(z)\,\dd z
\]
由参数连续积分定理得到，因为点态连续成立，且共同支配为
$\|f\|_\infty p_t\in L^1$。因此 $K_tf$ 连续，且
$|K_tf|\le\|f\|_\infty$。
\end{remark}

积分号下微分需要支配的不是单个导数值，而是邻近参数的全部差商。失败模型已经很具体：在
$(0,1)$ 上令 $f(t,x)=\1_{(0,t)}(x)$，则对每个固定 $x>0$，$t=0$ 处的右导数为零；可是差商
$t^{-1}\1_{(0,t)}$ 的积分恒为 $1$。因此下一定理把统一可积支配直接加在差商上；
例~\ref{ex:3-2-3-4} 随后用显式计算对比满足该条件与违反该条件的模型。

\begin{theorem}[积分号下微分]\label{theorem:3-2-3-4}
设 $I\subset\R$ 为开区间，$t_0\in I$，且 $f:I\times X\to\R$ 或 $\C$ 满足每个
$f(t,\cdot)$ 可测。假设 $f(t_0,\cdot)\in L^1(\mu)$，
$\partial_tf(t_0,x)$ 对几乎处处 $x$ 存在，并且存在 $\delta>0$、$g\in L^1(\mu)$ 与一个零集
$N$，使 $[t_0-\delta,t_0+\delta]\subset I$，且
\[
 \left|\frac{f(t_0+h,x)-f(t_0,x)}{h}\right|\le g(x)
 \tag{3.2.20}\label{eq:difference-quotient-domination}
\]
对所有 $x\notin N$ 及 $0<|h|<\delta$ 成立。则 $F(t)=\int_Xf(t,x)\,\dd\mu(x)$
在 $|t-t_0|<\delta$ 时有定义，并且
\[
 F'(t_0)=\int_X\partial_tf(t_0,x)\,\dd\mu(x).
 \tag{3.2.21}\label{eq:differentiate-under-integral}
\]
\end{theorem}

\begin{proof}
由 \eqref{eq:difference-quotient-domination}，
\[
 |f(t_0+h,x)|\le|f(t_0,x)|+|h|g(x)
\]
在同一零集外成立，所以 $f(t_0+h,\cdot)\in L^1$。取任意实数列
$h_k\to0$ 且 $0<|h_k|<\delta$。相应差商是可测函数，并在几乎处处收敛到
$\partial_tf(t_0,\cdot)$。把导数不存在的可测零集上的值定义为零；所得函数在零集补集上等于
可测差商列的极限，因而可测。由支配收敛定理，它属于
$L^1$，且
\[
 \begin{aligned}
 \lim_{k\to\infty}\frac{F(t_0+h_k)-F(t_0)}{h_k}
 &=\lim_{k\to\infty}\int_X
 \frac{f(t_0+h_k,x)-f(t_0,x)}{h_k}\,\dd\mu(x)\\
 &=\int_X\partial_tf(t_0,x)\,\dd\mu(x).
 \end{aligned}
\]
极限对每个趋于零的非零实数列都相同；实直线中的序列判别说明差商本身在 $h\to0$ 时有该极限。
这就是 \eqref{eq:differentiate-under-integral}。
\end{proof}

一个常用的充分条件是：对几乎处处 $x$，函数 $t\mapsto f(t,x)$ 在邻域中处处可微，而且
$|\partial_tf(t,x)|\le g(x)$。实值情形由一维中值定理得到
\eqref{eq:difference-quotient-domination}。复值情形也不需要复分析：若
$z=f(t_0+h,x)-f(t_0,x)\ne0$，取 $u=z/|z|$，对实函数
$t\mapsto\Re(\overline u f(t,x))$ 使用中值定理，得到
\[
 |z|=\Re(\overline uz)
 \le |h|\sup_{t\text{ 在线段上}}|\partial_tf(t,x)|.
\]
若 $z=0$，则左端为零而右端非负，估计仍成立。

\begin{example}[成功与失败的差商]\label{ex:3-2-3-4}
对 $t>0$ 令
\[
 F(t)=\int_0^\infty e^{-tx}\,\dd x.
\]
固定 $t_0>0$。当 $|t-t_0|<t_0/2$ 时，
\[
 |\partial_te^{-tx}|=xe^{-tx}\le xe^{-t_0x/2},
\]
作代换 $u=t_0x/2$，并用分部积分
$\int_0^\infty ue^{-u}\,\dd u=[-ue^{-u}]_0^\infty+\int_0^\infty e^{-u}\,\dd u=1$，可得
$\int_0^\infty xe^{-t_0x/2}\,\dd x=4/t_0^2$；以 $u=t_0x$ 作同一计算又得
$\int_0^\infty xe^{-t_0x}\,\dd x=t_0^{-2}$。因此定理~\ref{theorem:3-2-3-4} 给出
\[
 F'(t_0)=-\int_0^\infty xe^{-t_0x}\,\dd x=-t_0^{-2},
\]
与直接计算 $F(t)=t^{-1}$ 一致。

反过来，在 $x\in(0,1)$、$0<t<1$ 时令
$f(t,x)=\1_{(0,t)}(x)$，并置 $f(0,x)=0$。对每个固定 $x>0$，当 $0<t<x$ 时
$f(t,x)=0$，所以 $f(\cdot,x)$ 在 $0$ 的右导数为零。然而
\[
 \int_0^1f(t,x)\,\dd x=t
\]
在 $0$ 的右导数为一。差商
$t^{-1}\1_{(0,t)}$ 不可能有共同可积支配；否则支配收敛定理会迫使其积分趋于零，而其积分恒为一。
仅有逐点导数存在不足以交换微分与积分。
\end{example}

\begin{figure}[htbp]
\centering
\begin{tikzpicture}[node distance=18mm and 23mm]
  \node[book strong node,minimum width=25mm] (x) {$x\in X$};
  \node[book strong node,right=of x,minimum width=31mm] (y) {$K(x,\cdot)$\\在 $Y$ 上的质量云};
  \node[book strong node,right=of y,minimum width=31mm] (z) {$(KL)(x,\cdot)$\\在 $Z$ 上的质量云};
  \draw[book accent arrow] (x)--node[above]{$K$}(y);
  \draw[book accent arrow] (y)--node[above]{$L$}(z);
  \draw[book dashed arrow,bend right=27] (x.south)--node[below]{$KL$}(z.south);
  \foreach \p in {-.32,0,.32}{
    \fill[BookAccent!70] ($(y.center)+(0.45,\p)$) circle (1.2pt);
    \fill[BookWarm!75] ($(z.center)+(0.58,\p)$) circle (1.2pt);
  }
\end{tikzpicture}
\caption{核把一点散成一个测度；复合核对中间的质量云再作一次积分。}
\label{fig:3-2-3-kernel-cloud}
\end{figure}

\begin{remark}[正则条件分布与核]
在概率论中，正则条件分布是一个 Markov 核。任意两个可测空间之间未必存在这种核；常用的存在定理
要求目标空间是标准 Borel 空间（standard Borel space）等额外条件。这里“标准 Borel”指一个与某个
完备可分度量空间的 Borel 子集（赋予相对 Borel $\sigma$-代数）可测同构的可测空间。
本章不证明正则条件分布的存在定理，也不在后续论证中调用它；使用条件分布时必须另行核对相应的存在性假设。
\end{remark}

在统计中，边缘化就是让一个概率核作用在常数或似然函数上；梯度交换则必须逐项核对
定理~\ref{theorem:3-2-3-4} 的差商支配。积分算子的半群性质可由
命题~\ref{proposition:3-2-3-2} 化为核复合。三类应用共享同一条边界：可测性由核公理负责，
极限或微分的交换另需支配条件。

\begin{exercise}
证明确定性核满足 $K_SK_T=K_{T\circ S}$，其中映射方向取为
$X\xrightarrow{S}Y\xrightarrow{T}Z$；并用推前测度核对次序。
\end{exercise}
\begin{exercise}
设 $K(x,\dd y)=k(x,y)\nu(\dd y)$、
$L(y,\dd z)=\ell(y,z)\rho(\dd z)$ 是概率密度核。用 Tonelli 写出 $KL$ 相对于 $\rho$ 的密度，
并直接核对其积分为一。
\end{exercise}
\begin{exercise}
从定义证明三个概率核的结合律，并在有限状态情形把每一步写成矩阵求和。
\end{exercise}
\begin{exercise}
在例~\ref{ex:3-2-3-4} 的失败模型中证明：不存在 $g\in L^1(0,1)$ 同时支配所有
$t^{-1}\1_{(0,t)}$，$0<t<1/2$。
\end{exercise}

核之间的绝对连续关系仍缺少一个关键问题：如果每个 $K(x,\cdot)$ 都相对于同一测度绝对连续，
Radon--Nikodym 定理会对每个固定的 $x$ 给出密度的几乎处处类，却不会自动从这些等价类中选出
联合可测的 $(x,y)\mapsto k(x,y)$。联合可测密度的存在需要另外的可测 Radon--Nikodym
选择定理及相应的可数生成或标准 Borel 结构。下一节先建立每个标量测度的 RN 表示与唯一性，
并由此进入支撑后续条件期望构造的 $L^p$ 空间；条件期望的正式算子理论留在第~5~章。

\section{密度、\texorpdfstring{$L^p$}{Lp} 空间、Hahn--Banach 与对偶}\label{sec:03-03}

\subsection{Radon--Nikodym 与 Lebesgue 分解}\label{subsec:3-3-1}
\index{Radon--Nikodym 定理}\index{Lebesgue 分解}

在区间上，密度 $h$ 产生测度
\[
  \nu(E)=\int_E h\,\dd\mu .
\]
反过来，若我们只知道集合函数 $\nu$，能否找回这样的 $h$？一个必要条件立刻从公式中读出：
$\mu(E)=0$ 时必须有 $\nu(E)=0$。这个条件排除了相对于 Lebesgue 测度的 Dirac 质量，却没有
说明它是否已经充分。更麻烦的是，即便密度存在，我们也不应期待逐点唯一；改变零集上的值不会改变
任何积分。

\begin{definition}[绝对连续与奇异]\label{def:3-3-1-1}
设 $\mu,\nu$ 是同一可测空间 $(X,\mathcal A)$ 上的正测度。若
\[
  \mu(E)=0\quad\Longrightarrow\quad \nu(E)=0,
  \qquad E\in\mathcal A,
\]
则称 $\nu$ 关于 $\mu$ 绝对连续，记作 $\nu\ll\mu$。若存在 $N\in\mathcal A$ 使
$\mu(N)=0$ 且 $\nu(X\setminus N)=0$，则称 $\nu$ 与 $\mu$ 互相奇异，记作 $\nu\perp\mu$。
\end{definition}

\begin{definition}[Radon--Nikodym 导数]\label{def:3-3-1-2}
若 $\nu\ll\mu$ 且存在可测函数 $h:X\to[0,\infty]$，使
\[
  \nu(E)=\int_Eh\,\dd\mu\qquad(E\in\mathcal A),
\]
则把 $h$ 称为 $\nu$ 关于 $\mu$ 的 Radon--Nikodym 导数，记作
$h=\dd\nu/\dd\mu$。这个记号总按 $\mu$-a.e. 等价理解。
\end{definition}

\begin{definition}[Lebesgue 分解]\label{def:3-3-1-3}
若
\[
  \nu=\nu_a+\nu_s,\qquad \nu_a\ll\mu,\qquad \nu_s\perp\mu,
\]
则称右式为 $\nu$ 关于 $\mu$ 的 Lebesgue 分解。
\end{definition}

\begin{example}[密度与原子]\label{ex:3-3-1-1}
在 $\R$ 上记 Lebesgue 测度为 $m$。对 $0<p<1$，令
\[
  \nu=p\,m|_{[0,1]}+(1-p)\delta_0.
\]
第一项的密度是 $p\1_{[0,1]}$，所以它关于 $m$ 绝对连续；第二项集中在 $m$-零集
$\{0\}$ 上，所以它与 $m$ 奇异。特别地，$\nu$ 不可能整个写成 $h\,m$：否则
$1-p=\nu(\{0\})=\int_{\{0\}}h\,\dd m=0$。这正是 Lebesgue 分解中两类质量同时出现的
最小模型。
\end{example}


奇异质量未必由原子组成。下面的例子把这两个概念彻底分开。

\begin{example}[Cantor 测度]\label{ex:3-3-1-2}
令 $C:\R\to[0,1]$ 为标准 Cantor 函数在 $[0,1]$ 外作常值延拓，并令 $\gamma$ 为第~2.3 节
Stieltjes 构造给出的测度，即
\[
  \gamma((s,t])=C(t)-C(s).
\]
$C$ 从 $0$ 增长到 $1$，所以 $\gamma(\R)=1$。Cantor 集 $K$ 的补集是可数个两两不交的三分
开区间 $(a_j,b_j)$ 之并，且 $C$ 在每个 $[a_j,b_j]$ 上恒定。对固定 $j$，由测度从下连续与
Stieltjes 增量公式，
\[
 \gamma((a_j,b_j))
 =\lim_{n\to\infty}\gamma((a_j,b_j-1/n])
 =\lim_{n\to\infty}\bigl(C(b_j-1/n)-C(a_j)\bigr)=0,
\]
其中只取使 $b_j-1/n>a_j$ 的 $n$。因此可数可加性给出
$\gamma(\R\setminus K)=\sum_j\gamma((a_j,b_j))=0$，而第~2~章的长度计算给出
$m(K)=0$，所以 $\gamma\perp m$。另一方面，Stieltjes 测度在 $x$ 处的原子质量是跳跃
\[
  \gamma(\{x\})=C(x)-C(x-)=0,
\]
因为 $C$ 连续。因此 $\gamma$ 是无原子的奇异概率；“奇异”与“离散”并不是同义词。
\end{example}

Radon--Nikodym 定理的困难不在于猜一个公式，而在于证明候选密度已经吃掉全部质量。下面的证明
只使用非负积分的单调收敛定理以及有符号测度的 Hahn 分解。

历史上，Radon 与 Nikodym 的结果把一维微积分中“增量来自导数”的现象推广到了抽象测度。
现代教材也常借助 Hilbert 空间的表示定理证明有限情形；那条路线需要先建立更多泛函分析。
这里采用纯测度的最大密度法：先解决有限质量问题，再用可数分块补上无限总质量的缺口。

\begin{theorem}[Radon--Nikodym 定理]\label{theorem:3-3-1-1}
设 $\mu,\nu$ 是 $(X,\mathcal A)$ 上的 $\sigma$-有限正测度，且 $\nu\ll\mu$。则存在可测
$h\ge0$ 使 $\nu=h\mu$，而且 $h$ 在 $\mu$-a.e. 意义下唯一。若 $\nu(X)<\infty$，则
$h\in L^1(\mu)$ 且
\[
  \int_Xh\,\dd\mu=\nu(X).
\]
\end{theorem}

\begin{proof}
先设 $\mu(X),\nu(X)<\infty$。考虑
\[
 \mathcal H=\left\{g:X\to[0,\infty]\text{ 可测}:
     \int_Eg\,\dd\mu\le\nu(E)\text{ 对每个 }E\in\mathcal A\right\},
 \qquad
 \alpha=\sup_{g\in\mathcal H}\int_Xg\,\dd\mu.
\]
$0\in\mathcal H$，且每个 $g\in\mathcal H$ 满足 $\int g\le\nu(X)$，故
$0\le\alpha\le\nu(X)<\infty$。若 $g_1,g_2\in\mathcal H$，令
$A=\{g_1\ge g_2\}$。对任意 $E\in\mathcal A$，
\[
 \int_E(g_1\vee g_2)\,\dd\mu
 =\int_{E\cap A}g_1\,\dd\mu+\int_{E\setminus A}g_2\,\dd\mu
 \le\nu(E\cap A)+\nu(E\setminus A)=\nu(E).
\]
所以 $g_1\vee g_2\in\mathcal H$。

对每个 $n$ 选取 $g_n\in\mathcal H$，使
$\int g_n>\alpha-1/n$；这里对可数个非空候选集作了逐项选择。令
$u_n=g_1\vee\cdots\vee g_n$，则 $u_n\in\mathcal H$ 且 $u_n\uparrow h$，其中 $h$ 可测。
对任意 $E$，单调收敛定理给出
\[
 \int_Eh\,\dd\mu=\lim_n\int_Eu_n\,\dd\mu\le\nu(E),
\]
故 $h\in\mathcal H$，从而 $\int h\le\alpha$。另一方面，$u_n\ge g_n$，所以
$\int h\ge\int u_n\ge\int g_n>\alpha-1/n$ 对每个 $n$ 成立。令 $n\to\infty$，得到
$\int h\ge\alpha$，故 $\int h=\alpha$。

定义集合函数
\[
  \rho(E)=\nu(E)-\int_Eh\,\dd\mu.
\]
因为 $h\in\mathcal H$，右侧对每个 $E$ 都是两个有限数之差且非负。若 $(E_j)$ 两两不交，对每个
$N$ 有有限可加等式
\[
 \rho\!\left(\bigcup_{j=1}^NE_j\right)=\sum_{j=1}^N\rho(E_j).
\]
又因为 $\nu$ 与 $h\mu$ 都是总质量有限的测度，从下连续性给出
\[
 \begin{aligned}
 \rho\!\left(\bigcup_{j=1}^{\infty}E_j\right)
 &=\lim_{N\to\infty}
   \left[\nu\!\left(\bigcup_{j=1}^NE_j\right)
       -\int_{\cup_{j=1}^NE_j}h\,\dd\mu\right]\\
 &=\lim_{N\to\infty}\sum_{j=1}^N\rho(E_j)
 =\sum_{j=1}^{\infty}\rho(E_j).
 \end{aligned}
\]
因此 $\rho$ 是有限正测度。若 $\mu(E)=0$，则 $\nu(E)=0$ 且
$\int_Eh\,\dd\mu=0$，所以 $\rho(E)=0$；因此 $\rho\ll\mu$。我们证明 $\rho=0$。若不然，则
$\rho(X)>0$。若 $\mu(X)=0$，则 $\nu\ll\mu$ 给出 $\nu=0$，进而 $\rho=0$，与
$\rho(X)>0$ 矛盾。因此 $\mu(X)>0$。取
$\varepsilon=\rho(X)/(2\mu(X))>0$，则
\[
  (\rho-\varepsilon\mu)(X)>0;
\]
对有限有符号测度 $\sigma=\rho-\varepsilon\mu$ 应用 Hahn 分解，取正集 $A$ 与负集
$X\setminus A$。由于
$\sigma(X\setminus A)\le0$ 而 $\sigma(X)>0$，有
\[
 (\rho-\varepsilon\mu)(A)>0.
\]
特别地 $\rho(A)>0$；由 $\rho\ll\mu$ 可知 $\mu(A)>0$。正集的定义还给出：对每个可测
$B\subset A$，
\[
  \rho(B)\ge\varepsilon\mu(B).
\]
于是对任意 $E\in\mathcal A$，利用 $\rho$ 的正性，
\[
 \rho(E)\ge\rho(E\cap A)\ge\varepsilon\mu(E\cap A).
\]
因此
\[
 \int_E(h+\varepsilon\1_A)\,\dd\mu
 \le\int_Eh\,\dd\mu+\rho(E)=\nu(E),
\]
即 $h+\varepsilon\1_A\in\mathcal H$。可是
$\int(h+\varepsilon\1_A)=\alpha+\varepsilon\mu(A)>\alpha$，与 $\alpha$ 的定义矛盾。
故 $\rho=0$，有限情形得证。

现在设 $\mu,\nu$ 仅为 $\sigma$-有限。分别取可测覆盖 $(A_r)$、$(B_s)$，使
$\mu(A_r)<\infty$、$\nu(B_s)<\infty$。把所有交集 $A_r\cap B_s$ 排成一列 $(C_n)$，再令
\[
 D_1=C_1,
 \qquad D_n=C_n\setminus\bigcup_{k<n}C_k\quad(n\ge2).
\]
删去空块后，$(D_n)$ 两两不交、并为 $X$，而且每块同时具有有限的 $\mu$-质量和 $\nu$-质量。
把两测度限制到 $D_n$；它们都是有限测度，而且
$\nu|_{D_n}\ll\mu|_{D_n}$。有限情形给出可测密度 $h_n$，在 $D_n$ 外把它置零。令
$h=\sum_nh_n\1_{D_n}$。这是可测函数。对任意 $E$，可数可加性和非负积分对不交分块的可加性给出
\[
 \nu(E)=\sum_n\nu(E\cap D_n)
 =\sum_n\int_{E\cap D_n}h_n\,\dd\mu
 =\int_Eh\,\dd\mu.
\]

最后证明唯一性。设 $h,k$ 都是密度。在每个有限 $\nu$-质量分块 $D_j$ 上，两者的积分有限，
故它们取无穷值的集合都是 $\mu$-零集；先删去这些集合的可数并，再令
\[
 E_{j,m}=D_j\cap\{h\ge k+1/m\}.
\]
在 $E_{j,m}$ 上两边积分都等于有限数 $\nu(E_{j,m})$，所以
\[
 0=\int_{E_{j,m}}(h-k)\,\dd\mu\ge m^{-1}\mu(E_{j,m}).
\]
故每个 $E_{j,m}$ 都是 $\mu$-零集。对 $j,m$ 取可数并得到
$\mu(\{h>k\})=0$；交换 $h,k$ 即得 $h=k$ a.e.。若 $\nu$ 有限，在表示式中取 $E=X$
便得到最后一项结论。
\end{proof}

任意两个 $\sigma$-有限测度的和仍然 $\sigma$-有限：分别取有限质量覆盖，交叉后即可得到对两者
都有限的可数覆盖。这个小观察使 RN 定理可以直接产生 Lebesgue 分解。

\begin{theorem}[Lebesgue 分解定理]\label{theorem:3-3-1-2}
若 $\mu,\nu$ 是 $(X,\mathcal A)$ 上的 $\sigma$-有限正测度，则存在唯一的正测度
$\nu_a,\nu_s$ 使
\[
  \nu=\nu_a+\nu_s,
  \qquad \nu_a\ll\mu,
  \qquad \nu_s\perp\mu.
\]
而且 $\nu_a=h\mu$，其中 $h\ge0$。
\end{theorem}

\begin{proof}
令 $\tau=\mu+\nu$。由前述交叉分块，$\tau$ 是 $\sigma$-有限测度；又有
$\mu\ll\tau$ 与 $\nu\ll\tau$。由 Radon--Nikodym 定理，存在 $a,b\ge0$ 使
\[
  \mu=a\tau,
  \qquad \nu=b\tau.
\]
两式相加给出，对每个 $E\in\mathcal A$，
\[
 \int_E(a+b)\,\dd\tau=\mu(E)+\nu(E)=\tau(E)=\int_E1\,\dd\tau.
\]
用 RN 导数唯一性，得到 $a+b=1$ $\tau$-a.e.。在这个共同零集上把 $a,b$ 都改为
零；此后 $a,b$ 处处有限，而且在零集外属于 $[0,1]$。
令 $N=\{a=0\}$。因为 $\mu(N)=\int_Na\,\dd\tau=0$，可定义
\[
  h(x)=
  \begin{cases}
    b(x)/a(x),&x\notin N,\\
    0,&x\in N,
  \end{cases}
  \qquad
  \nu_a(E)=\int_Eh\,\dd\mu,
  \qquad
  \nu_s(E)=\int_{E\cap N}b\,\dd\tau.
\]
在 $X\setminus N$ 上有 $ha=b$，所以
\[
 \nu_a(E)=\int_{E\setminus N}\frac ba\,a\,\dd\tau
 =\int_{E\setminus N}b\,\dd\tau.
\]
从而 $\nu=\nu_a+\nu_s$。按定义 $\nu_a\ll\mu$，而 $\nu_s$ 集中在 $\mu$-零集 $N$ 上，
故存在性成立。

设另有 $\nu=h_1\mu+\sigma_1=h_2\mu+\sigma_2$，其中 $\sigma_i\perp\mu$。取
$\mu$-零集 $N_i$ 使 $\sigma_i(X\setminus N_i)=0$，并令 $N=N_1\cup N_2$。在 $X\setminus N$
上，两条等式给出 $h_1\mu=h_2\mu$；在 $N$ 上两个绝对连续测度都为零。因此
$h_1\mu=h_2\mu$ 在整个 $X$ 上成立，继而由总测度相等得到 $\sigma_1=\sigma_2$。
RN 导数的唯一性还给出 $h_1=h_2$ $\mu$-a.e.。
\end{proof}

对复测度，密度还携带一个相位。下一个定理中的模长结论不是记号约定，而是全变差定义的实质后果。

\begin{theorem}[复测度极分解]\label{theorem:3-3-1-3}
若 $\nu$ 是有限复测度，则存在 $|\nu|$-可测函数 $u$，使
\[
  |u|=1\quad |\nu|\text{-a.e.},
  \qquad \dd\nu=u\,\dd|\nu|.
\]
因此，每个关于 $\nu$ 可积的函数 $f$ 都满足
\[
  \int f\,\dd\nu=\int fu\,\dd|\nu|.
\]
\end{theorem}

\begin{proof}
记 $\lambda=|\nu|$。若 $\lambda(E)=0$，则对每个 $A\subset E$ 有
$|\nu(A)|\le\lambda(A)=0$。因此 $\Re\nu$、$\Im\nu$ 的 Jordan 正负部分都关于 $\lambda$
绝对连续。分别应用 Radon--Nikodym 定理并合并实部、虚部，可得某个
$w\in L^1(\lambda;\C)$，使
\[
  \nu(E)=\int_Ew\,\dd\lambda.
\]
在一个 $\lambda$-零集上把 $w$ 改为零后，可将它视为处处有限的复值可测函数。

我们从全变差的分割定义证明
\begin{equation}\label{eq:3-3-1-variation-density}
  |\nu|(E)=\int_E|w|\,\dd\lambda
  \qquad(E\in\mathcal A).
\end{equation}
对 $E$ 的任意有限可测分割 $(E_j)$，
\[
 \sum_j|\nu(E_j)|
 \le\sum_j\int_{E_j}|w|\,\dd\lambda
 =\int_E|w|\,\dd\lambda,
\]
故左边取上确界后得到“$\le$”。为证反向，令
$v=w/|w|$ 于 $\{w\ne0\}$，并在 $\{w=0\}$ 上令 $v=1$。给定
$0<\delta<\pi/2$，把单位圆分成有限个 Borel 弧，使每段的角直径不超过 $\delta$；取每段中点
$\zeta_j$，并令 $E_j=E\cap\{v\text{ 落在第 }j\text{ 段}\}$。于是
\[
 \begin{aligned}
 |\nu|(E)
 &\ge\sum_j\left|\int_{E_j}w\,\dd\lambda\right|
 \ge\sum_j\Re\left(\overline{\zeta_j}\int_{E_j}w\,\dd\lambda\right)\\
 &\ge \cos\delta\sum_j\int_{E_j}|w|\,\dd\lambda
 =\cos\delta\int_E|w|\,\dd\lambda.
 \end{aligned}
\]
令 $\delta\downarrow0$，即得 \eqref{eq:3-3-1-variation-density}。

可是 $\lambda=|\nu|$，故 \eqref{eq:3-3-1-variation-density} 对每个 $E$ 给出
$\int_E|w|\,\dd\lambda=\int_E1\,\dd\lambda$。取
$E_n^+=\{|w|>1+1/n\}$，则
\[
  0=\int_{E_n^+}(|w|-1)\,\dd\lambda\ge n^{-1}\lambda(E_n^+),
\]
故 $\lambda(E_n^+)=0$。同样，对 $E_n^-=\{|w|<1-1/n\}$ 有
\[
  0=\int_{E_n^-}(1-|w|)\,\dd\lambda\ge n^{-1}\lambda(E_n^-),
\]
故 $\lambda(E_n^-)=0$。再对 $n$ 取可数并，得到
$|w|=1$ $\lambda$-a.e.。令 $u=w$ 即得极分解。最后的积分公式先对简单函数由定义成立，
再由第~3.1 节的复积分逼近推广到所有可积 $f$。
\end{proof}

密度可以逐层相乘，但“约掉密度”需要两个测度拥有相同零集。

\begin{proposition}[RN 链式法则与倒数]\label{proposition:3-3-1-4}
设 $\lambda,\nu,\mu$ 均为 $\sigma$-有限正测度。若
$\lambda\ll\nu\ll\mu$，则
\[
 \frac{\dd\lambda}{\dd\mu}
 =\frac{\dd\lambda}{\dd\nu}\frac{\dd\nu}{\dd\mu}
 \quad\mu\text{-a.e.},
\]
其中乘积在 $\dd\nu/\dd\mu=0$ 处取零。若还满足 $\mu\ll\nu$，则
\[
  \frac{\dd\mu}{\dd\nu}
  =\left(\frac{\dd\nu}{\dd\mu}\right)^{-1}
  \quad\nu\text{-a.e.}
\]
\end{proposition}

\begin{proof}
记 $a=\dd\lambda/\dd\nu$、$b=\dd\nu/\dd\mu$。先固定代表。由 $\nu$ 的 $\sigma$-有限性，可用
$\nu(E_j)<\infty$ 的可数集合覆盖 $X$；在每个 $E_j$ 上
$\int_{E_j}b\,\dd\mu<\infty$，故 $b<\infty$ $\mu$-a.e.。在所得共同 $\mu$-零集上把 $b$ 置零。
再用 $\lambda(F_j)<\infty$ 的可数集合覆盖 $X$；等式
$\int_{F_j}a\,\dd\nu=\lambda(F_j)$ 说明 $a<\infty$ $\nu$-a.e.。在所得共同 $\nu$-零集上把
$a$ 置零。
这些修改不改变相应 RN 导数。于是下列乘积都是普通的非负实数乘积。

对非负简单函数 $s$，密度定义直接给出
$\int s\,\dd\nu=\int sb\,\dd\mu$；用单调简单逼近可将此式推广到所有非负可测 $s$。
因此对任意 $E$，
\[
 \lambda(E)=\int_Ea\,\dd\nu=\int_Eab\,\dd\mu.
\]
RN 导数的唯一性给出第一式。

再设 $\mu\ll\nu$。集合 $Z=\{b=0\}$ 满足 $\nu(Z)=0$，故 $\mu(Z)=0$，于是
$b>0$ $\mu$-a.e.。把 $1/b$ 在 $Z$ 上定义为零。对任意 $E$，再次使用非负函数的密度换算，
\[
 \int_E\frac1b\,\dd\nu
 =\int_{E\cap\{b>0\}}\frac1b b\,\dd\mu
 =\mu(E).
\]
这证明倒数公式。反之，$b>0$ $\mu$-a.e. 时，$\nu(E)=0$ 蕴含
$\int_Eb\,\dd\mu=0$，从而 $\mu(E)=0$；所以该正性恰好对应 $\mu$ 与 $\nu$ 等价。
\end{proof}

\begin{example}[链式密度]\label{ex:3-3-1-3}
在 $\R$ 上令
\[
 \dd\nu=2x\1_{(0,1)}(x)\,\dd x,
 \qquad \dd\lambda=6x^2\1_{(0,1)}(x)\,\dd x.
\]
则 $\dd\nu/\dd x=2x\1_{(0,1)}(x)$，而
$\dd\lambda/\dd\nu=3x$ 于 $(0,1)$；在补集上取零。相乘得到
$6x^2\1_{(0,1)}(x)$，即
\[
  \frac{\dd\lambda}{\dd x}
  =\frac{\dd\lambda}{\dd\nu}\frac{\dd\nu}{\dd x}.
\]
补集上的比值只在 $\nu$-零集上被指定，因此不影响链式法则。另一方面，Lebesgue 测度并不关于
$\nu$ 绝对连续，例如 $(2,3)$ 有正 Lebesgue 测度而 $\nu((2,3))=0$；所以不能在全空间写
$\dd x/\dd\nu=(2x\1_{(0,1)})^{-1}$。这说明倒数公式中的测度等价性不能删去。
\end{example}

\begin{figure}[htbp]
  \centering
  \begin{tikzpicture}[node distance=9mm and 13mm]
    \node[book strong node] (nu) {$\nu$\quad 总质量};
    \node[book node,below left=of nu] (ac) {$h\mu$\quad 绝对连续部分};
    \node[book warning node,below right=of nu] (sing) {$\nu_s$\quad 奇异部分};
    \node[book warm node,below=of ac] (base) {$\mu$\quad 参考测度};
    \node[book node,below=of sing] (null) {$N$\quad $\mu(N)=0$};
    \draw[book accent arrow] (nu)--node[left,book label] {密度} (ac);
    \draw[book arrow] (nu)--node[right,book label] {剩余} (sing);
    \draw[book accent arrow] (base)--node[right,book label] {$\times h$} (ac);
    \draw[book arrow] (sing)--node[right,book label] {集中} (null);
  \end{tikzpicture}
  \caption{Lebesgue 分解把能由参考测度看见的质量写成密度；其余质量只能集中在参考测度的零集上。}
  \label{fig:3-3-1-lebesgue-decomposition}
\end{figure}

RN 导数统一了几个看似不同的公式。若 $Q\ll P$ 是两个总质量为一的测度，令
$L=\dd Q/\dd P$。只要 $Q(B)>0$，限制测度归一化便给出
\[
  \frac{Q(A\cap B)}{Q(B)}
  =\frac{\int_{A\cap B}L\,\dd P}{\int_BL\,\dd P}.
\]
这只是密度公式的代入，分母非零的条件不能省略。在概率记号中，左端写作 $Q(A\mid B)$，
$L$ 称为似然比。

再设 $(\Omega,\mathcal F,\Prob)$ 是总质量为一的测度空间，$X\in L^1(\Prob)$，且
$\mathcal G\subset\mathcal F$ 是子 $\sigma$-代数。实值情形中，对
$A\in\mathcal G$ 定义有限有符号测度
\[
  \eta(A)=\int_AX\,\dd\Prob.
\]
对 $A\in\mathcal G$ 令
$\eta_\pm(A)=\int_A X^\pm\,\dd\Prob$。这是有限正测度，且
$\eta_\pm\ll\Prob|_{\mathcal G}$。分别应用 RN 定理，得到 $\mathcal G$-可测函数
$Y_\pm\ge0$，满足 $\eta_\pm=Y_\pm\Prob|_{\mathcal G}$。由于
$\int Y_\pm\,\dd\Prob=\eta_\pm(\Omega)=\int X^\pm\,\dd\Prob<\infty$，函数
$Y=Y_+-Y_-$ 可积，并且
\[
  \int_AY\,\dd\Prob=\int_AX\,\dd\Prob
  \qquad(A\in\mathcal G).
\]
若另一个 $\mathcal G$-可测可积函数 $Z$ 也满足同一等式，对
$A_n^+=\{Y-Z>1/n\}$ 取积分得
$0=\int_{A_n^+}(Y-Z)\,\dd\Prob\ge n^{-1}\Prob(A_n^+)$。对 $n$ 求并并对
$Z-Y$ 重复这一论证，得 $Y=Z$ a.e.。因此 $Y$ 在 a.e. 意义下唯一，记作
$Y_{X,\mathcal G}$。条件期望理论把这一函数记为
$\E[X\mid\mathcal G]$，并进一步研究其正性和 $L^1$ 压缩性。复值情形可分别处理实部与虚部。
若还要把条件积分表示为给定随机变量取值后的概率核，则需要标准 Borel 空间等额外结构；标量
RN 定理本身不保证这种核表示存在。

一维绝对连续函数也落在同一框架。本段采用如下自足定义：
若对每个 $\varepsilon>0$ 都存在 $\delta>0$，使任意有限族两两不交开区间
$(a_k,b_k)\subset[a,b]$ 在 $\sum_k(b_k-a_k)<\delta$ 时满足
$\sum_k|F(b_k)-F(a_k)|<\varepsilon$，就称 $F:[a,b]\to\R$ 绝对连续；并定义
\[
 \operatorname{Var}(F;[a,b])
 :=\sup_{a=t_0<\cdots<t_N=b}\sum_{j=1}^N|F(t_j)-F(t_{j-1})|.
\]
该上确界有限时称 $F$ 有界变差。第~4~章将从微分的角度进一步研究这两个概念。

若 $F:[a,b]\to\R$ 绝对连续，则 $F$ 有界变差。为核实这一点，
对 $\varepsilon=1$ 取绝对连续性给出的 $\delta>0$，再把 $[a,b]$ 分成有限个长度小于
$\delta$ 的区间，设共有 $R$ 段。对任意分割 $a=t_0<\cdots<t_N=b$ 加入这些固定端点。
由三角不等式，加入端点只会使增量绝对值之和不减，所以只需估计加细后的分割。
若固定小区间为 $[c_{r-1},c_r]$，则落在其中的分割子区间两两不交，长度和不超过
$c_r-c_{r-1}<\delta$，所以对应的函数增量绝对值之和小于 $1$。把 $R$ 段估计相加得
\[
 \sum_{j=1}^N|F(t_j)-F(t_{j-1})|<R.
\]
右侧与原分割无关，故 $F$ 有界变差。令 $V(t)$ 为 $F$ 在 $[a,t]$ 上的全变差。
变差对相邻区间可加：把一个分割限制到两个子区间给出一个方向，把两边任意接近上确界的分割
在公共端点拼接给出反向。所以对 $s<t$，
\[
  V(t)-V(s)=\operatorname{Var}(F;[s,t])\ge |F(t)-F(s)|,
\]
故 $(V+F)/2$ 与 $(V-F)/2$ 都单调不减。绝对连续性还说明 $V$ 连续：若 $t-s<\delta$，
$[s,t]$ 内任意分割的子区间总长度小于 $\delta$，所以 $V(t)-V(s)$ 小于相应的
$\varepsilon$。令 $G_+=(V+F)/2$、$G_-=(V-F)/2$，则二者连续且单调不减。在
$[a,b]$ 外将它们作常值延拓，第~2.3 节的 Stieltjes 构造给出有限测度
$\mu_+$、$\mu_-$。令 $\nu_F=\mu_+-\mu_-$，则对 $a\le s<t\le b$，
\[
 \nu_F((s,t])
 =[G_+(t)-G_+(s)]-[G_-(t)-G_-(s)]
 =F(t)-F(s).
\]
还需核对 $\nu_F\ll m$。测度 $\tau=\mu_++\mu_-$ 支配 $|\nu_F|$，且
\[
 \tau((s,t])=V(t)-V(s)=\operatorname{Var}(F;[s,t]).
\]
给定 $\varepsilon>0$，取绝对连续性对应的 $\delta$。若 $m(E)=0$，取开集
$O\supset E$ 使 $m(O)<\delta$，并把 $O\cap[a,b]$ 写成可数个两两不交的相对开区间
$(I_j)$ 之并，并把 $I_j$ 在 $[a,b]$ 中的两个端点记为 $\alpha_j<\beta_j$。对前 $N$ 个分支，
在每个 $[\alpha_j,\beta_j]$ 内取任意有限分割；所有分割子区间的内部仍两两不交，
且总长度不超过 $m(O)<\delta$。令每个分支内的分割逼近相应变差的上确界，绝对连续性因而给出
\[
 \sum_{j=1}^N\operatorname{Var}(F;[\alpha_j,\beta_j])\le\varepsilon.
\]
由于 $G_\pm$ 连续，它们的 Stieltjes 测度在分支端点没有原子，所以
\[
 \tau(I_j)=\operatorname{Var}(F;[\alpha_j,\beta_j]).
\]
令 $N\to\infty$，再用测度的从下连续性，得
$\tau(O\cap[a,b])=\sum_j\tau(I_j)\le\varepsilon$。由于常值延拓不在端点产生原子，
\[
 |\nu_F|(E)\le\tau(E)\le\tau(O\cap[a,b])\le\varepsilon.
\]
令 $\varepsilon\downarrow0$ 得 $|\nu_F|(E)=0$。于是 $\nu_F$ 的 Jordan 两部分都关于 $m$
绝对连续。分别应用 RN 定理并作差，
得到 $h\in L^1([a,b])$，使
\[
  F(t)-F(s)=\int_s^t h(x)\,\dd x.
\]
反向也成立：若 $h\in L^1$，先取 $M>0$ 使
$\int_{\{|h|>M\}}|h|<\varepsilon/2$，再令
$\delta=\varepsilon/(2M)$。若 $I_1,\dots,I_N\subset[a,b]$ 两两不交且
$\sum_{k=1}^N|I_k|<\delta$，记 $A=\bigcup_{k=1}^NI_k$，则
\[
 \begin{aligned}
 \sum_{k=1}^N\left|\int_{I_k}h\,\dd m\right|
 &\le \int_A|h|\,\dd m\\
 &\le \int_{A\cap\{|h|\le M\}}M\,\dd m
      +\int_{\{|h|>M\}}|h|\,\dd m\\
 &\le M\sum_{k=1}^N|I_k|+\frac{\varepsilon}{2}
 <M\delta+\frac{\varepsilon}{2}=\varepsilon.
 \end{aligned}
\]
因此其不定积分绝对连续。第~4 章的微分定理还将进一步识别 $h=F'$ a.e.

这里的边界值得单独记住：绝对连续是 RN 表示的必要条件；密度只在 a.e. 意义下确定；在无限总质量
情形，证明必须落到同时具有有限 $\mu$、$\nu$ 质量的可数分块上，不能对两个无穷总质量作差。

\begin{exercise}[混合测度的分解]\label{exercise:3-3-1-1}
令 $\nu=\tfrac13m|_{[-1,1]}+2\delta_0+\delta_2$。分别求它关于
$m$ 以及关于 $m+\delta_0$ 的 Lebesgue 分解和 RN 密度。
\end{exercise}

\begin{exercise}[RN 导数的零集]\label{exercise:3-3-1-2}
设 $\nu\ll\mu$ 且 $h=\dd\nu/\dd\mu$。证明
$\nu\perp\mu$ 当且仅当 $h=0$ $\mu$-a.e.；解释为什么在两个测度都非零时，这与
$\nu\ll\mu$ 同时成立会造成矛盾。
\end{exercise}

\begin{exercise}[链式法则]\label{exercise:3-3-1-3}
不直接引用命题~\ref{proposition:3-3-1-4}，从简单函数的密度换算公式出发重证链式法则，并说明
三个导数各自在哪个测度的 a.e. 意义下定义。
\end{exercise}

\begin{exercise}[复测度的模长]\label{exercise:3-3-1-4}
设 $\nu=\sum_{j=1}^Nc_j\delta_{x_j}$，其中 $x_j$ 两两不同。计算 $|\nu|$ 和极分解中的 $u$，
并直接用有限分割定义验证答案。
\end{exercise}

密度函数已经出现；接下来需要一个能够精确衡量它们大小、容纳极限并支持积分配对的函数空间。


\subsection{\texorpdfstring{$L^p$}{Lp} 不等式、完备性与几何}\label{subsec:3-3-2}
\index{$L^p$ 空间}\index{Holder@\Holder{} 不等式}\index{Minkowski 不等式}

点态上很高但支撑很窄的函数，究竟算“大”还是“小”？答案取决于我们给尖峰多少惩罚。对
$\alpha\in\R$、$\beta>0$，先看
\[
  f_n=n^\alpha\1_{(0,n^{-\beta})}\qquad((0,1)\text{ 上}).
\]
若以 $p$ 次平均衡量它，则
\begin{equation}\label{eq:3-3-2-spike-norm}
  \int_0^1|f_n|^p\,\dd x=n^{\alpha p-\beta},
  \qquad
  \|f_n\|_p=n^{\alpha-\beta/p}.
\end{equation}
同一列函数可以在一个指数下趋于零、在另一个指数下发散。积分的代数因此要与指数的几何配合起来。

\begin{definition}[$L^p$ 空间]\label{def:3-3-2-1}
对 $1\le p<\infty$，令 $L^p(X,\mathcal A,\mu)$ 为满足
$\int|f|^p\,\dd\mu<\infty$ 的可测实（或复）函数关于 a.e. 相等的商空间，并定义
\[
  \|f\|_p=\left(\int_X|f|^p\,\dd\mu\right)^{1/p}.
\]
\end{definition}

\begin{definition}[$L^\infty$]\label{def:3-3-2-2}
可测函数 $f$ 的本质上确界为
\[
 \|f\|_\infty
 =\esssup_X|f|
 =\inf\{M\ge0:|f|\le M\ \mu\text{-a.e.}\}.
\]
本质有界函数关于 a.e. 相等的商空间记作 $L^\infty$。
\end{definition}

\begin{definition}[共轭指数]\label{def:3-3-2-3}
若 $p,q\in[1,\infty]$ 且 $1/p+1/q=1$，则称 $p,q$ 互为共轭指数；约定
$1/\infty=0$。
\end{definition}

\begin{definition}[Banach 空间]\label{def:3-3-2-banach}
若范数空间中的每个范数 Cauchy 序列都收敛到该空间中的元素，则称它为 Banach 空间。
\end{definition}

有限维范数空间以及配以上确界范数的 $C(K)$ 都是 Banach 空间。完备性只保证 Cauchy 近似有极限，
并不提供有界序列的收敛子列：$\ell^2$ 中的标准单位向量满足
$\|e_j-e_k\|_2=\sqrt2$（$j\ne k$），所以它们甚至没有 Cauchy 子列。

$L^2$ 的早期发展与 Fourier 系数及 Hilbert 空间紧密相连，随后 Riesz 把这种平方可积几何
推广到 $L^p$ 及其对偶。把 $L^p$ 只描述成 Riemann 可积函数的完备化，会遮住测度密度与积分配对
在其形成中的作用。不等式与完备性给出空间的几何结构，对偶表示则识别其中所有连续线性观测。

\begin{example}[尖峰的 $p$ 敏感性]\label{ex:3-3-2-1}
在 \eqref{eq:3-3-2-spike-norm} 中取 $\alpha=1/3$、$\beta=1$。则
\[
 \|f_n\|_2=n^{-1/6}\longrightarrow0,
 \qquad
 \|f_n\|_3=1,
 \qquad
 \|f_n\|_4=n^{1/12}\longrightarrow\infty.
\]
所以“支撑趋于零”并不单独决定平均大小；振幅与指数共同决定结果。
\end{example}

\begin{lemma}[Young 不等式]\label{lemma:3-3-2-1}
设 $a,b\ge0$，$1<p,q<\infty$ 且 $1/p+1/q=1$，则
\[
  ab\le\frac{a^p}{p}+\frac{b^q}{q}.
\]
\end{lemma}

\begin{proof}
固定 $b\ge0$，考虑 $\Phi(t)=t^p/p-bt$，$t\ge0$。若 $b=0$，所需不等式化为
$0\le a^p/p$。若 $b>0$，
$\Phi'(t)=t^{p-1}-b$，故唯一极小点为
$t_0=b^{1/(p-1)}=b^{q-1}$。由于 $(q-1)p=q$ 且 $q/p=q-1$，
\[
 \Phi(t)\ge\Phi(t_0)
 =\frac{b^q}{p}-b^q=-\frac{b^q}{q}.
\]
令 $t=a$ 并移项即得结论。
\end{proof}

\begin{theorem}[\Holder{} 与 Minkowski]\label{theorem:3-3-2-2}
若 $p,q\in[1,\infty]$ 互为共轭指数，$f\in L^p$、$g\in L^q$，则
\[
  \|fg\|_1\le\|f\|_p\|g\|_q.
  \tag{H}\label{eq:3-3-2-holder}
\]
若 $1\le p\le\infty$ 且 $f,g\in L^p$，则
\[
  \|f+g\|_p\le\|f\|_p+\|g\|_p.
  \tag{M}\label{eq:3-3-2-minkowski}
\]
\end{theorem}

\begin{proof}
先证 \Holder{}。若 $1<p,q<\infty$，且某个范数为零，则相应函数 a.e. 为零，结论成立。
其余情形令 $F=|f|/\|f\|_p$、$G=|g|/\|g\|_q$。逐点应用 Young 不等式并积分，得到
\[
 \int FG\,\dd\mu
 \le\frac1p\int F^p\,\dd\mu+\frac1q\int G^q\,\dd\mu
 =\frac1p+\frac1q=1.
\]
乘回两个范数即得 \eqref{eq:3-3-2-holder}。若 $(p,q)=(1,\infty)$，则
$|fg|\le\|g\|_\infty|f|$ a.e.；积分即可。$(p,q)=(\infty,1)$ 对称。

再证 Minkowski。$p=1$ 时由点态三角不等式积分；$p=\infty$ 时，由
$|f+g|\le\|f\|_\infty+\|g\|_\infty$ a.e. 得到结论。设 $1<p<\infty$，并令
$q=p/(p-1)$。不等式
\[
 |f+g|^p\le2^{p-1}(|f|^p+|g|^p)
\]
先保证 $f+g\in L^p$。若 $\|f+g\|_p=0$，则
$0=\|f+g\|_p\le\|f\|_p+\|g\|_p$；以下设该范数非零。由
\[
 |f+g|^p
 \le |f|\,|f+g|^{p-1}+|g|\,|f+g|^{p-1}
\]
以及已经证明的 \Holder{} 不等式，
\[
 \begin{aligned}
 \|f+g\|_p^p
 &\le(\|f\|_p+\|g\|_p)
       \bigl\||f+g|^{p-1}\bigr\|_q\\
 &=(\|f\|_p+\|g\|_p)\|f+g\|_p^{p-1}.
 \end{aligned}
\]
除以非零的 $\|f+g\|_p^{p-1}$ 得到 \eqref{eq:3-3-2-minkowski}。整个推导只在确认左侧有限后
才使用它，因而没有循环调用三角不等式。
\end{proof}

\begin{theorem}[Riesz--Fischer 完备性]\label{theorem:3-3-2-3}
对任意测度空间以及每个 $1\le p\le\infty$，$L^p$ 都是 Banach 空间。若
$f_n\to f$ 于 $L^p$，则存在子列 $(f_{n_k})$ 使
$f_{n_k}(x)\to f(x)$ 对 $\mu$-a.e. 的 $x$ 成立。
\end{theorem}

\begin{proof}
先设 $1\le p<\infty$，并令 $(f_n)$ 为 $L^p$ 中的 Cauchy 序列。递归选取严格递增的指标
$n_k$，使
\[
  \|f_{n_{k+1}}-f_{n_k}\|_p\le2^{-k}.
\]
记 $g_k=|f_{n_{k+1}}-f_{n_k}|$，$S_N=\sum_{k=1}^Ng_k$。Minkowski 不等式给出
\[
  \|S_N\|_p\le\sum_{k=1}^N2^{-k}\le1.
\]
因为 $S_N\uparrow S:=\sum_{k=1}^\infty g_k$，单调收敛定理给出
\[
  \int S^p\,\dd\mu
  =\lim_N\int S_N^p\,\dd\mu\le1.
\]
所以 $S<\infty$ a.e.，增量级数
\[
 f_{n_1}+\sum_{k=1}^\infty(f_{n_{k+1}}-f_{n_k})
\]
在这些点绝对收敛。它的有限部分和都是可测函数，而收敛点集
$\{S<\infty\}$ 也可测；以其和定义 $f$，并在补集上令 $f=0$，便得到可测函数。若
$H_k=\sum_{j=k}^\infty g_j$。对有限尾和使用 Minkowski，再对其 $p$ 次方用 MCT，得
\[
 \begin{aligned}
 \|H_k\|_p^p
 &=\lim_{N\to\infty}
   \left\|\sum_{j=k}^Ng_j\right\|_p^p\\
 &\le\lim_{N\to\infty}
   \left(\sum_{j=k}^N\|g_j\|_p\right)^p
 \le\left(\sum_{j=k}^{\infty}2^{-j}\right)^p.
 \end{aligned}
\]
因而 $\|H_k\|_p\le\sum_{j=k}^\infty2^{-j}=2^{1-k}$。
又有 $|f-f_{n_k}|\le H_k$ a.e.，故 $f-f_{n_k}\in L^p$ 且
$\|f-f_{n_k}\|_p\le2^{1-k}$。由 $f_{n_k}\in L^p$ 及
$f=f_{n_k}+(f-f_{n_k})$，还得到 $f\in L^p$。原序列为 Cauchy，因此给定 $\varepsilon>0$，先取 $N$，使
$m,n\ge N$ 时 $\|f_n-f_m\|_p<\varepsilon/2$；再选 $k$ 使 $n_k\ge N$ 且
$\|f-f_{n_k}\|_p<\varepsilon/2$。于是对每个 $n\ge N$，三角不等式给出
$\|f_n-f\|_p<\varepsilon$，故 $f_n\to f$ 于 $L^p$。

若 $p=\infty$，同样选取快速子列。对每个 $k$ 选代表后，存在零集 $N_k$，使
\[
 |f_{n_{k+1}}(x)-f_{n_k}(x)|\le2^{-k}\qquad(x\notin N_k).
\]
删去可数并 $N=\bigcup_kN_k$ 后，增量级数由 Weierstrass 判别法在 $X\setminus N$ 上一致收敛；
以其和定义 $f$，并在 $N$ 上令 $f=0$。于是
\[
  |f(x)-f_{n_k}(x)|\le2^{1-k}\qquad(x\notin N).
\]
再取零集 $N_0'$，使
$|f_{n_1}(x)|\le\|f_{n_1}\|_\infty+1$ 于 $X\setminus N_0'$。于是对
$x\notin N\cup N_0'$，
\[
 |f(x)|\le |f_{n_1}(x)|+\sum_{k=1}^\infty2^{-k}
 \le\|f_{n_1}\|_\infty+2.
\]
因此 $f\in L^\infty$，而上面的尾估计还给出
$\|f-f_{n_k}\|_\infty\le2^{1-k}$。给定 $\varepsilon>0$，先由原列的 Cauchy 性取
$N_0$，使 $m,n\ge N_0$ 时 $\|f_n-f_m\|_\infty<\varepsilon/2$；再选 $k$ 使
$n_k\ge N_0$ 且 $\|f-f_{n_k}\|_\infty<\varepsilon/2$。于是对每个 $n\ge N_0$，三角不等式给出
$\|f_n-f\|_\infty<\varepsilon$，
故整列收敛于 $f$。

最后设 $f_n\to f$ 于 $L^p$。当 $p<\infty$ 时选子列使
$\|f_{n_k}-f\|_p\le2^{-k}$。令
$G_N=\sum_{k=1}^N|f_{n_k}-f|$，则 Minkowski 不等式给出
\[
 \|G_N\|_p\le\sum_{k=1}^N2^{-k}\le1.
\]
由 $G_N\uparrow G:=\sum_{k=1}^\infty|f_{n_k}-f|$ 及单调收敛，$G\in L^p$，故
$G(x)<\infty$ a.e.；收敛级数的通项趋于零，即 $f_{n_k}(x)\to f(x)$ a.e.
当 $p=\infty$ 时，也选取 $\|f_{n_k}-f\|_\infty\le2^{-k}$。对每个 $k$ 存在零集 $Z_k$，使
$|f_{n_k}(x)-f(x)|\le2^{-k}$ 于 $X\setminus Z_k$；在
$X\setminus\bigcup_kZ_k$ 上直接令 $k\to\infty$，同样得到逐点收敛。
\end{proof}

完备性与依测度收敛是两件不同的事；上面的“快速子列”把范数信息转成可和的逐点增量，正是二者
之间可重复使用的桥梁。

\begin{proposition}[Jensen 不等式]\label{proposition:3-3-2-4}
设 $\mu(X)=1$，$I\subset\R$ 是区间，$f$ 是取值于 $I$ 的可积实函数，
$\varphi:I\to\R$ 为凸函数，且 $\varphi\circ f$ 可积。则
\[
  \varphi\left(\int_Xf\,\dd\mu\right)
  \le\int_X\varphi(f)\,\dd\mu.
\]
\end{proposition}

\begin{proof}
记 $m=\int f\,\dd\mu$。先说明 $m\in I$。若 $a=\inf I$ 是有限左端点，则
$f\ge a$ a.e.，所以 $m\ge a$；如果 $a\notin I$ 而 $m=a$，则 $f-a>0$ a.e.，零积分判据却会从
$\int(f-a)=0$ 推出 $f=a$ a.e.，矛盾。若 $b=\sup I$ 是有限右端点，则
$f\le b$ a.e. 给出 $m\le b$；如果 $b\notin I$ 而 $m=b$，则 $b-f>0$ a.e.，但
$\int(b-f)=0$ 同样违反零积分判据。无穷端点不可能等于有限实数 $m$。
因此 $m$ 要么是 $I$ 的内点，要么是属于 $I$ 的有限端点。先设 $m$ 是内点。凸函数在内点有支撑直线：存在
$c\in\R$ 使
\[
  \varphi(t)\ge\varphi(m)+c(t-m),
  \qquad t\in I.
\]
具体地，可取任意
\[
  c\in\left[
  \sup_{t<m}\frac{\varphi(m)-\varphi(t)}{m-t},
  \inf_{t>m}\frac{\varphi(t)-\varphi(m)}{t-m}
  \right];
\]
凸性保证左端不超过右端。令 $t=f(x)$ 并积分，得到
\[
 \int\varphi(f)\,\dd\mu
 \ge\varphi(m)+c\int(f-m)\,\dd\mu=\varphi(m).
\]

若 $m$ 是 $I$ 中的左端点 $a$，则 $f-a\ge0$ 且积分为零，故 $f=a$ a.e.，等式成立。若 $m$ 是
$I$ 中的右端点 $b$，则 $b-f\ge0$ 且积分为零，故 $f=b$ a.e.，等式也成立。
\end{proof}

两点的总质量为一的测度空间给出最直接的检验：令
$X=\{0,1\}$、$\mu=\theta\delta_0+(1-\theta)\delta_1$，并取
$f(0)=a$、$f(1)=b$。Jensen 正好
化为
\[
 \varphi(\theta a+(1-\theta)b)
 \le\theta\varphi(a)+(1-\theta)\varphi(b),
\]
即凸性的定义。

\begin{example}[有限测度嵌入]\label{ex:3-3-2-2}
若 $\mu(X)<\infty$ 且 $q>p$，则 $L^q\subset L^p$。当 $q<\infty$ 时，对
$|f|^p$ 与 $1$ 使用指数 $q/p$、$q/(q-p)$ 的 \Holder{} 不等式，得到
\[
 \int|f|^p\,\dd\mu
 \le\left(\int|f|^q\,\dd\mu\right)^{p/q}
      \mu(X)^{1-p/q}.
\]
开 $p$ 次方即
\[
  \|f\|_p\le\mu(X)^{1/p-1/q}\|f\|_q.
\]
$q=\infty$ 时由 $|f|\le\|f\|_\infty$ a.e. 直接得到同一公式。

无限测度空间上两个方向都可能失败。设 $q>p$ 且 $q<\infty$。在 $(0,\infty)$ 上取
\[
 f(x)=x^{-a}\1_{(1,\infty)}(x),
 \qquad \frac1q<a\le\frac1p.
\]
因为 $\int_1^\infty x^{-ar}\,\dd x<\infty$ 当且仅当 $ar>1$，所以
$f\in L^q\setminus L^p$。另一方面，取
\[
 g(x)=x^{-b}\1_{(0,1)}(x),
 \qquad \frac1q\le b<\frac1p.
\]
由 $\int_0^1x^{-br}\,\dd x<\infty$ 当且仅当 $br<1$，可得
$g\in L^p\setminus L^q$。若 $q=\infty$，取 $0<a\le1/p$ 的第一个尾部函数，它属于
$L^\infty\setminus L^p$；再取 $0<b<1/p$ 的第二个原点函数，它属于
$L^p\setminus L^\infty$。
\end{example}

\begin{example}[本质上确界]\label{ex:3-3-2-3}
在 $[0,1]$ 上定义
\[
 f(x)=1\quad(x\ne0),
 \qquad f(0)=10^{100}.
\]
则 $\|f\|_\infty=1$，因为异常值只落在零测单点上。相反，在计数测度下单点不是零集，同一修改
会把 $\ell^\infty$ 范数改成 $10^{100}$。所以 $L^\infty$ 的元素是等价类；未经选择代表，点值
$f(x)$ 并不是空间元素所固有的数据。
\end{example}

当 $0<p<1$ 时，公式 $(\int|f|^p)^{1/p}$ 一般不满足三角不等式。若 $A,B$ 是两个质量为 $1$ 的
不交集合，则
\[
 \|\1_A+\1_B\|_p=2^{1/p}>2=\|\1_A\|_p+\|\1_B\|_p.
\]
这个例子的单位球甚至不是凸集：$\1_A$、$\1_B$ 在单位球中，而它们的中点范数为
$2^{1/p-1}>1$。因此在一般测度空间上，这一范围只能得到拟范数结构；单一原子之类的
退化有限维情形可能例外。以下 $L^p$ 理论均取 $1\le p\le\infty$。

\begin{figure}[htbp]
  \centering
  \begin{tikzpicture}[scale=1.1]
    \draw[->,BookInk!70] (-2.1,0)--(2.1,0) node[right] {$x$};
    \draw[->,BookInk!70] (0,-2.1)--(0,2.1) node[above] {$y$};
    \draw[BookWarning,semithick] (-1.65,-1.65) rectangle (1.65,1.65);
    \draw[BookAccent,semithick] (0,0) circle (1.65);
    \draw[BookWarm,semithick] (0,1.65)--(1.65,0)--(0,-1.65)--(-1.65,0)--cycle;
    \node[book label,text=BookWarm] at (1.08,0.72) {$p=1$};
    \node[book label,text=BookAccent] at (1.55,1.02) {$p=2$};
    \node[book label,text=BookWarning] at (1.72,-1.83) {$p=\infty$};
  \end{tikzpicture}
  \caption{二维 $\ell^p$ 单位球的截面。指数越大，坐标尖峰受到的相对惩罚越弱；有限维单位球的
  形状本身不蕴含无限维空间中的紧性。}
  \label{fig:3-3-2-lp-balls}
\end{figure}

这些结论很快会在几个方向中反复出现。$L^2$ 的完备性使平方可积函数成为 Hilbert 空间，并为
Fourier 的 Plancherel 理论提供容纳极限的空间；\Holder{} 不等式控制参数积分、卷积和两个密度的
配对；Riesz--Fischer 的快速子列方法还可用于证明 Banach 值 $L^p$ 空间的完备性。

\begin{figure}[htbp]
  \centering
  \begin{tikzpicture}[node distance=8mm and 8mm]
    \node[book strong node] (f1) {$f_{n_1}$};
    \node[book node,right=of f1] (f2) {$f_{n_2}$};
    \node[book node,right=of f2] (f3) {$f_{n_3}$};
    \node[book node,right=of f3] (dots) {$\cdots$};
    \node[book warm node,right=of dots] (f) {$f$};
    \draw[book accent arrow] (f1)--node[above,book label] {$\|\Delta_1\|_p\le2^{-1}$} (f2);
    \draw[book accent arrow] (f2)--node[above,book label] {$\|\Delta_2\|_p\le2^{-2}$} (f3);
    \draw[book accent arrow] (f3)-- (dots);
    \draw[book arrow] (dots)--node[above,book label] {$\sum|\Delta_k|<\infty$ a.e.} (f);
  \end{tikzpicture}
  \caption{快速 Cauchy 子列把范数增量压成可和级数；点态绝对收敛产生候选极限，尾和估计再把它
  拉回 $L^p$ 范数。}
  \label{fig:3-3-2-fast-cauchy}
\end{figure}

\begin{exercise}[\Holder{} 的等号]\label{exercise:3-3-2-1}
设 $1<p,q<\infty$，$f\in L^p$、$g\in L^q$ 均非零。沿 Young 不等式的等号条件证明：
\Holder{} 取等号当且仅当存在常数 $c>0$ 使 $|f|^p=c|g|^q$ a.e.，并另行说明复相位对
$\int fg$ 取模等号的影响。
\end{exercise}

\begin{exercise}[$L^p$ 收敛与依测度收敛]\label{exercise:3-3-2-2}
若 $1\le p<\infty$ 且 $f_n\to f$ 于 $L^p$，证明对每个 $\varepsilon>0$，
\[
 \mu(|f_n-f|>\varepsilon)
 \le\varepsilon^{-p}\|f_n-f\|_p^p\longrightarrow0.
\]
给出一个依测度收敛但不在 $L^p$ 收敛的尖峰例。
\end{exercise}

\begin{exercise}[$L^\infty$ 的完备性]\label{exercise:3-3-2-3}
把 Riesz--Fischer 定理的 $p=\infty$ 证明改写为不抽取快速子列的直接证明：先为 Cauchy 列选
一串误差阈值，再删去一个可数并零集，证明代表在该零集外一致 Cauchy。
\end{exercise}

\begin{exercise}[有限测度中的严格性]\label{exercise:3-3-2-4}
在 $([0,1],m)$ 上，对任意 $1\le p<q<\infty$ 构造
$f\in L^p\setminus L^q$。再证明若 $\mu(X)<\infty$ 且 $|f|$ a.e. 为常数，则有限测度嵌入中的
范数估计取等号。
\end{exercise}

$L^p$ 已经是完备的几何空间。还剩两个识别问题：可否用简单或连续函数逼近其中每个元素？所有连续
线性观测是否都来自积分配对？第一个问题由稠密性回答，第二个问题最终又回到 RN 密度。


\subsection{稠密性、可分性与 \texorpdfstring{$L^p$}{Lp} 对偶}\label{subsec:3-3-3}
\index{Hahn--Banach 定理}\index{$L^p$ 对偶}\index{可分性!$L^p$}

抽象的 a.e. 等价类不适合直接计算。我们希望先在有限分割上的简单函数中完成计算，再由范数极限
进入整个空间。在有拓扑的 Radon 空间上，还希望把测试函数换成 $C_c$。另一方面，一个线性观测
若只在某个子空间上给出，也需要知道它能否在不增大范数的前提下延到全空间。这两条路线分别通向
稠密性和 Hahn--Banach 定理，并在 $L^p$ 对偶表示中会合。

Riesz 的 $L^p$ 对偶理论与矩问题一同发展：线性观测先由有限测试确定，再被一个可积密度统一表示。
这也解释了下面证明的顺序。我们不会先假设抽象对偶已经可表示，而是让泛函在示性函数上生成测度，
随后调用刚刚证明的 RN 定理。

\begin{definition}[连续对偶与对偶范数]\label{def:3-3-3-dual}
范数空间 $E$ 的连续对偶 $E^*$ 是所有连续线性映射
$\Lambda:E\to\C$（实空间中以 $\R$ 为值域）组成的空间，范数为
\[
  \|\Lambda\|=\sup_{\|x\|\le1}|\Lambda(x)|.
\]
线性泛函连续当且仅当存在 $C<\infty$ 使
$|\Lambda(x)|\le C\|x\|$ 对所有 $x$ 成立；满足此式的最小 $C$ 就是 $\|\Lambda\|$。
\end{definition}

\begin{definition}[积分配对]\label{def:3-3-3-1}
对共轭指数 $p,q$，定义
\[
  \langle f,g\rangle=\int_X f\overline g\,\dd\mu
\]
（实值情形去掉共轭）。全书约定内积或配对对第一变量线性。因此
$\Lambda_g(f)=\langle f,g\rangle$ 关于 $f$ 线性，而 \Holder{} 不等式给出
$|\Lambda_g(f)|\le\|f\|_p\|g\|_q$。
\end{definition}

\begin{definition}[测试类稠密]\label{def:3-3-3-2}
若子空间 $D\subset L^p$ 满足：对每个 $f\in L^p$ 与 $\varepsilon>0$，存在 $d\in D$ 使
$\|f-d\|_p<\varepsilon$，则称 $D$ 在 $L^p$ 中稠密。若 $D$ 最初由点值函数组成，这里所说的是
它们经 a.e. 商映射所得的像。
\end{definition}

\begin{definition}[对偶表示映射]\label{def:3-3-3-3}
对 $g\in L^q$，令
\[
  J(g)=\Lambda_g,
  \qquad \Lambda_g(f)=\int f\overline g\,\dd\mu.
\]
\Holder{} 保证 $J:L^q\to(L^p)^*$ 有定义。对偶问题是判断 $J$ 是否等距、是否满射。
\end{definition}

在讨论具体的 $L^p$ 前，先解决一般范数空间中的延拓障碍。给定 $x_0\ne0$，直线
$M=\Span\{x_0\}$ 上的公式
$f(\alpha x_0)=\alpha\|x_0\|$ 范数为 $1$；无限维中却没有可供使用的正交投影把它延出去。
一维扩张可以算清，但把所有一维扩张拼成全空间需要一个真正的选择原理。

另一个更尖锐的问题已经出现在 $\ell^\infty$。在
$c_0+\Span\{\mathbf1\}$ 上取出“常数尾巴”的泛函
$x+a\mathbf1\mapsto a$；由 $x_n\to0$ 可知它有界。若它能延到整个 $\ell^\infty$，就会产生一个
不能由 $\ell^1$ 求和配对表示的观测。下面的定理先保证延拓存在，定理后再完成这个反例。

\begin{theorem}[Hahn--Banach 延拓定理]\label{theorem:3-3-3-hahn-banach}
设 $E$ 为实向量空间，$P:E\to\R$ 为次线性函数，即
\[
 P(x+y)\le P(x)+P(y),
 \qquad P(tx)=tP(x)\quad(t\ge0).
\]
若 $M\subset E$ 为线性子空间，线性泛函 $f:M\to\R$ 满足 $f(m)\le P(m)$，则存在
$E$ 上的线性泛函 $F$ 延拓 $f$，并满足 $F(x)\le P(x)$。特别地，实或复范数空间子空间上的
有界线性泛函都可保持范数延拓到全空间。
\end{theorem}

\begin{proof}
先做一维扩张。设 $x_0\notin M$，希望在 $M+\R x_0$ 上定义
\[
  F_\alpha(m+tx_0)=f(m)+t\alpha.
\]
当 $t=1$ 与 $t=-1$ 时，受 $P$ 控制所要求的上下界分别汇成
\[
 L:=\sup_{m\in M}\{f(m)-P(m-x_0)\}
 \le \alpha\le
 U:=\inf_{n\in M}\{P(n+x_0)-f(n)\}.
\]
先核对两个端点之间确实有实数。在下面的不等式中取 $n=0$ 可得
$f(m)-P(m-x_0)\le P(x_0)$，故 $L<\infty$；取 $m=0$ 可得
$P(n+x_0)-f(n)\ge-P(-x_0)$，故 $U>-\infty$。另一方面，对任意 $m,n\in M$，
\[
 f(m)+f(n)=f(m+n)
 \le P(m+n)
 \le P(m-x_0)+P(n+x_0),
\]
故 $f(m)-P(m-x_0)\le P(n+x_0)-f(n)$；先对 $m$ 取上确界、再对 $n$ 取下确界便得
$L\le U$。因而 $[L,U]$ 含有实数，取 $\alpha\in[L,U]$。若 $t>0$，由上界应用于 $m/t$，
\[
 f(m)+t\alpha
 \le tP(m/t+x_0)=P(m+tx_0).
\]
若 $t<0$，写 $s=-t>0$，由下界应用于 $m/s$，
\[
 f(m)-s\alpha
 \le sP(m/s-x_0)=P(m-sx_0).
\]
$t=0$ 正是原假设。因此 $F_\alpha\le P$。

令 $\mathscr E$ 为所有受 $P$ 控制的线性延拓对 $(N,g)$，其中
$M\subset N\subset E$，并按定义域包含且取值相容排序。任意链的定义域之并仍为线性子空间；链上
泛函相容，所以其并给出该链的上界。Zorn 引理于是给出极大元 $(N_*,g_*)$。若
$N_*\ne E$，选 $x_0\in E\setminus N_*$，刚才的一维步骤会把 $g_*$ 继续延拓，违背极大性。
所以 $N_*=E$。这里使用 Zorn 引理，因而使用了相应的选择公理形式；无限次延拓并不是一次有限递归。

若 $E$ 为实范数空间且 $|f(m)|\le C\|m\|$，取 $P(x)=C\|x\|$。定理给出
$F(x)\le C\|x\|$；把 $x$ 换成 $-x$ 又得 $-F(x)\le C\|x\|$，故
$|F(x)|\le C\|x\|$。取 $C=\|f\|$ 可保持范数。

最后设 $E$ 为复范数空间。把 $g=\Re f$ 看成 $M$ 上的实线性泛函。它满足
$\|g\|=\|f\|$：一个方向来自 $|\Re f(x)|\le|f(x)|$；反向则对每个 $x$ 乘一个模为 $1$ 的
标量，使 $f(x)$ 旋转到非负实轴。由实情形把 $g$ 保范延拓为实线性泛函 $G$，再定义
\[
  F(x)=G(x)-iG(ix).
\]
实线性直接给出 $F(ix)=iF(x)$，故 $F$ 复线性；对 $x\in M$，
$G(ix)=\Re(if(x))=-\Im f(x)$，所以 $F(x)=f(x)$。给定 $x$，取 $|\zeta|=1$ 使
$\zeta F(x)=|F(x)|$。复线性与 $\Re F(y)=G(y)$ 给出
\[
 |F(x)|=\Re F(\zeta x)=G(\zeta x)
 \le\|f\|\,\|x\|.
\]
于是复情形也保持范数。
\end{proof}

\begin{corollary}[一点范数化与子空间分离]\label{corollary:3-3-3-norming}
若 $x_0\ne0$ 属于范数空间 $E$，则存在 $\phi\in E^*$ 使
\[
  \|\phi\|=1,
  \qquad \phi(x_0)=\|x_0\|.
\]
更一般地，子空间上的连续线性泛函可保持范数延拓到 $E$。
\end{corollary}

\begin{proof}
在 $\Span\{x_0\}$ 上定义 $\phi_0(\alpha x_0)=\alpha\|x_0\|$。它的范数为 $1$，由
Hahn--Banach 定理保范延拓即可。第二句就是定理的范数空间版本。
\end{proof}

例如，在 $(\R^2,\|\cdot\|_\infty)$ 中，对 $x_0=(1,-1)$ 取
\[
  \phi(a,b)=\frac{a-b}{2}.
\]
由 $|a-b|/2\le\max\{|a|,|b|\}$ 可知 $\|\phi\|\le1$；在 $(1,-1)$ 处取等，故
$\|\phi\|=1$ 且 $\phi(x_0)=1=\|x_0\|_\infty$。这里支撑超平面可以写成公式，定理保证无限维中
也有同样的范数化观测。

现在可以闭合 $L^\infty$ 对偶的第一个失败现场。

\begin{example}[$L^\infty$ 的额外对偶]\label{ex:3-3-3-3}
在 $\N$ 的计数测度上，$L^\infty=\ell^\infty$。令 $c_0$ 为趋于零的序列空间，
$\mathbf1=(1,1,\ldots)$，并在
$M=c_0+\Span\{\mathbf1\}$ 上定义
\[
  F_0(x+a\mathbf1)=a.
\]
由于 $\mathbf1\notin c_0$，分解唯一。又因为 $x_n\to0$，
\[
 |a|=\lim_n|x_n+a|\le\|x+a\mathbf1\|_\infty,
\]
而 $F_0(\mathbf1)=1$，所以 $\|F_0\|=1$。Hahn--Banach 定理把它保范延拓为
$F\in(\ell^\infty)^*$。若存在 $g\in\ell^1$ 使
$F(x)=\sum_nx_n\overline{g_n}$，则
$0=F(e_n)=\overline{g_n}$ 对每个 $n$ 成立，故 $g=0$；但这与
$F(\mathbf1)=1$ 矛盾。因此 $(L^\infty)^*$ 一般严格大于由 $L^1$ 配对产生的部分。
\end{example}

\begin{example}[阶梯函数]\label{ex:3-3-3-1}
设 $\mu$ 是 $\R$ 上的 Radon 测度，$E$ 是有限测度 Borel 集，$1\le p<\infty$。给定
$\varepsilon>0$，正则性给出紧集 $K$ 与开集 $U$，使
\[
  K\subset E\subset U,
  \qquad \mu(U\setminus K)<\varepsilon^p.
\]
在每个有界区间上，$\mu$ 的原子至多可数：质量至少 $1/r$ 的原子只有有限个，再对 $r$ 取可数并。
因此可为每个 $x\in K$ 选开区间 $(a_x,b_x)$，使
\[
  x\in(a_x,b_x),
  \qquad [a_x,b_x]\subset U,
  \qquad \mu(\{a_x\})=\mu(\{b_x\})=0.
\]
端点恰是局部分布函数的连续点。紧性给出有限子覆盖
$I=\bigcup_{j=1}^N(a_j,b_j)$。函数 $s=\1_I$ 是只有有限个跳点的阶梯函数，而且
$K\subset I\subset U$，故
\[
  \|\1_E-s\|_p^p
  =\mu(E\mathbin{\triangle} I)
  \le\mu(U\setminus K)<\varepsilon^p.
\]
这给出一次带明确误差的阶梯逼近；端点取连续点还保证改变开闭端点不改变其 $L^p$ 等价类。
\end{example}

一般的简单函数稠密只依赖测度，而从简单函数升级到连续函数需要 Radon 正则性和局部 Urysohn
截断。

\begin{theorem}[简单函数与 $C_c$ 稠密]\label{theorem:3-3-3-1}
设 $1\le p<\infty$。
\begin{enumerate}
  \item 在任意测度空间上，具有有限测度支撑的简单函数在 $L^p$ 中稠密；
  \item 若 $X$ 是局部紧 Hausdorff 空间，$\mu$ 是其上的 Radon 测度，则
        $C_c(X)$ 在 $L^p(\mu)$ 中的像稠密。
\end{enumerate}
\end{theorem}

\begin{proof}
先设 $f\in L^p$。令
\[
  E_n=\{1/n\le|f|\le n\},
  \qquad f_n=f\1_{E_n}.
\]
在 $E_n$ 上有 $n^{-p}\1_{E_n}\le |f|^p$，故积分单调性直接给出
$\mu(E_n)\le n^p\|f\|_p^p<\infty$。又有
$|f-f_n|^p\to0$ a.e. 且 $|f-f_n|^p\le|f|^p$，支配收敛定理给出
$f_n\to f$ 于 $L^p$。固定 $n$ 后，$f_n$ 有界且支撑测度有限。把其实部、虚部的有界值域分别
分成网格不超过 $1/k$ 的有限区间，并在每个原像上取相应网格值，得到有限值简单函数 $s_k$，使
\[
 |s_k-f_n|\le\frac{\sqrt2}{k}\1_{E_n}\quad\text{a.e.}
\]
（实值情形去掉 $\sqrt2$）。所以
$\|s_k-f_n\|_p\le\sqrt2\mu(E_n)^{1/p}/k\to0$，第一项得证。

现在设 $X$ 局部紧 Hausdorff，$\mu$ 为 Radon 测度。只需逼近一个有限测度集合 $E$ 的示性函数。
给定 $\varepsilon>0$，由内外正则性取紧集 $K$ 和开集 $U$，使
\[
  K\subset E\subset U,
  \qquad \mu(U\setminus K)<\varepsilon^p.
\]
局部 Urysohn 引理给出 $\phi\in C_c(X,[0,1])$，满足
$\phi=1$ 于 $K$ 且 $\supp\phi\subset U$。于是
\[
  |\1_E-\phi|\le\1_{U\setminus K},
  \qquad
  \|\1_E-\phi\|_p<\varepsilon.
\]
有限线性组合再配合三角不等式，可逼近每个有限测度支撑简单函数；第一项完成对任意
$f\in L^p$ 的逼近。
\end{proof}

特别地，在 $\R$ 上对简单函数的每个有限测度水平集使用例~\ref{ex:3-3-3-1}，再作有限线性组合，
可知有限区间示性函数张成的阶梯函数在任意 Radon 测度的 $L^p(\R)$ 中稠密。

\begin{example}[欧氏空间中的光滑逼近]\label{ex:3-3-3-2}
在 $\R^n$ 的 Lebesgue 测度下，上一定理已经给出 $C_c(\R^n)$ 的稠密性。取非负
$\rho\in C_c^\infty(\R^n)$，满足 $\int\rho=1$，并置
$\rho_\varepsilon(x)=\varepsilon^{-n}\rho(x/\varepsilon)$。预期的第二步是
\[
  f_\varepsilon=\rho_\varepsilon*f\in C_c^\infty,
  \qquad \|f_\varepsilon-f\|_p\longrightarrow0
  \quad(f\in C_c).
\]
这个极限由卷积、平移连续性和近似恒等族推出，完整证明见
定理~\ref{theorem:4-3-1-3}。一般 Radon 空间未必具有平移或光滑结构，所以其中适用的稠密类是
$C_c(X)$，不能无条件替换为 $C_c^\infty(X)$。
\end{example}

第二可数性把连续逼近进一步压缩成可数测试族。

\begin{proposition}[$L^p$ 可分性]\label{proposition:3-3-3-2}
若 $X$ 是第二可数的局部紧 Hausdorff 空间，$\mu$ 是 Radon 测度，且
$1\le p<\infty$，则 $L^p(\mu)$ 可分。
\end{proposition}

\begin{proof}
由第~1.4 节，可取相对紧开集组成的可数基 $\mathcal B$，以及紧耗尽
\[
 K_1\subset\interior{K_2}\subset K_2\subset\interior{K_3}\subset\cdots,
 \qquad \bigcup_nK_n=X.
\]
对每一对 $V,U\in\mathcal B$ 且 $\overline V\subset U$，由局部 Urysohn 引理选取
$\psi_{V,U}\in C_c(X,[0,1])$，满足 $\psi_{V,U}=1$ 于 $\overline V$ 且
$\supp\psi_{V,U}\subset U$。由 $K_n\subset\interior{K_{n+1}}$，对紧集 $K_n$ 与开集
$\interior{K_{n+1}}$ 再用一次局部 Urysohn 引理，选取 $\eta_n\in C_c(X,[0,1])$，使
\[
 \eta_n=1\text{ 于 }K_n,
 \qquad \supp\eta_n\subset\interior{K_{n+1}}\subset K_{n+1}.
\]
令 $\mathcal A_0$ 为这些函数生成的、不另添全空间
常数函数 $1$ 的有理系数代数；因而 $\mathcal A_0\subset C_c(X)$。复情形
允许实部、虚部均为有理数的系数，并加入所有生成元的共轭。于是 $\mathcal A_0$ 可数。它的一致闭包
对实数（复情形对复数）数乘封闭，因为任意标量都可由相应的有理标量逼近；故这个闭包是
Stone--Weierstrass 定理所需的实或复代数。

这些帽函数分离点。事实上，若 $x\ne y$，由 Hausdorff 正则性和基的性质可取
$x\in V$、$\overline V\subset U$ 且 $y\notin U$，于是
$\psi_{V,U}(x)=1$、$\psi_{V,U}(y)=0$。固定 $n$，把 $\mathcal A_0$ 限制到紧空间
$K_{n+1}$；该代数分离点，并且某个更大的 $\eta_m$ 在 $K_{n+1}$ 上恒为 $1$。实或复
Stone--Weierstrass 定理因此说明它在 $C(K_{n+1})$ 中一致稠密。

给定 $f\in C_c(X)$，取 $n$ 使 $\supp f\subset K_n$。对任意 $\delta>0$，可取
$a\in\mathcal A_0$ 使
\[
  \sup_{K_{n+1}}|a-f|<\delta.
\]
函数 $d=\eta_na$ 仍属于 $\mathcal A_0$，支撑包含在 $K_{n+1}$，且因为
$\eta_nf=f$，
\[
 \|d-f\|_p
 \le\delta\,\mu(K_{n+1})^{1/p}.
\]
Radon 测度在紧集上有限，所以可令右侧任意小。这证明可数集 $\mathcal A_0$ 在 $C_c$ 的
$L^p$ 像中稠密；再用定理~\ref{theorem:3-3-3-1}，它在整个 $L^p$ 中稠密。
\end{proof}

第二可数性不能从陈述中无声删除。若 $I$ 是不可数集，给它离散拓扑和计数测度，则该测度是 Radon，
而 $L^p=\ell^p(I)$ 中的单位向量族 $(e_i)_{i\in I}$ 满足
$\|e_i-e_j\|_p=2^{1/p}$。任意可数子集都不可能逼近全部这些向量，故 $\ell^p(I)$ 不可分。

现在进入对偶表示。核心步骤是让泛函作用在示性函数上，从而产生一个测度；RN 定理再把这个测度
还原成函数。$\sigma$-有限分块同时保证示性函数属于 $L^p$，并避免对无限质量集合直接操作。

\begin{theorem}[$L^p$ 对偶]\label{theorem:3-3-3-3}
设 $(X,\mathcal A,\mu)$ 为 $\sigma$-有限测度空间，$1\le p<\infty$，$q$ 为共轭指数。
则每个 $\Lambda\in(L^p)^*$ 都唯一表示为
\[
  \Lambda(f)=\int_Xf\overline g\,\dd\mu,
  \qquad g\in L^q,
\]
而且 $\|\Lambda\|=\|g\|_q$。
\end{theorem}

\begin{proof}
取两两不交的可测分块 $X=\bigsqcup_{j=1}^\infty X_j$，使
$\mu(X_j)<\infty$。对 $E\subset X_j$ 定义
\[
  \nu_j(E)=\Lambda(\1_E).
\]
若 $(E_n)$ 是 $X_j$ 中两两不交的可测集，令 $E=\bigcup_nE_n$，则
\[
 \left\|\1_E-\sum_{n=1}^N\1_{E_n}\right\|_p^p
 =\mu\left(E\setminus\bigcup_{n=1}^NE_n\right)\longrightarrow0
\]
（$p=1$ 也包含在此公式中）。由 $\Lambda$ 的连续性，
$\nu_j(E)=\sum_n\nu_j(E_n)$。还要核对全变差有限。若 $(A_\ell)_{\ell=1}^r$ 是 $X_j$ 的有限可测分割，
对每个 $\ell$ 选取 $|\zeta_\ell|=1$ 使
$\zeta_\ell\nu_j(A_\ell)=|\nu_j(A_\ell)|$ （零值时任取）。则
\[
 \sum_{\ell=1}^r|\nu_j(A_\ell)|
 =\left|\Lambda\!\left(\sum_{\ell=1}^r\zeta_\ell\1_{A_\ell}\right)\right|
 \le \|\Lambda\|\,\mu(X_j)^{1/p}.
\]
对分割取上确界，得 $|\nu_j|(X_j)<\infty$。所以 $\nu_j$ 是有限复测度；
实标量情形取 $\zeta_\ell\in\{-1,1\}$，得到有限有符号测度。
若 $\mu(E)=0$，则 $\1_E$ 在 $L^p$ 中为零，故 $\nu_j(E)=0$，即 $\nu_j\ll\mu|_{X_j}$。
而且对每个可测 $A\subset E$ 同样有 $\nu_j(A)=0$；全变差的分割定义于是给出
$|\nu_j|(E)=0$，所以实部、虚部的 Jordan 正负部分也都关于 $\mu|_{X_j}$ 绝对连续。
把 $\Re\nu_j$、$\Im\nu_j$ 分成 Jordan 正负部并分别应用 RN 定理，得到
$h_j\in L^1(\mu|_{X_j};\C)$，使
\[
  \nu_j(E)=\int_Eh_j\,\dd\mu.
\]
实情形的 $h_j$ 为实值。令 $h=h_j$ 于 $X_j$。于是对只跨有限多个分块的简单函数 $s$，
\begin{equation}\label{eq:3-3-3-simple-representation}
  \Lambda(s)=\int_Xsh\,\dd\mu.
\end{equation}
在单个有限测度块上，任何有界可测函数都可由一致有界的简单函数逐点逼近；左侧在 $L^p$ 中收敛，
右侧由 $|h_j|$ 支配收敛。因此 \eqref{eq:3-3-3-simple-representation} 也适用于有限多个块上
有界支撑的可测函数。

设 $1<p<\infty$。记 $C=\|\Lambda\|$、$F_N=\bigcup_{j\le N}X_j$，并令
\[
 A_{N,r}=F_N\cap\{|h|\le r\},
 \qquad
 \phi_{N,r}=\overline{\sgn h}\,|h|^{q-1}\1_{A_{N,r}},
\]
其中 $\sgn h=h/|h|$ 于 $h\ne0$，零处取 $0$。函数 $\phi_{N,r}$ 有界并支撑于有限测度集合，
故可代入上式。置 $I_{N,r}=\int_{A_{N,r}}|h|^q\,\dd\mu<\infty$。因为
$(q-1)p=q$，
\[
 I_{N,r}=\Lambda(\phi_{N,r})
 \le C\|\phi_{N,r}\|_p
 =C I_{N,r}^{1/p}.
\]
若 $I_{N,r}>0$，除以 $I_{N,r}^{1/p}$ 得 $I_{N,r}^{1/q}\le C$；若它为零，同一结论成立。
令 $N,r\to\infty$ 并用单调收敛定理，得到
\[
  \int_X|h|^q\,\dd\mu\le C^q.
\]
故 $h\in L^q$ 且 $\|h\|_q\le C$。

设 $p=1$。固定 $j$ 与 $\varepsilon>0$，令
$A=X_j\cap\{|h|>C+\varepsilon\}$。函数
$\phi=\overline{\sgn h}\1_A$ 属于 $L^1$，且
\[
  \int_A|h|\,\dd\mu
  =\Lambda(\phi)
  \le C\|\phi\|_1=C\mu(A).
\]
若 $\mu(A)>0$，则
\[
 A=\bigcup_{n=1}^{\infty}
   \bigl\{x\in A:|h(x)|\ge C+\varepsilon+1/n\bigr\};
\]
因而其中某个集合 $A_n$ 有正测度。这就给出
\[
 \int_A|h|\,\dd\mu
 \ge(C+\varepsilon)\mu(A)+n^{-1}\mu(A_n)
 >(C+\varepsilon)\mu(A),
\]
与上式矛盾。因此
$|h|\le C+\varepsilon$ 于 $X_j$ 上 a.e. 成立。对 $j$ 取可数并、再令
$\varepsilon\downarrow0$，得到 $h\in L^\infty$ 且 $\|h\|_\infty\le C$。

两种情形中都令 $g=\overline h$；这是第一变量线性约定所要求的共轭。有限测度支撑简单函数在
$L^p$ 中稠密，故 \eqref{eq:3-3-3-simple-representation} 与连续性给出
\[
  \Lambda(f)=\int f h\,\dd\mu=\int f\overline g\,\dd\mu
  \qquad(f\in L^p).
\]
上面的截断测试已经给出 $\|g\|_q\le\|\Lambda\|$；\Holder{} 不等式给出反向估计
$\|\Lambda\|\le\|g\|_q$，故范数相等。

若 $g_1,g_2$ 给出同一泛函，则 $h=\overline{g_1-g_2}$ 满足
$\int_Eh\,\dd\mu=0$ 对每个包含于某个 $X_j$ 的可测 $E$ 成立。分别取
$E=\{\Re h>1/n\}\cap X_j$、$\{\Re h<-1/n\}\cap X_j$，以及虚部对应的集合，可得
$h=0$ a.e.，故 $g_1$ 与 $g_2$ 作为 $L^q$ 元素相等。
\end{proof}

这里的满射性由 RN 定理给出；Hahn--Banach 定理在一般范数空间中的作用，则是保证有足够多的
连续线性泛函。定理的 $p=1$ 端点确实给出 $(L^1)^*=L^\infty$；但
例~\ref{ex:3-3-3-3} 已经说明把箭头反过来写成 $(L^\infty)^*=L^1$ 会失败。对
$1<p<\infty$，在比 $\sigma$-有限更一般的测度空间上还可通过半有限化处理，或在局部化
（localizable）测度空间上陈述相应版本；这里采用上述无需隐藏附加条件的统一表述。

\begin{proposition}[稠密测试类上的收敛升级]\label{proposition:3-3-3-4}
设 $D\subset L^p$ 是稠密线性子空间，$(\Lambda_n)\subset(L^p)^*$ 满足
$\sup_n\|\Lambda_n\|\le C<\infty$，并且对每个 $d\in D$，数列
$\Lambda_n(d)$ 收敛。则其在 $D$ 上的极限唯一延拓为某个
$\Lambda\in(L^p)^*$，且
$\Lambda_n(f)\to\Lambda(f)$ 对每个 $f\in L^p$ 成立。
\end{proposition}

\begin{proof}
在 $D$ 上定义 $\Lambda_0(d)=\lim_n\Lambda_n(d)$。极限保持线性，并且
\[
 |\Lambda_0(d)|\le C\|d\|_p.
\]
若 $f\in L^p$，取 $d_k\in D$ 使 $d_k\to f$。上式说明
$(\Lambda_0(d_k))$ 是 Cauchy 列，其极限与逼近列无关；定义该极限为 $\Lambda(f)$，便得到
$\Lambda\in(L^p)^*$ 且 $\|\Lambda\|\le C$。

给定 $f$ 与 $\varepsilon>0$。若 $C=0$，则每个 $\Lambda_n$ 的范数为零，因而
$\Lambda_n=0=\Lambda$；以下设 $C>0$。取 $d\in D$ 使
$\|f-d\|_p<\varepsilon/(3C)$，再取 $n$ 足够大使
$|\Lambda_n(d)-\Lambda(d)|<\varepsilon/3$。于是
\[
 \begin{aligned}
 |\Lambda_n(f)-\Lambda(f)|
 &\le|\Lambda_n(f-d)|+|\Lambda_n(d)-\Lambda(d)|+|\Lambda(d-f)|\\
 &\le C\|f-d\|_p+\frac\varepsilon3+C\|f-d\|_p<\varepsilon.
 \end{aligned}
\]
范数的一致有界性正是同时控制首尾两项所必需的条件。
\end{proof}

\begin{figure}[htbp]
  \centering
  \begin{tikzpicture}[node distance=13mm and 32mm]
    \node[book node] (simple) {有限测度支撑\\简单函数};
    \node[book strong node,below=of simple] (cc) {$C_c(X)$};
    \node[book warm node,below=of cc] (smooth) {$C_c^\infty(\R^n)$};
    \node[book node,right=of cc] (lp) {$L^p$};
    \node[book warning node,below=of lp] (dual) {$\Lambda\in(L^p)^*$};
    \node[book strong node,below=of dual] (density) {$g\in L^q$};
    \draw[book accent arrow] (simple)--node[left,book label,align=center] {Radon\\正则性} (cc);
    \draw[book dashed arrow] (cc)--node[left,book label,align=center] {mollifier\\第~4.3.1 小节} (smooth);
    \draw[book accent arrow] (simple.east)-- (lp.west);
    \draw[book accent arrow] (cc.east)-- (lp.west);
    \draw[book dashed arrow] (lp)--node[right,book label] {$f\mapsto\int f\overline g$} (dual);
    \draw[book arrow] (dual)--node[right,book label] {RN} (density);
  \end{tikzpicture}
  \caption{左侧是逐层加强的测试类；虚线所示的光滑层要等卷积建立后才可使用。右侧中，连续线性
  观测先在简单函数上产生测度，再由 RN 定理识别为密度。}
  \label{fig:3-3-3-density-duality}
\end{figure}

对偶语言可用来定义弱收敛与弱星收敛：前者测试所有连续线性泛函，后者测试一个指定的前对偶。
命题~\ref{proposition:3-3-3-4} 给出一条常用原则：先在稠密类上验证收敛，再借助一致范数界延伸到
整个空间。条件期望算子、分布与谱测度的逼近论证都会采用这一原则。

\begin{exercise}[简单函数稠密]\label{exercise:3-3-3-1}
给出另一种证明：先把 $f$ 限制到
$E_n=\{1/n\le |f|\le n\}$，再对实值函数的正负部分分别作二进制层逼近；复值情形还要分别处理
实部与虚部。证明 $\mu(E_n)<\infty$，并证明所得有限值简单函数具有有限测度支撑且在 $L^p$ 中收敛到
$f$。
\end{exercise}

\begin{exercise}[$C_c$ 逼近的误差]\label{exercise:3-3-3-2}
设 $s=\sum_{k=1}^Na_k\1_{E_k}$，其中 $E_k$ 两两不交且测度有限。用 Radon 正则性分别构造
$\phi_k\in C_c$，并给出保证
$\|s-\sum_ka_k\phi_k\|_p<\varepsilon$ 的明确误差分配。
\end{exercise}

\begin{exercise}[$L^1$ 对偶的端点测试]\label{exercise:3-3-3-3}
在 $p=1$ 的对偶证明中，把集合
$\{|h|>C+\varepsilon\}$ 分别截到有限测度块上是必要的。写出若直接取整个集合时
$\overline{\sgn h}\1_A$ 可能不属于 $L^1$ 的例子，并完成分块论证。
\end{exercise}

\begin{exercise}[一致有界不可删]\label{exercise:3-3-3-4}
在 $\ell^1$ 中令 $D$ 为有限支撑序列，并定义
$\Lambda_n(x)=n x_n$。证明 $\Lambda_n(d)\to0$ 对每个 $d\in D$ 成立，但存在
$x\in\ell^1$ 使 $(\Lambda_n(x))$ 不收敛。指出命题~\ref{proposition:3-3-3-4} 的哪一项假设失败。
\end{exercise}

范数收敛给出强控制，却不是分析中唯一的极限模式。下一节将比较 a.e. 收敛、依测度收敛与
$L^p$ 收敛，并找出把依测度收敛升级为积分收敛所需的额外条件。

\section{收敛模式与向量值积分}\label{sec:03-04}

前三节把标量积分从非负简单函数一直推进到 $L^p$ 空间与对偶表示。现在出现两道新问题。第一道来自
极限：坏点集合即使越来越小，也可能承载越来越高的尖峰；逐点收敛、依测度收敛和 $L^1$ 收敛因而
不能相互替代。第二道来自值域：若被积函数取值于无限维 Banach 空间，逐坐标可测并不自动产生一个
可以由简单函数逼近的向量场。本节先找出防止质量集中的条件，再建立向量值积分；最后把标量
Fubini 的截面论证升级到 Bochner 积分。


\subsection{依测度收敛、一致可积与 Vitali 定理}\label{subsec:3-4-1}
\index{依测度收敛}\index{一致可积}\index{Vitali 收敛定理}

在有限测度空间上，逐点误差、坏点集合的大小与平均误差分别对应三种不同的收敛语言。最短的反例
已经把它们之间的裂缝全部暴露出来。

Vitali 定理在概率极限定理中往往比 DCT 灵活：它不要求寻找一支同时支配整列的函数，而是直接
禁止一族函数把质量压进越来越小的集合。这个条件也成为概率论中控制积分极限的基本紧性条件。

\begin{example}[缩窄而升高的尖峰]\label{ex:3-4-1-1}
在 $([0,1],\Borel([0,1]),m)$ 上令
\[
  f_n(x)=n\1_{(0,1/n)}(x).
\]
对每个 $\varepsilon>0$，当 $n>\varepsilon$ 时
\[
 m\{|f_n|>\varepsilon\}=\frac1n\longrightarrow0,
\]
但 $\|f_n\|_1=1$。而且给定任意 $K>0$，只要取整数 $n>K$，便有
\[
 \int_{\{|f_n|>K\}}|f_n|\,\dd m=1.
\]
因此坏集虽在缩小，全部质量却集中在其中。仅仅知道坏集的测度趋于零，还不足以控制积分。
\end{example}

\begin{definition}[依测度收敛]\label{def:3-4-1-1}
设 $f_n,f$ 是有限测度空间 $(X,\mathcal A,\mu)$ 上 a.e. 取有限值的可测函数。若对每个
$\varepsilon>0$，
\[
  \mu\{x:|f_n(x)-f(x)|>\varepsilon\}\longrightarrow0,
\]
则称 $f_n$ \emph{依测度收敛}于 $f$，记作 $f_n\xrightarrow{\mu}f$。
\end{definition}

当 $\mu(X)=1$ 时，同一概念在概率论中称为\emph{依概率收敛}（convergence in probability）。

\begin{definition}[一致可积]\label{def:3-4-1-2}
设 $\mu(X)<\infty$。族 $\mathcal F\subset L^1(\mu)$ 称为\emph{一致可积}，若
\[
  \lim_{K\to\infty}\sup_{f\in\mathcal F}
  \int_{\{|f|>K\}}|f|\,\dd\mu=0.
  \tag{3.4.1}\label{eq:ui-tail-definition}
\]
\end{definition}

这个定义既控制高度，也控制小集合上的质量。下面给出以后反复使用的等价形式。

\begin{lemma}[一致可积的两种口径]\label{lemma:3-4-1-ui-equivalence}
在有限测度空间上，\eqref{eq:ui-tail-definition} 等价于下列两项同时成立：
\begin{enumerate}
  \item $\sup_{f\in\mathcal F}\|f\|_1<\infty$；
  \item 对每个 $\varepsilon>0$，存在 $\delta>0$，使得 $\mu(E)<\delta$ 时
  \[
    \sup_{f\in\mathcal F}\int_E|f|\,\dd\mu<\varepsilon.
  \]
\end{enumerate}
\end{lemma}

\begin{proof}
先假设 \eqref{eq:ui-tail-definition}。取 $K_0$ 使其后的统一尾积分小于 $1$，则
\[
  \|f\|_1
  \le K_0\mu(X)+\int_{\{|f|>K_0\}}|f|\,\dd\mu
  \le K_0\mu(X)+1,
\]
所以第一项成立。给定 $\varepsilon>0$，再取 $K$ 使统一尾积分小于 $\varepsilon/2$，并令
$\delta=\varepsilon/(2K)$。当 $\mu(E)<\delta$ 时，
\[
 \int_E|f|
 \le \int_{E\cap\{|f|\le K\}}|f|+
      \int_{\{|f|>K\}}|f|
 <K\delta+\frac\varepsilon2=\varepsilon.
\]

反过来，设两项成立，并记 $M=\sup_{f\in\mathcal F}\|f\|_1$。给定 $\varepsilon>0$，由第二项
取相应的 $\delta$。对任意 $K>\max\{1,M/\delta\}$，点态不等式
$\1_{\{|f|>K\}}\le |f|/K$ 给出
\[
  \mu\{|f|>K\}\le \frac{\|f\|_1}{K}<\delta
\]
对所有 $f\in\mathcal F$ 同时成立，因而相应尾积分小于 $\varepsilon$。这就是
\eqref{eq:ui-tail-definition}。
\end{proof}

一致 $L^1$ 有界只给出上面第一项，\cref{ex:3-4-1-1} 正说明第二项不能省略。一个较高次矩的
一致控制则足以排除尖峰。

\begin{example}[$L^p$ 有界族一致可积]\label{ex:3-4-1-2}
设 $\mu(X)<\infty$、$p>1$、$\mathcal F\subset L^p(\mu)$，且
$C:=\sup_{f\in\mathcal F}\|f\|_p<\infty$。有限测度嵌入先给出
$\mathcal F\subset L^1(\mu)$。在 $\{|f|>K\}$ 上有
$|f|\le K^{1-p}|f|^p$，故
\[
 \sup_{f\in\mathcal F}\int_{\{|f|>K\}}|f|\,\dd\mu
 \le K^{1-p}C^p\longrightarrow0.
\]
因此 $\mathcal F$ 一致可积。概率空间只是这个计算中最常见的情形。
\end{example}

依测度收敛也不能保证整列在每一点安定下来。

\begin{example}[打字机序列]\label{ex:3-4-1-3}
对 $m\ge0$ 与 $0\le k<2^m$ 令
\[
 I_{m,k}=[k2^{-m},(k+1)2^{-m}),\qquad f_{2^m+k}=\1_{I_{m,k}}
\]
（在端点 $1$ 处统一置零）。若第 $n$ 项位于第 $m$ 层，则
$m\{|f_n|>1/2\}=2^{-m}$；当 $n\to\infty$ 时 $m\to\infty$，所以 $f_n\xrightarrow{m}0$。
可是每个 $x\in[0,1)$ 在每一层恰落入一个 $I_{m,k}$。沿整列观察，函数值在每层出现一次 $1$，
其余各项为 $0$，故整列在任何 $x\in[0,1)$ 都不收敛。依测度收敛的正确逐点结论只能落在子列上。
\end{example}

\begin{definition}[依测度 Cauchy]\label{def:3-4-1-3}
若对每个 $\varepsilon>0$，
\[
 \mu\{|f_n-f_m|>\varepsilon\}\longrightarrow0
 \qquad(m,n\to\infty),
\]
则称 $(f_n)$ 是\emph{依测度 Cauchy 列}。
\end{definition}

先记录逐点收敛与依测度收敛之间的一条有限测度结论。有限性在支配函数 $1$ 可积这一处使用。

\begin{proposition}[几乎处处收敛蕴含依测度收敛]\label{proposition:3-4-1-ae-measure}
若 $\mu(X)<\infty$ 且 $f_n\to f$ a.e.，则 $f_n\xrightarrow{\mu}f$。
\end{proposition}

\begin{proof}
固定 $\varepsilon>0$。指标函数
$\1_{\{|f_n-f|>\varepsilon\}}$ 几乎处处趋于零，并由可积函数 $1$ 支配。标量 DCT 给
\[
 \mu\{|f_n-f|>\varepsilon\}
 =\int_X\1_{\{|f_n-f|>\varepsilon\}}\,\dd\mu\longrightarrow0.
\]
\end{proof}

\begin{proposition}[子列原理]\label{proposition:3-4-1-1}
设 $\mu(X)<\infty$。则 $f_n\xrightarrow{\mu}f$，当且仅当 $(f_n)$ 的每个子列都有一个进一步
子列几乎处处收敛于 $f$。
\end{proposition}

\begin{proof}
先设 $f_n\xrightarrow{\mu}f$，并从任意子列中仍以 $(f_n)$ 记之。递归选取 $n_k$，使
\[
  \mu(E_k)<2^{-k},\qquad
  E_k:=\{|f_{n_k}-f|>2^{-k}\}.
\]
对每个 $N$，次可加性给
\[
 \mu\!\left(\bigcup_{k\ge N}E_k\right)
 \le\sum_{k\ge N}2^{-k}=2^{1-N}.
\]
集合 $\bigcap_N\bigcup_{k\ge N}E_k$ 是递减集合
$\bigcup_{k\ge N}E_k$ 的交；第一个集合测度有限，故测度的上连续性给其测度为零。对交集外的
$x$，存在 $N(x)$ 使 $k\ge N(x)$ 时 $x\notin E_k$，于是
$|f_{n_k}(x)-f(x)|\le2^{-k}\to0$。

反之，若原列不依测度收敛，则存在 $\varepsilon,\delta>0$ 及一个子列 $(f_{n_k})$，使
\[
  \mu\{|f_{n_k}-f|>\varepsilon\}\ge\delta\qquad(k\ge1).
\]
按照假设，该子列有进一步子列 $(f_{n_{k_j}})$ 几乎处处收敛于 $f$。由
\cref{proposition:3-4-1-ae-measure}，它又依测度收敛于 $f$，所以上述测度应趋于零，与其不小于
$\delta$ 矛盾。
\end{proof}

现在可以精确填补尖峰反例留下的缺口。

\begin{theorem}[Vitali 收敛定理]\label{theorem:3-4-1-2}
设 $\mu(X)<\infty$。若 $f_n\xrightarrow{\mu}f$ 且 $\{f_n:n\ge1\}$ 一致可积，则
$f\in L^1(\mu)$ 且
\[
  \|f_n-f\|_1\longrightarrow0.
\]
反之，$L^1$ 收敛蕴含依测度收敛，并且 $\{f_n:n\ge1\}$ 一致可积。
\end{theorem}

\begin{proof}
先证正向。一致可积给出共同的 $L^1$ 上界：由
\cref{lemma:3-4-1-ui-equivalence}，存在 $M<\infty$ 使 $\|f_n\|_1\le M$。子列原理允许从
原列抽取 $f_{n_j}\to f$ a.e. 的子列。Fatou 引理给
\[
  \int_X|f|\,\dd\mu
  \le\liminf_{j\to\infty}\int_X|f_{n_j}|\,\dd\mu\le M,
\]
故 $f\in L^1$。

还需把统一尾控制传给这个新出现的 $f$，这里不能预先调用待证的 $L^1$ 收敛。固定 $K>0$。在
$f_{n_j}(x)\to f(x)$ 的点上，若 $|f(x)|>K$，则最终有 $|f_{n_j}(x)|>K$，因而
\[
 |f(x)|\1_{\{|f(x)|>K\}}
 \le\liminf_{j\to\infty}
 |f_{n_j}(x)|\1_{\{|f_{n_j}(x)|>K\}}.
\]
再次使用 Fatou 引理，得到
\[
 \int_{\{|f|>K\}}|f|\,\dd\mu
 \le\liminf_j\int_{\{|f_{n_j}|>K\}}|f_{n_j}|\,\dd\mu
 \le\sup_n\int_{\{|f_n|>K\}}|f_n|\,\dd\mu.
 \tag{3.4.2}\label{eq:vitali-limit-tail}
\]

给定 $\varepsilon>0$，由一致可积及 \eqref{eq:vitali-limit-tail} 取 $K$，使 $f_n$ 的统一尾积分
以及 $f$ 的尾积分都小于 $\varepsilon/4$。令
\[
 h_n=|f_n-f|\wedge 2K.
\]
$h_n\xrightarrow{\mu}0$。对任意 $0<\eta<2K$，
\[
 \int_Xh_n\,\dd\mu
 \le \eta\mu(X)+2K\mu\{h_n>\eta\}.
 \tag{3.4.3}\label{eq:bounded-in-measure-l1}
\]
先取 $\eta$ 使第一项小于 $\varepsilon/4$，再令 $n$ 足够大使第二项小于 $\varepsilon/4$，便得
$\int h_n<\varepsilon/2$。逐点有
\[
 |f_n-f|
 \le |f_n|\1_{\{|f_n|>K\}}
    +|f|\1_{\{|f|>K\}}+h_n;
\]
这个不等式可逐情形核对：当两端都不超过 $K$ 时 $h_n=|f_n-f|$；当两端都超过 $K$ 时两个尾项
之和不小于 $|f_n-f|$；若只有一端超过 $K$，则 $|f_n-f|\le2K$ 时 $h_n$ 已等于差值，而
$|f_n-f|>2K$ 时，相应尾项加 $2K$ 不小于两端绝对值之和，因而不小于差值。积分后令
$n\to\infty$，上面三项之和小于 $\varepsilon$，得到
$\|f_n-f\|_1\to0$。

现在设 $\|f_n-f\|_1\to0$。由
$\1_{\{|f_n-f|>\varepsilon\}}\le|f_n-f|/\varepsilon$ 直接积分，得到
\[
 \mu\{|f_n-f|>\varepsilon\}
 \le\varepsilon^{-1}\|f_n-f\|_1\longrightarrow0,
\]
所以依测度收敛。为证一致可积，注意在集合 $A_n(K)=\{|f_n|>K\}$ 上，若
$|f|\le K/2$，则 $|f_n-f|>K/2$，从而 $|f_n|\le2|f_n-f|$；若 $|f|>K/2$，则
$|f_n|\le|f_n-f|+|f|$。故
\[
 \int_{A_n(K)}|f_n|
 \le2\|f_n-f\|_1+
      \int_{\{|f|>K/2\}}|f|.
 \tag{3.4.4}\label{eq:l1-convergence-ui-tail}
\]
先取 $N$ 使 $n\ge N$ 时第一项一致很小，再令 $K\to\infty$ 控制第二项。有限多个
$f_1,\dots,f_{N-1}$ 各自的尾积分也随 $K\to\infty$ 趋零，故整族一致可积。
\end{proof}

图 \ref{fig:3-4-1-vitali-split} 展示证明中的三项分工：式
\eqref{eq:bounded-in-measure-l1} 把截断后的主体积分压到零，一致可积则同时控制两端的高度尾部。

\begin{figure}[htbp]
\centering
\begin{tikzpicture}[x=1cm,y=1cm]
  \draw[->,BookInk!70] (-.2,0)--(9.4,0) node[right,book label] {$x$};
  \draw[->,BookInk!70] (0,-.2)--(0,3.6) node[above,book label] {误差大小};
  \draw[BookWarning,very thick] plot[smooth] coordinates
    {(.3,.25) (1,.45) (1.55,3.1) (2.05,.6) (3,.3) (4.2,.55) (5.1,.35)
     (6,.65) (6.6,2.75) (7.1,.45) (8.8,.25)};
  \draw[BookWarm,dashed] (0,1.7)--(9,1.7) node[right,book label] {$2K$};
  \fill[BookWarning!18] (1.27,0) rectangle (1.83,3.08);
  \fill[BookWarning!18] (6.35,0) rectangle (6.87,2.72);
  \node[book warning node] at (4.1,3.1) {高尾部\\由 UI 控制};
  \draw[book arrow] (3.15,2.82)--(1.72,2.42);
  \draw[book arrow] (5.05,2.82)--(6.55,2.25);
  \node[book strong node] at (4.25,1.05) {截断主体 $h_n$\\坏集测度趋零};
  \node[book warm node] at (7.8,.72) {低于 $\eta$ 的部分\\至多 $\eta\mu(X)$};
\end{tikzpicture}
\caption{Vitali 证明的三块误差：高度尾部、截断后的坏集，以及阈值以下的均匀小量。}
\label{fig:3-4-1-vitali-split}
\end{figure}

一致可积还可以由一个超线性代价函数检测；这个判据经常比直接估计尾积分方便。

\begin{proposition}[de la \ValleePoussin{} 判据]\label{proposition:3-4-1-3}
设 $\mu(X)<\infty$ 且 $\mathcal F\subset L^1(\mu)$。若存在凸的单调递增函数
$\Phi:[0,\infty)\to[0,\infty)$，满足 $\Phi(0)=0$、
\[
 \frac{\Phi(t)}t\longrightarrow\infty
 \quad(t\to\infty),
 \qquad
 \sup_{f\in\mathcal F}\int_X\Phi(|f|)\,\dd\mu<\infty,
 \tag{3.4.5}\label{eq:dlvp-assumption}
\]
则 $\mathcal F$ 一致可积。反过来，在有限测度空间上（特别在概率空间上），任一一致可积族都存在满足
\eqref{eq:dlvp-assumption} 的 $\Phi$。
\end{proposition}

\begin{proof}
凸性与 $\Phi(0)=0$ 蕴含 $t\mapsto\Phi(t)/t$ 在 $(0,\infty)$ 上单调不减：若
$0<s<t$，则 $s=(s/t)t+(1-s/t)0$，所以
$\Phi(s)\le(s/t)\Phi(t)$。对充分大的 $K$ 有 $\Phi(K)>0$，并且在 $|f|>K$ 上有
\[
  |f|\le \frac{K}{\Phi(K)}\Phi(|f|).
\]
若 \eqref{eq:dlvp-assumption} 中积分的共同上界为 $C$，则
\[
 \sup_{f\in\mathcal F}\int_{\{|f|>K\}}|f|
 \le C\frac{K}{\Phi(K)}\longrightarrow0.
\]

反向设 $\mu(X)<\infty$ 且 $\mathcal F$ 一致可积。递归选取阈值 $K_j$，使
\[
 K_j\ge j,\qquad K_{j+1}>K_j,
 \qquad
 \sup_{f\in\mathcal F}\int_{\{|f|>K_j\}}|f|\,\dd\mu\le2^{-j}
 \qquad(j\ge1).
\]
定义
\[
 \Phi(t)=\sum_{j=1}^{\infty}(t-K_j)_+,
 \qquad (u)_+=\max\{u,0\}.
 \tag{3.4.6}\label{eq:dlvp-construction}
\]
因为 $K_j\ge j$，对每个固定 $t$，和中只有有限个非零项；每个有限部分和都凸且递增，且对固定
$t$ 最终等于 $\Phi(t)$，所以 $\Phi$ 凸、递增并满足 $\Phi(0)=0$。给定 $L>0$，取整数
$M>2L$。当 $t\ge2K_M$ 时，对 $1\le j\le M$ 有
$(t-K_j)_+\ge t/2$，故 $\Phi(t)/t\ge M/2>L$。这按定义证明
$\Phi(t)/t\to\infty$。最后由
Tonelli 定理及 $(|f|-K_j)_+\le |f|\1_{\{|f|>K_j\}}$，
\[
 \int_X\Phi(|f|)\,\dd\mu
 \le\sum_{j=1}^{\infty}
      \int_{\{|f|>K_j\}}|f|\,\dd\mu
 \le\sum_{j=1}^{\infty}2^{-j}=1,
\]
且上界与 $f$ 无关。
\end{proof}

依测度收敛本身也有完备性。其证明再次把一列坏集测度估计压缩成一个几乎处处命题。

\begin{proposition}[依测度完备性]\label{proposition:3-4-1-4}
在有限测度空间中，每个依测度 Cauchy 列都有依测度极限；极限在 a.e. 等价意义下唯一。
\end{proposition}

\begin{proof}
递归选取 $n_1<n_2<\cdots$，使
\[
 \mu\{|f_{n_{k+1}}-f_{n_k}|>2^{-k}\}<2^{-k}.
 \tag{3.4.7}\label{eq:measure-cauchy-fast-subsequence}
\]
与子列原理证明相同，尾并的测度估计说明：除去零集后，存在 $k_0(x)$，使 $k\ge k_0(x)$ 时
\[
 |f_{n_{k+1}}(x)-f_{n_k}(x)|\le2^{-k}.
\]
因此向量不是问题，此处只是标量级数：
$\sum_k|f_{n_{k+1}}(x)-f_{n_k}(x)|<\infty$，所以 $(f_{n_k}(x))$ 收敛。把极限记为 $f(x)$，
并在例外零集上置零。它是可测函数，且由
\cref{proposition:3-4-1-ae-measure}，$f_{n_k}\xrightarrow{\mu}f$。

给定 $\varepsilon,\delta>0$。先由原列的 Cauchy 性取 $N$，使 $m,n\ge N$ 时
$\mu\{|f_n-f_m|>\varepsilon/2\}<\delta/2$；再取 $k$ 足够大，使 $n_k\ge N$ 且
$\mu\{|f_{n_k}-f|>\varepsilon/2\}<\delta/2$。于是对所有 $n\ge N$，
\[
 \mu\{|f_n-f|>\varepsilon\}
 \le\mu\{|f_n-f_{n_k}|>\varepsilon/2\}
   +\mu\{|f_{n_k}-f|>\varepsilon/2\}<\delta.
\]
故整列依测度收敛于 $f$。若同时依测度收敛于 $g$，则
\[
 \mu\{|f-g|>\varepsilon\}
 \le\mu\{|f-f_n|>\varepsilon/2\}
   +\mu\{|f_n-g|>\varepsilon/2\}\longrightarrow0,
\]
所以 $f=g$ a.e.
\end{proof}

\begin{remark}[无限空间上要把尖峰与逃逸分开]
若 $\mu(X)=\infty$，a.e. 收敛未必蕴含这里定义的全局依测度收敛。例如
$\1_{[n,n+1]}$ 在 $\R$ 上逐点趋零，但每一项的坏集测度均为 $1$。对一个有限测度耗尽
$E_m\uparrow X$，可以改谈每个 $E_m$ 上的局部依测度收敛；然而局部收敛加高度尾控制仍不足以推出
全局 $L^1$ 收敛，因为质量可能逃向无穷远。还需另加空间紧性条件
\[
 \lim_{m\to\infty}\sup_n\int_{X\setminus E_m}|f_n|\,\dd\mu=0.
\]
$n\1_{(0,1/n)}$ 违反高度尾控制，$\1_{[n,n+1]}$ 则违反空间紧性；两种失效机制需要分别控制。
\end{remark}

Vitali 定理可直接用于幂函数的积分收敛。例如在总质量为一的测度空间上，若 $r>0$、可测函数
$X_n\xrightarrow{\mu}X$，且
$\{|X_n|^r:n\ge1\}$ 一致可积，则子列原理表明每个子列都有 a.e. 收敛于 $X$ 的进一步子列；
在该子列上取连续函数 $t\mapsto|t|^r$，再用子列原理的反向，得到
$|X_n|^r\xrightarrow{\mu}|X|^r$，因而
$\int|X_n|^r\to\int|X|^r$。在概率记号中，这给出依概率收敛加一致可积所蕴含的矩收敛。

\begin{exercise}[尖峰与一致 $L^1$ 界]\label{exercise:3-4-1-1}
构造一个非负函数族，使其 $L^1$ 范数一致有界而不一致可积；再用
\cref{lemma:3-4-1-ui-equivalence} 指出失败的是哪一项。
\end{exercise}
\begin{exercise}[$L^p$ 控制]\label{exercise:3-4-1-2}
设 $p>1$。分别用尾积分估计与小测集绝对连续性证明 $L^p$ 有界族一致可积。
\end{exercise}
\begin{exercise}[打字机子列]\label{exercise:3-4-1-3}
从 \cref{ex:3-4-1-3} 中每层恰取一项。证明所得任意这样的子列都几乎处处趋于零，
并特别写出选取每层第一项时的点态极限。
\end{exercise}
\begin{exercise}[局部 Vitali]\label{exercise:3-4-1-4}
设 $E_m\uparrow X$、$\mu(E_m)<\infty$。在局部依测度收敛、统一高度尾控制及上述空间紧性条件下，
证明 $L^1$ 收敛的对应版本。
\end{exercise}

标量函数可在每个点直接比较大小；向量值函数没有正负部，甚至“每个连续线性泛函看起来都可测”
也可能遗漏值域的不可分部分。下一小节从这个失败处开始。


\subsection{强可测性、Pettis 定理与 Bochner 积分}\label{subsec:3-4-2}
\index{强可测}\index{弱可测}\index{Pettis 可测性定理}\index{Bochner 积分}

本小节中的底层测度空间取完成化。这样，在零集上修改向量值函数仍保持可测性，也与积分的 a.e.
等价类约定一致。若从未完备的 $\sigma$-代数出发，以下陈述应理解为：存在一个 a.e. 相等的代表，
在完成化上满足相应性质。

Hilbert 与 Schmidt 研究层势积分方程时，常把 Fourier 系数组织成 $\ell^2$ 中的向量和无穷矩阵；
有限截断之外还需用完备性补上极限。这里先处理一个更基础的问题：参数化核截面何时真能作为
Banach 或 Hilbert 值函数积分。

例如，对 Hilbert 空间中的规范正交序列 $(e_n)$，形式级数
$u(x)=\sum_na_n(x)e_n$ 即使每个坐标 $a_n$ 都可测，也还不能直接写 $\int u$：必须先证明
部分和确实在范数中收敛，所得极限能由有限值简单函数逼近，并且
$\int(\sum_n|a_n|^2)^{1/2}<\infty$。这说明“逐坐标可测”与“向量可积”之间还隔着范数收敛、
强可测性和范数可积三道条件。

\begin{definition}[强可测]\label{def:3-4-2-1}
设 $B$ 为 Banach 空间。函数 $f:X\to B$ 称为\emph{强可测}，若存在取有限多个值的可测简单函数
$s_n:X\to B$，使 $s_n(x)\to f(x)$ a.e.（收敛关于 $B$ 的范数）。
\end{definition}

\begin{definition}[弱可测与几乎可分值]\label{def:3-4-2-2}
若对每个 $\phi\in B^*$，标量函数 $\phi\circ f$ 都可测，则称 $f$ \emph{弱可测}。若存在
可分闭子空间 $B_0\subset B$，使 $f(x)\in B_0$ a.e.，则称 $f$ \emph{几乎可分值}。
\end{definition}

这里需要的 Hahn--Banach 工具只有\cref{corollary:3-3-3-norming} 已经证明的一点范数化形式。
可分性把逐点可选的泛函压缩成一个可数族。

\begin{lemma}[可分子空间的可数范数化族]\label{lemma:3-4-2-1}
若 $B_0\subset B$ 为可分闭子空间，则存在 $\phi_j\in B^*$、$\|\phi_j\|\le1$，使
\[
  \|b\|=\sup_{j\ge1}|\phi_j(b)|\qquad(b\in B_0).
  \tag{3.4.8}\label{eq:countable-norming-family}
\]
\end{lemma}

\begin{proof}
若 $B_0=\{0\}$，取零泛函即可。否则在 $B_0$ 的单位球面取稠密列 $(u_j)$。由一点
Hahn--Banach，对每个 $j$ 存在 $\psi_j\in B_0^*$，满足
$\|\psi_j\|=1$ 与 $\psi_j(u_j)=1$；再以相同范数把 $\psi_j$ 延拓到 $B^*$，记为 $\phi_j$。
于是 $|\phi_j(b)|\le\|b\|$，所以 \eqref{eq:countable-norming-family} 右端不超过左端。

若 $b\ne0$，给定 $\varepsilon>0$，选 $j$ 使
$\|u_j-b/\|b\|\|<\varepsilon$。则
\[
 |\phi_j(b)|
 =\|b\|\,|\psi_j(b/\|b\|)|
 \ge\|b\|(1-\varepsilon).
\]
令 $\varepsilon\downarrow0$ 即得反向不等式。
\end{proof}

\begin{theorem}[Pettis 可测性定理]\label{theorem:3-4-2-1}
Banach 值函数 $f:X\to B$ 强可测，当且仅当它弱可测且几乎可分值。
\end{theorem}

\begin{proof}
先设 $f$ 强可测，并取简单函数 $s_n\to f$ a.e.。所有 $s_n$ 的值合在一起是可数集；令 $B_0$
为其闭线性包，则 $B_0$ 可分，而且在收敛成立的点上 $f(x)\in B_0$。对任意 $\phi\in B^*$，
$\phi(s_n)$ 可测且在共同零集外趋于 $\phi(f)$。测度空间完备，故这个 a.e. 极限的任意零集修改仍
可测，于是 $f$ 弱可测。

反之，设 $f$ 弱可测且几乎取值于可分闭子空间 $B_0$。在例外零集上把 $f$ 改为 $0$；完成性
保证修改后的函数仍弱可测。取 $B_0$ 中的稠密列 $(b_k)_{k\ge1}$。由
\cref{lemma:3-4-2-1}，对每个 $k$，
\[
 x\longmapsto\|f(x)-b_k\|
 =\sup_{j\ge1}|\phi_j(f(x)-b_k)|
\]
是可数个可测函数的上确界，故可测。对 $n\ge1$，令 $s_n(x)$ 为
$b_1,\dots,b_n$ 中离 $f(x)$ 最近的点；若有多个最近点，取指标最小者。精确地，值层
\[
 \{s_n=b_k\}
 =\bigcap_{\ell<k}\{\|f-b_k\|<\|f-b_\ell\|\}
  \cap\bigcap_{k<\ell\le n}\{\|f-b_k\|\le\|f-b_\ell\|\}
\]
可测，所以 $s_n$ 是有限值简单函数。又因 $(b_k)$ 稠密，
\[
 \|s_n(x)-f(x)\|=\min_{1\le k\le n}\|b_k-f(x)\|\longrightarrow0
\]
对每个 $x$ 成立。这就证明了强可测性，且整个构造没有调用可测选择定理。
\end{proof}

弱可测缺少可分值条件时，结论确实失败。

\begin{example}[不可分正交族]\label{ex:3-4-2-2}
取 $X=[0,1]$ 配 Lebesgue 测度，$H=\ell^2([0,1])$，并以 $e_t$ 表示第 $t$ 个标准单位向量。
定义 $f(t)=e_t$。对任意 $\phi\in H^*$，令
$c_t=\phi(e_t)$。对任意有限集 $F\subset[0,1]$，在线性组合的系数中选取适当相位，并用
Cauchy--Schwarz，可得
\[
 \left(\sum_{t\in F}|c_t|^2\right)^{1/2}
 =\left|\phi\!\left(
   \sum_{t\in F}\frac{\overline{c_t}}{(\sum_{s\in F}|c_s|^2)^{1/2}}e_t
   \right)\right|
 \le\|\phi\|,
\]
其中分母为零时不需计算。于是对每个 $n$，集合
$\{t:|c_t|\ge1/n\}$ 必为有限集，否则可选任意大的有限子集使左端超过 $\|\phi\|$。因此
$\{t:c_t\ne0\}$ 可数，而
$t\mapsto\phi(f(t))=c_t$ 除可数集外为零，故可测。这对每个 $\phi\in H^*$ 成立，所以 $f$ 弱可测。

另一方面，每个 $h\in H$ 的支撑都可数，因为
$\{t:|h_t|\ge1/n\}$ 对每个 $n$ 必为有限集，而支撑是这些集合的可数并。任一可分子空间
$H_0\subset H$ 因而都包含在某个可数坐标子空间中：取 $H_0$ 的稠密列，把各向量的可数支撑取并
得到可数集 $J$，再取闭包便有 $H_0\subset\ell^2(J)$。于是除可数个
$t\in J$ 外，$e_t\notin H_0$。所以 $f$ 不几乎可分值，依
\cref{theorem:3-4-2-1} 也不强可测。
\end{example}

\begin{example}[连续映射的强可测性]\label{ex:3-4-2-3}
设 $X$ 是可分度量空间，其测度结构包含 Borel 集，$B$ 是 Banach 空间，且 $f:X\to B$ 连续。
取 $X$ 的可数稠密集 $D$。连续性给
$f(X)\subset\overline{f(D)}$，所以 $f$ 几乎可分值；对每个 $\phi\in B^*$，$\phi\circ f$ 连续，
因而 Borel 可测。Pettis 定理遂说明 $f$ 强可测。这一结论以后可用于连续参数的核截面。
\end{example}

现在进入积分。逐个对偶坐标可积并不保证这些坐标积分来自同一个 $B$ 中向量，所以不能把它当定义。

\begin{definition}[Bochner 可积]\label{def:3-4-2-3}
强可测函数 $f:X\to B$ 若满足
\[
  \int_X\|f(x)\|\,\dd\mu(x)<\infty,
\]
则称 $f$ \emph{Bochner 可积}。若
$s=\sum_{k=1}^m b_k\1_{E_k}$ 的 $E_k$ 两两不交，且每个 $b_k\ne0$ 所在的 $E_k$ 都有有限测度，
忽略零值层后，定义
\[
  \int_Xs\,\dd\mu=\sum_{k:\,b_k\ne0} b_k\mu(E_k).
  \tag{3.4.9}\label{eq:simple-bochner-integral}
\]
一般 $f$ 的积分由下面的 $L^1(B)$ 逼近构造。
\end{definition}

\begin{lemma}[简单函数逼近与 Bochner 积分的构造]\label{lemma:3-4-2-bochner-construction}
每个 Bochner 可积函数 $f$ 都存在可积简单函数 $s_n$，使
\[
  \int_X\|s_n-f\|\,\dd\mu\longrightarrow0.
  \tag{3.4.10}\label{eq:bochner-simple-l1-approximation}
\]
而且 $\int s_n$ 在 $B$ 中收敛，其极限与逼近列无关。
\end{lemma}

\begin{proof}
先处理简单函数。若
$s=\sum_i b_i\1_{E_i}=\sum_jc_j\1_{F_j}$ 是两个不交层表示，先删去两边的零值层。
由可积性，余下每个值层都有有限测度。集合
$E_i\cap F_j$ 构成共同细分；在每个非空细分层上有 $b_i=c_j$。因此
\[
 \sum_i b_i\mu(E_i)
 =\sum_{i,j}b_i\mu(E_i\cap F_j)
 =\sum_{i,j}c_j\mu(E_i\cap F_j)
 =\sum_jc_j\mu(F_j),
\]
故
\eqref{eq:simple-bochner-integral} 与表示无关。对 $s-t$ 的值层作同样的共同细分，并在有限向量和上
使用三角不等式，得到
\[
 \left\|\int_X(s-t)\,\dd\mu\right\|
 \le\int_X\|s-t\|\,\dd\mu
 \tag{3.4.11}\label{eq:simple-bochner-estimate}
\]
对任意可积简单函数 $s,t$ 成立。

令 $f$ Bochner 可积。由 Pettis 定理的正向部分，修改一个零集后可设 $f(X)$ 包含于可分闭子空间
$B_0$。在 $B_0$ 中取稠密列 $(b_k)_{k\ge1}$，并令 $b_1=0$。对 $f(x)\ne0$，令 $q_n(x)$ 为首个满足
\[
  \|q_n(x)-f(x)\|\le2^{-n}\|f(x)\|
  \tag{3.4.12}\label{eq:relative-countable-approximation}
\]
的 $b_k$；当 $f(x)=0$ 时令 $q_n(x)=0$。稠密性保证候选点存在。函数
$x\mapsto\|f(x)-b_k\|$ 可测，例如可由强简单逼近取连续范数的点态极限；“首个”规则因此使
$q_n$ 的每个值层可测。更具体地，若
$C_{n,k}=\{\|f-b_k\|\le2^{-n}\|f\|\}$，则首个指标为 $k$ 的值层是
$C_{n,k}\setminus\bigcup_{j<k}C_{n,j}$（再把 $\{f=0\}$ 归入零值层）。于是 $q_n$ 是可数值可测
函数，并且
\[
 \|q_n-f\|\le2^{-n}\|f\|,
 \qquad \|q_n\|\le(1+2^{-n})\|f\|.
\]

把 $A_{n,k}$ 取为“首个合格指标恰为 $k$”的两两不交层，并把 $\{f=0\}$ 归入 $A_{n,1}$；于是
$q_n=\sum_kb_k\1_{A_{n,k}}$。由非负积分的可数可加性，
\[
 \sum_{k\ge1}\|b_k\|\mu(A_{n,k})=\int_X\|q_n\|\,\dd\mu<\infty.
\]
故可选 $m(n)$，使这一级数在 $k>m(n)$ 的尾和小于 $2^{-n}$。令
\[
 s_n=\sum_{k\le m(n)}b_k\1_{A_{n,k}},
\]
并在其余值层上置零。它是真正的有限值可积简单函数，且
\[
 \int\|s_n-f\|
 \le\int\|s_n-q_n\|+\int\|q_n-f\|
 \le2^{-n}+2^{-n}\int\|f\|\longrightarrow0.
\]

由 \eqref{eq:simple-bochner-estimate}，
\[
 \left\|\int s_n-\int s_m\right\|
 \le\int\|s_n-f\|+\int\|s_m-f\|\longrightarrow0.
\]
$B$ 完备，所以 $\int s_n$ 有极限。若 $(t_n)$ 是另一套满足
\eqref{eq:bochner-simple-l1-approximation} 的逼近，则
\[
 \left\|\int s_n-\int t_m\right\|
 \le\int\|s_n-f\|+\int\|t_m-f\|,
\]
右端在 $n,m\to\infty$ 时趋零，因此两套极限相同。
\end{proof}

以该极限定义 $\int f$。简单函数估计与 Banach 空间的完备性还给出以下基本性质。

\begin{theorem}[Bochner 积分基本估计]\label{theorem:3-4-2-2}
Bochner 积分线性。若 $f$ Bochner 可积，则
\[
 \left\|\int_Xf\,\dd\mu\right\|
 \le\int_X\|f\|\,\dd\mu,
 \qquad
 \phi\!\left(\int_Xf\,\dd\mu\right)
 =\int_X\phi(f(x))\,\dd\mu(x)
 \tag{3.4.13}\label{eq:bochner-basic-estimates}
\]
对每个 $\phi\in B^*$ 成立。
\end{theorem}

\begin{proof}
线性先对简单函数由有限和成立。若 $s_n\to f$、$t_n\to g$ 于 $L^1(B)$，则
$as_n+bt_n\to af+bg$ 于 $L^1(B)$；在简单函数等式中取极限便得一般线性。

取 \eqref{eq:bochner-simple-l1-approximation} 中的 $s_n$。由
\eqref{eq:simple-bochner-estimate}，
\[
 \left\|\int s_n\right\|\le\int\|s_n\|,
 \qquad
 \left|\int\|s_n\|-\int\|f\|\right|
 \le\int\|s_n-f\|\longrightarrow0.
\]
令 $n\to\infty$ 得第一个估计。对简单 $s_n$，有限和直接给
$\phi(\int s_n)=\int\phi(s_n)$，而
\[
 \int|\phi(s_n)-\phi(f)|
 \le\|\phi\|\int\|s_n-f\|\longrightarrow0.
\]
$\phi$ 的连续性允许左端同时取极限，得到第二式。
\end{proof}

\begin{theorem}[Bochner DCT]\label{theorem:3-4-2-3}
设 $f_n:X\to B$ 强可测，$f_n(x)\to f(x)$ 在范数意义下 a.e.，且
$\|f_n(x)\|\le g(x)$ a.e.，其中 $g\in L^1(\mu)$ 非负。则 $f$ Bochner 可积，
\[
 \int_X\|f_n-f\|\,\dd\mu\longrightarrow0,
 \qquad
 \int_Xf_n\,\dd\mu\longrightarrow\int_Xf\,\dd\mu
\]
（后一收敛在 $B$ 的范数中）。
\end{theorem}

\begin{proof}
对每个 $n$，强可测性给一个几乎可分值子空间 $B_n$。把所有 $B_n$ 的可数稠密集取并，再取闭
线性包，得到一个可分闭子空间 $B_0$；在共同零集外，所有 $f_n(x)$ 以及其极限 $f(x)$ 都属于
$B_0$。又对每个 $\phi\in B^*$，$\phi(f_n)\to\phi(f)$ a.e.，所以完成化上 $f$ 弱可测。
Pettis 定理给出 $f$ 的强可测性。由 $\|f\|\le g$，它范数可积；同样，假设
$\|f_n\|\le g$ 与 $f_n$ 的强可测性说明每个 $f_n$ 都 Bochner 可积。

现在 $\|f_n-f\|\to0$ a.e. 且 $\|f_n-f\|\le2g$。标量 DCT 给第一项收敛；基本估计于是给
\[
 \left\|\int f_n-\int f\right\|
 \le\int\|f_n-f\|\longrightarrow0.
\]
\end{proof}

\begin{example}[Hilbert 值级数]\label{ex:3-4-2-1}
设 $(e_n)$ 是 Hilbert 空间 $H$ 中的规范正交序列，标量函数 $a_n$ 可测，并且
\[
 G(x):=\left(\sum_{n=1}^{\infty}|a_n(x)|^2\right)^{1/2}\in L^1(\mu).
\]
对几乎每个 $x$，正交性说明部分和
$u_N(x)=\sum_{n\le N}a_n(x)e_n$ 在 $H$ 中 Cauchy；记其极限为
$u(x)=\sum_na_n(x)e_n$，并在余下的可测零集上令 $u=0$；则 $\|u(x)\|=G(x)$ a.e.。

每个 $u_N$ 都强可测，因为把有限个可测坐标同时截断并量化即可得到有限值简单逼近。极限 $u$ 的
强可测性也可以由一条显式对角列看出：把 $u_N$ 的前 $N$ 个复坐标的实部与虚部分别量化到网格
$N^{-1}\Z$；当 $\|u_N(x)\|>N$ 时把量化值置零。所得 $v_N$ 只有有限多个值且可测。在任意
$G(x)<\infty$ 的点，当 $N>G(x)$ 后截断不再发生，而坐标量化误差至多
$\sqrt{2N}/N$。因此
\[
 \|v_N(x)-u(x)\|
 \le\|v_N(x)-u_N(x)\|+\|u_N(x)-u(x)\|\longrightarrow0.
\]
所以 $u$ 强可测，并因 $\|u\|=G$ 而 Bochner 可积。又
$\|u_N-u\|\to0$ a.e. 且 $\|u_N-u\|\le2G$，Bochner DCT 给
\[
 \int_Xu_N\,\dd\mu\longrightarrow\int_Xu\,\dd\mu.
\]
这为 Hilbert 值 Fourier 级数逐项积分提供了一个完全可核查的入口。
\end{example}

\begin{proposition}[Banach 值 $L^p$ 的完备性]\label{proposition:3-4-2-4}
对 $1\le p<\infty$，强可测函数模 a.e. 相等、赋范
\[
 \|f\|_{L^p(B)}=\left(\int_X\|f(x)\|^p\,\dd\mu(x)\right)^{1/p},
\]
构成 Banach 空间。取本质上确界范数时，$L^\infty(X;B)$ 也完备。
\end{proposition}

\begin{proof}
强可测函数对线性组合封闭；标量 Minkowski 不等式应用于点态范数后给三角不等式，而
$\|f\|_{L^p(B)}=0$ 等价于 $f=0$ a.e.。因此只需证明完备性。

设 $(f_n)$ 在 $L^p(B)$ 中 Cauchy，$p<\infty$。抽取子列 $(f_{n_k})$，使
\[
 \sum_{k=1}^{\infty}\|f_{n_{k+1}}-f_{n_k}\|_{L^p(B)}<\infty.
 \tag{3.4.14}\label{eq:banach-lp-fast-subsequence}
\]
令 $h_k(x)=\|f_{n_{k+1}}(x)-f_{n_k}(x)\|$。标量 Minkowski 不等式给
\[
 \left\|\sum_{k=1}^Nh_k\right\|_p
 \le\sum_{k=1}^N\|h_k\|_p.
\]
右端由 \eqref{eq:banach-lp-fast-subsequence} 一致有界。对左端的 $p$ 次方用 MCT，得到
$H:=\sum_kh_k\in L^p$，特别地 $H(x)<\infty$ a.e.。在这些点，$B$ 的完备性使
$f_{n_k}(x)$ 收敛于某个 $f(x)$。

每个 $f_{n_k}$ 几乎取值于一个可分闭子空间；把这些子空间的可数稠密集取并并闭线性张成，得到
共同可分闭子空间 $B_0$，而 $f$ 也几乎取值其中。作为强可测函数的 a.e. 点态极限，或由 Pettis
定理，$f$ 强可测。对每个 $j$，
\[
 \|f_{n_j}(x)-f(x)\|
 \le\sum_{k\ge j}h_k(x).
\]
对右端的有限部分用 Minkowski，再令部分和增大并用 MCT，得到
\[
 \|f_{n_j}-f\|_{L^p(B)}
 \le\sum_{k\ge j}\|h_k\|_p\longrightarrow0.
\]
原列为 Cauchy。给定 $\varepsilon>0$，先取 $N$，使 $m,n\ge N$ 时
$\|f_n-f_m\|_{L^p(B)}<\varepsilon/2$；再选 $j$ 使 $n_j\ge N$ 且
$\|f_{n_j}-f\|_{L^p(B)}<\varepsilon/2$。于是对每个 $n\ge N$，三角不等式给出
$\|f_n-f\|_{L^p(B)}<\varepsilon$。

若 $p=\infty$，抽取子列使
$\|f_{n_{k+1}}-f_{n_k}\|_\infty\le2^{-k}$。去掉可数个例外零集后，增量在每一点一致可和，
从而在完备空间 $B$ 中定义极限 $f$，并且
$\|f_{n_j}-f\|_\infty\le\sum_{k\ge j}2^{-k}$。强可测性仍由共同可分值空间与 Pettis 定理得到，
而
\[
  \|f(x)\|\le \|f_{n_1}(x)\|+\sum_{k\ge1}2^{-k}
  \quad\text{a.e.},
\]
所以 $f\in L^\infty(X;B)$。给定 $\varepsilon>0$，先由原列的 Cauchy 性取 $N$，使
$m,n\ge N$ 时 $\|f_n-f_m\|_\infty<\varepsilon/2$；再选 $j$ 使 $n_j\ge N$ 且
$\|f_{n_j}-f\|_\infty<\varepsilon/2$。于是对每个 $n\ge N$，三角不等式给出
$\|f_n-f\|_\infty<\varepsilon$，故整列收敛于 $f$。
\end{proof}

\begin{remark}[选读边界：Bochner 积分与向量测度]
Pettis 可测性说明何时能由简单函数逐点逼近，Bochner 积分则积分一个已经给定的强可测函数；二者
都不蕴含任意绝对连续的 $B$-值向量测度具有 Bochner 密度。后一个性质取决于值域是否具有
Radon--Nikodym 性质（RNP）。自反空间具有 RNP，而某些 $L^\infty$ 型空间不具有 RNP。因此，
从向量测度恢复密度是一个额外的值域几何问题，不能由 Pettis 可测性或 Bochner 积分的定义推出。
本章不使用上述两个 RNP 分类结论，也不把它们当作任何后续证明的前提；这里只用它们标出一个必须在
向量测度理论中另行定义和证明的边界。
\end{remark}

\begin{figure}[htbp]
\centering
\begin{tikzpicture}[node distance=14mm and 16mm]
  \node[book strong node] (f) {向量场 $f:X\to B$};
  \node[book node,right=of f] (weak) {所有坐标 $\phi\circ f$\\可测};
  \node[book warm node,below=of weak] (sep) {值域几乎落在\\可分子空间 $B_0$};
  \node[book strong node,right=18mm of weak] (dist) {可数距离函数\\$\|f-b_k\|$};
  \node[book strong node,below=of dist] (simp) {有限最近点逼近\\$s_n\to f$};
  \draw[book accent arrow] (f)--node[above,book label]{标量观察}(weak);
  \draw[book accent arrow] (weak)--(dist);
  \draw[book accent arrow] (sep)--(dist);
  \draw[book accent arrow] (dist)--(simp);
  \draw[book dashed arrow] (simp.west)--node[below,book label]{强可测}(f.south east);
\end{tikzpicture}
\caption{Pettis 定理的可数化机制：弱可测给坐标，可分值把不可数距离检验压成可数族。}
\label{fig:3-4-2-pettis}
\end{figure}

Bochner 积分允许把参数化核截面视作 $L^2$ 或一般 Banach 空间中的向量，也允许在可积支配下把
向量极限移过积分号。随机 \Ito{} 积分虽然最终也是 $L^2$ 极限，却还需要适应性与等距构造，因而
不是 Bochner 积分的直接特例。

\begin{exercise}[简单向量积分]\label{exercise:3-4-2-1}
从共同分割出发，补全简单函数积分的表示无关性，并证明
\eqref{eq:simple-bochner-estimate}。
\end{exercise}
\begin{exercise}[连续映射]\label{exercise:3-4-2-2}
不用直接引用 \cref{ex:3-4-2-3}，从可分度量空间的一组可数基构造连续 Banach 值函数的简单
逼近。
\end{exercise}
\begin{exercise}[不可分值反例]\label{exercise:3-4-2-3}
补全 \cref{ex:3-4-2-2} 中“任一可分子空间包含在可数坐标子空间中”的闭包论证。
\end{exercise}
\begin{exercise}[Banach 值 $L^p$]\label{exercise:3-4-2-4}
在 \cref{proposition:3-4-2-4} 的证明中，逐步说明 MCT、Minkowski 与 Pettis 定理各自承担的
作用，并验证极限的 a.e. 等价类不依赖所选快速子列。
\end{exercise}

最后还要回答：若 $f$ 同时依赖两个变量，把一个变量积分掉后所得 Banach 值函数是否强可测？逐个
对偶坐标应用标量 Fubini 只能识别一个已经存在的向量，不能自动制造这个向量。下一小节必须单独
完成截面可测性与简单逼近两项工作。


\subsection{向量值 Fubini 与 Hilbert 值应用}\label{subsec:3-4-3}
\index{Bochner Fubini}\index{Hilbert 空间}\index{Schur 测试}

设 $(X,\mathcal A,\mu)$ 与 $(Y,\mathcal B,\nu)$ 都是 $\sigma$-有限测度空间，并采用
\cref{subsec:3-2-1} 中的乘积测度约定。标量 Fubini 已经告诉我们如何控制
$\|f(x,y)\|$，但它本身没有说明 $y\mapsto f(x,y)$ 是否强可测，也没有说明
$x\mapsto\int_Yf(x,y)\,\dd\nu(y)$ 是否是强可测向量场。这两个问题都要从简单函数逼近中获得。

\begin{definition}[Hilbert 空间与规范正交族]\label{def:3-4-3-hilbert}
内积空间 $H$ 若关于范数 $\|h\|=\sqrt{\langle h,h\rangle}$ 完备，则称为\emph{Hilbert 空间}。
族 $(e_j)$ 若满足 $\langle e_j,e_k\rangle=\delta_{jk}$，称为\emph{规范正交族}；若其闭线性包为
$H$，则称为\emph{规范正交基}。本书在复 Hilbert 空间中约定内积对第一变量线性。
\end{definition}

$\ell^2$ 与本章已经证明完备的 $L^2(\mu)$ 是基本例子。在 $H=\C^m$ 中，Bochner 积分就是逐坐标
积分，正交恒等式也只是有限平方和；无限维新增的义务是证明坐标极限仍落在 $H$ 中，并证明积分可与
该极限交换。

\begin{example}[Hilbert 值指示函数]\label{ex:3-4-3-1}
若 $E\in\mathcal A$、$\mu(E)<\infty$ 且 $h\in H$，则
$f(x)=\1_E(x)h$ 是可积简单函数，
\[
  \int_Xf\,\dd\mu=\mu(E)h,
  \qquad
  \int_X\|f\|\,\dd\mu=\mu(E)\|h\|.
\]
若 $h\ne0$ 且 $\mu(E)=\infty$，则第二个积分为无穷；此时第一个表达式不是一个“无穷倍向量”，
而是根本没有 Bochner 积分。
\end{example}

\begin{example}[轨道平均]\label{ex:3-4-3-2}
设 $(U_t)_{t\in\R}$ 是一族酉算子，即每个 $U_t:H\to H$ 线性且保持内积，并且对每个固定
$\xi\in H$，映射
$t\mapsto U_t\xi$ 在范数中连续。若 $g\in L^1(\R)$，则
\[
  t\longmapsto g(t)U_t\xi
\]
强可测：连续轨道的像具有可分闭线性包，乘以标量后仍落在该可分闭线性包中；对每个
$\phi\in H^*$，标量函数 $g(t)\phi(U_t\xi)$ 可测，故可用 Pettis 定理。其范数等于
$|g(t)|\|\xi\|$，所以
$\int_\R g(t)U_t\xi\,\dd t$ 是一个确定的 Hilbert 向量。若 $V:H\to H$ 有界线性，则
\cref{proposition:3-4-3-2} 将严格证明
\[
 V\!\left(\int_\R g(t)U_t\xi\,\dd t\right)
 =\int_\R g(t)VU_t\xi\,\dd t.
\]
\end{example}

\begin{example}[总质量为一时的 Bochner 平均]\label{ex:3-4-3-3}
在总质量为一的测度空间 $(\Omega,\mathcal F,\Prob)$ 上，强可测 Banach 值函数 $Z$ 若满足
$\int\|Z\|\,\dd\Prob<\infty$，可记
\[
  \operatorname{Av}_{\Prob}(Z):=\int_\Omega Z(\omega)\,\dd\Prob(\omega)
\]
为其 Bochner 平均；在概率论中也称为 Bochner 期望。不可分值域中，单有 Borel 可测不能替代
强可测。从绝对连续向量测度恢复密度则是 RNP 问题，与给定强可测函数的 Bochner 积分是不同问题。
\end{example}

\begin{definition}[乘积空间上的 Bochner 可积]\label{def:3-4-3-1}
若 $f:X\times Y\to B$ 关于 $\mathcal A\otimes\mathcal B$ 强可测，并且
\[
 \int_{X\times Y}\|f(x,y)\|\,\dd(\mu\otimes\nu)(x,y)<\infty,
 \tag{3.4.15}\label{eq:bochner-product-integrable}
\]
则称 $f$ 在乘积空间上 Bochner 可积。
\end{definition}

\begin{definition}[积分算子]\label{def:3-4-3-2}
给定标量核 $K:X\times Y\to\C$ 与函数 $g:Y\to\C$，形式表达式
\[
 (Tg)(x)=\int_YK(x,y)g(y)\,\dd\nu(y)
 \tag{3.4.16}\label{eq:integral-operator-definition}
\]
称为由 $K$ 生成的\emph{积分算子}。该式是否 a.e. 绝对存在、是否定义有界算子，都必须由核估计
证明，不能只靠记号宣告。
\end{definition}

\begin{remark}[弱等式的识别原则]\label{def:3-4-3-3}
若 $u,v\in B$ 满足 $\phi(u)=\phi(v)$ 对每个 $\phi\in B^*$ 成立，则一点 Hahn--Banach
范数化推论给 $u=v$。在 Hilbert 空间中，可以等价地检验
$\langle u,h\rangle=\langle v,h\rangle$ 对每个 $h\in H$ 成立。这个原则只能识别两个已经存在的
向量；逐坐标积分等式本身不能保证强可测性或 Bochner 积分的存在。
\end{remark}

\begin{theorem}[Bochner Fubini]\label{theorem:3-4-3-1}
设 $\mu,\nu$ $\sigma$-有限，且 $f:X\times Y\to B$ 关于
$\mathcal A\otimes\mathcal B$ 强可测，并满足
\eqref{eq:bochner-product-integrable}。则对几乎每个 $x$，截面
$f_x(y)=f(x,y)$ 强可测且 Bochner 可积；在例外处置零后，
\[
 F(x):=\int_Yf(x,y)\,\dd\nu(y)
\]
强可测并属于 $L^1(X;B)$，而且
\[
 \int_{X\times Y}f\,\dd(\mu\otimes\nu)
 =\int_X\left(\int_Yf(x,y)\,\dd\nu(y)\right)\dd\mu(x).
 \tag{3.4.17}\label{eq:bochner-fubini}
\]
交换 $X,Y$ 后的结论与迭代积分等式也成立。
\end{theorem}

\begin{proof}
先说明完成化问题。若强可测性相对于乘积测度的完成化给出，将第 $n$ 个简单逼近的
两两不交值层记为 $E_{n,1},\ldots,E_{n,r_n}$。对每层选取
$\mathcal A\otimes\mathcal B$-可测集 $A_{n,k}$，使 $E_{n,k}\mathbin\triangle A_{n,k}$ 为乘积零集，
再依次以 $A_{n,k}\setminus\bigcup_{j<k}A_{n,j}$ 取代 $A_{n,k}$。原值层两两不交，所以这个去交操作
只再改变一个零集，且得到真正的可测简单函数 $\widetilde u_n$。将所有值层的对称差与
原逼近的收敛例外集都包含在某个共同的
$\mathcal A\otimes\mathcal B$-可测乘积零集 $N$ 中；在 $N^c$ 上有 $\widetilde u_n=u_n\to f$。
令 $\widetilde f=\lim_n\widetilde u_n$ 于 $N^c$，并在 $N$ 上置零，则 $\widetilde f$ 具有乘积可测的简单逼近，且
$\widetilde f=f$ a.e.。标量 Tonelli 应用于 $\1_N$，
说明几乎每个截面只改变一个 $\nu$-零集。因此只需对具有乘积可测简单逼近的代表证明结论。

取简单函数 $u_n\to f$ a.e. 于 $X\times Y$。乘积零集的截面几乎处处为 $\nu$-零集，而每个
$(u_n)_x$ 都是 $Y$ 上的可测简单函数。因此对几乎每个 $x$，$(u_n)_x\to f_x$ a.e.，从而
$f_x$ 强可测。对非负标量函数 $\|f\|$ 使用 Tonelli，由
\eqref{eq:bochner-product-integrable} 得
\[
 r(x):=\int_Y\|f(x,y)\|\,\dd\nu(y)<\infty
\]
对几乎每个 $x$ 成立，并且 $\int_Xr\,\dd\mu=\int_{X\times Y}\|f\|$。所以这些截面
Bochner 可积。

还需证明 $F$ 强可测。由\cref{lemma:3-4-2-bochner-construction} 先取一列 $L^1(B)$ 简单逼近，
再抽取子列并重新编号，可使
\[
 \sum_{n=1}^{\infty}\int_{X\times Y}\|f-s_n\|\,\dd(\mu\otimes\nu)<\infty.
 \tag{3.4.18}\label{eq:bochner-fubini-fast-approximation}
\]
若这些简单函数的值层只在乘积完成化中可测，就按上述逐层取代并去交的方法换成
$\mathcal A\otimes\mathcal B$-可测代表；每个 $s_n$ 只改变于乘积零集，所以上式保持不变。
删去零值层后，把 $s_n$ 写成
$s_n=\sum_{k=1}^{m_n}b_{n,k}\1_{E_{n,k}}$，其中 $b_{n,k}\ne0$；可积性保证每个
$(\mu\otimes\nu)(E_{n,k})<\infty$。在所有截面测度有限的点定义
\[
 S_n(x)=\int_Ys_n(x,y)\,\dd\nu(y)
       =\sum_{k=1}^{m_n}b_{n,k}\nu((E_{n,k})_x),
 \tag{3.4.19}\label{eq:simple-section-integral}
\]
并在例外点置零。标量截面可测性定理说明
$x\mapsto\nu((E_{n,k})_x)$ 可测。有限个可测标量函数乘固定向量之和是强可测的：把这些标量函数
同时截断并按越来越细的网格量化，就得到有限值简单逼近。因此 $S_n$ 强可测。

令
\[
 \rho_n(x)=\int_Y\|f(x,y)-s_n(x,y)\|\,\dd\nu(y).
\]
标量 Fubini 及 \eqref{eq:bochner-fubini-fast-approximation} 给
\[
 \sum_n\int_X\rho_n\,\dd\mu<\infty.
\]
对非负级数使用 Tonelli，得 $\sum_n\rho_n(x)<\infty$ a.e.，特别地 $\rho_n(x)\to0$。在这些点，
Bochner 基本估计给
\[
 \|S_n(x)-F(x)\|\le\rho_n(x)\longrightarrow0.
\]
所有 $S_n$ 的几乎可分值子空间取可数闭线性包仍可分，且每个 $\phi\in B^*$ 都满足
$\phi(S_n)\to\phi(F)$ a.e.；Pettis 定理因此给出 $F$ 的强可测性。并且
\[
 \int_X\|S_n-F\|\,\dd\mu
 \le\int_X\rho_n\,\dd\mu
 =\int_{X\times Y}\|s_n-f\|\,\dd(\mu\otimes\nu)\longrightarrow0.
 \tag{3.4.20}\label{eq:section-integrals-l1-converge}
\]
同时 $\|F(x)\|\le r(x)$，所以 $F\in L^1(X;B)$。

对简单函数 $s_n$，标量 Cavalieri 公式逐个值层给
\[
 \int_XS_n(x)\,\dd\mu(x)
 =\sum_kb_{n,k}\int_X\nu((E_{n,k})_x)\,\dd\mu(x)
 =\sum_kb_{n,k}(\mu\otimes\nu)(E_{n,k})
 =\int_{X\times Y}s_n.
\]
由 \eqref{eq:section-integrals-l1-converge} 与 Bochner 基本估计，左端趋于 $\int_XF$；由
$s_n\to f$ 于 $L^1(B)$，右端趋于 $\int_{X\times Y}f$。这证明
\eqref{eq:bochner-fubini}。转置函数 $\widetilde f(y,x)=f(x,y)$ 关于
$\mathcal B\otimes\mathcal A$ 强可测，且其范数积分仍等于 \eqref{eq:bochner-product-integrable}
中的积分；把已经证明的等式应用于 $(Y,\nu)\times(X,\mu)$，便得到另一迭代次序。
\end{proof}

图 \ref{fig:3-4-3-bochner-fubini} 中两条路径之所以汇合，不是因为逐坐标写出了同样的形式式子，而是
因为可积简单函数在两条路径上严格相等，且 Bochner 基本估计允许把等式传到 $L^1(B)$ 极限。

\begin{figure}[htbp]
\centering
\begin{tikzpicture}[node distance=16mm and 23mm]
  \node[book strong node] (prod) {$f\in L^1(X\times Y;B)$};
  \node[book strong node,right=of prod] (sect) {$F(x)=\int_Yf(x,y)\,\dd\nu$\\$F\in L^1(X;B)$};
  \node[book warm node,below=of prod] (direct) {$\int_{X\times Y}f$};
  \node[book warm node,below=of sect] (iter) {$\int_XF(x)\,\dd\mu$};
  \draw[book accent arrow] (prod)--node[above,book label]{沿 $Y$ 积分}(sect);
  \draw[book accent arrow] (prod)--node[left,book label]{直接积分}(direct);
  \draw[book accent arrow] (sect)--node[right,book label]{沿 $X$ 积分}(iter);
  \draw[book accent arrow] (direct)--node[below,book label]{简单函数逼近}(iter);
\end{tikzpicture}
\caption{Bochner Fubini：标量 Tonelli 控制范数，简单函数逼近负责向量等式与强可测性。}
\label{fig:3-4-3-bochner-fubini}
\end{figure}

一旦积分向量存在，有界线性算子便可以安全地穿过积分号。

\begin{proposition}[有界算子与积分交换]\label{proposition:3-4-3-2}
若 $T:B\to C$ 有界线性，$f:X\to B$ Bochner 可积，则 $Tf$ Bochner 可积，并且
\[
  T\!\left(\int_Xf\,\dd\mu\right)=\int_XTf\,\dd\mu.
\]
\end{proposition}

\begin{proof}
由 $f$ 的强简单逼近以及 $T$ 的连续性，$Tf$ 强可测；又有
$\|Tf(x)\|\le\|T\|\,\|f(x)\|$ a.e.，故 $Tf$ 范数可积。
若 $s_n$ 是 $f$ 的 $L^1(B)$ 简单逼近，则 $Ts_n$ 是简单函数，且
\[
 \int\|Ts_n-Tf\|\le\|T\|\int\|s_n-f\|\longrightarrow0.
\]
对简单函数，等式由有限和及 $T$ 的线性成立；令 $n\to\infty$，
$T$ 的连续性与 Bochner 基本估计分别允许等式两端取极限。
\end{proof}

\begin{proposition}[Hilbert 内积交换与 $L^2$ 配对]\label{proposition:3-4-3-3}
若 $f:X\to H$ Bochner 可积，则对每个 $h\in H$，
\[
 \left\langle\int_Xf\,\dd\mu,h\right\rangle
 =\int_X\langle f(x),h\rangle\,\dd\mu(x).
 \tag{3.4.21}\label{eq:hilbert-inner-product-integral}
\]
若 $f,g\in L^2(X;H)$，则 $x\mapsto\langle f(x),g(x)\rangle$ 可积，且
\[
 \int_X|\langle f(x),g(x)\rangle|\,\dd\mu(x)
 \le\|f\|_{L^2(H)}\|g\|_{L^2(H)}.
 \tag{3.4.22}\label{eq:hilbert-l2-pairing}
\]
\end{proposition}

\begin{proof}
映射 $T_h(v)=\langle v,h\rangle$ 是范数为 $\|h\|$ 的有界线性泛函（本书内积对第一变量
线性），所以 \eqref{eq:hilbert-inner-product-integral} 是
\cref{proposition:3-4-3-2} 的特例。

取有限值可测简单函数 $s_n\to f$、$t_n\to g$ a.e.，并把两个例外零集合并。
每个 $\langle s_n,t_n\rangle$ 是标量简单函数，因而可测；内积的连续性给出
$\langle s_n(x),t_n(x)\rangle\to\langle f(x),g(x)\rangle$ 于该零集外。完成测度空间上的零集修改仍可测，
故 $\langle f,g\rangle$ 可测。逐点 Cauchy--Schwarz 及标量 \Holder{} 不等式给
\[
 \int|\langle f,g\rangle|
 \le\int\|f\|\,\|g\|
 \le\left(\int\|f\|^2\right)^{1/2}
      \left(\int\|g\|^2\right)^{1/2},
\]
这就是 \eqref{eq:hilbert-l2-pairing}。注意在无限测度空间上，$f\in L^2$ 本身并不保证
$\int f$ 存在；合法的对象是这里的 $L^2$ 配对。
\end{proof}

积分算子的第一个实用估计完全不需要 Fourier 变换。它把每一行、每一列的 $L^1$ 控制组合成
$L^2$ 算子范数。

\begin{theorem}[Schur 型积分算子估计]\label{theorem:3-4-3-4}
设 $K:X\times Y\to\C$ 联合可测，并存在 $0\le A,B<\infty$，使
\[
 \int_Y|K(x,y)|\,\dd\nu(y)\le A\quad\text{对几乎处处的 }x,
 \qquad
 \int_X|K(x,y)|\,\dd\mu(x)\le B\quad\text{对几乎处处的 }y.
 \tag{3.4.23}\label{eq:schur-row-column}
\]
则 \eqref{eq:integral-operator-definition} 先在有有限测度支撑的简单函数上定义，并唯一延拓为有界
线性算子
\[
 T:L^2(Y,\nu)\longrightarrow L^2(X,\mu),
 \qquad \|T\|\le\sqrt{AB}.
\]
对每个 $g\in L^2(Y)$，该延拓仍与 a.e. 绝对收敛的核积分
$\int_YK(x,y)g(y)\,\dd\nu(y)$ 一致。
\end{theorem}

\begin{proof}
有有限测度支撑的简单函数在 $L^2(Y)$ 中稠密；这由 $\sigma$-有限耗尽、截断和标量简单函数逼近
得到。先取这样的 $g$。由非负 Tonelli 及列估计，
\[
 \int_X\int_Y|K(x,y)|\,|g(y)|^2\,\dd\nu(y)\dd\mu(x)
 \le B\|g\|_2^2<\infty.
 \tag{3.4.24}\label{eq:schur-column-tonelli}
\]
故对几乎每个 $x$，
$H_g(x):=\int_Y|K(x,y)||g(y)|^2\,\dd\nu(y)<\infty$。在这些点，加权
Cauchy--Schwarz 给
\[
\int_Y|K(x,y)g(y)|\,\dd\nu(y)
 \le\left(\int_Y|K(x,y)|\,\dd\nu(y)\right)^{1/2}H_g(x)^{1/2}<\infty
\]
以及
\[
 |Tg(x)|^2
 \le A\int_Y|K(x,y)||g(y)|^2\,\dd\nu(y).
\]
函数 $(x,y)\mapsto K(x,y)g(y)$ 联合可测；分别对其实部、虚部的正负部作参数积分，Cavalieri--Tonelli
说明这些扩展实值函数关于 $x$ 可测。因而在上述绝对可积集合上取积分、在其补集置零所定义的
$Tg$ 是可测函数。
积分并使用 \eqref{eq:schur-column-tonelli}，得到
\[
 \|Tg\|_2^2\le AB\|g\|_2^2.
 \tag{3.4.25}\label{eq:schur-dense-bound}
\]
因此 $T$ 唯一连续延拓到全部 $L^2(Y)$，且范数不超过 $\sqrt{AB}$。若 $A=0$ 或 $B=0$，
\eqref{eq:schur-dense-bound} 已给出零算子结论；而相应的行积分或列积分为零也说明 $K=0$
于乘积几乎处处，故一般 $g$ 的核积分也几乎处处为零。以下可假设 $A,B>0$。

剩下的是识别延拓的点态公式。固定 $g\in L^2(Y)$，取上述稠密子空间中的 $g_n$，满足
$\|g_n-g\|_2\le2^{-n}$，并令
\[
 H_n(x)=\int_Y|K(x,y)|\,|g_n(y)-g(y)|^2\,\dd\nu(y).
\]
Tonelli 与列界给
\[
 \sum_{n=1}^{\infty}\int_XH_n\,\dd\mu
 \le B\sum_{n=1}^{\infty}4^{-n}<\infty.
\]
因此 $\sum_nH_n(x)<\infty$ a.e.，特别地 $H_n(x)\to0$。同样由
\eqref{eq:schur-column-tonelli} 对一般 $g$ 的版本，$H_g(x)<\infty$ a.e.。行估计与加权
Cauchy--Schwarz 因而同时给出
\[
 \int_Y|K(x,y)g(y)|\,\dd\nu(y)<\infty,
 \qquad
 \int_Y|K(x,y)||g_n(y)-g(y)|\,\dd\nu(y)
 \le\sqrt{A H_n(x)}\longrightarrow0
 \tag{3.4.26}\label{eq:schur-pointwise-kernel-limit}
\]
对几乎每个 $x$ 成立。于是核积分 $T_0g_n(x)\to T_0g(x)$ a.e.
这个几乎处处类与 $g$ 的代表无关：若 $g$ 只在 $\nu$-零集 $N$ 上被修改，则由
Tonelli 与列估计，
\[
 \int_X\int_N|K(x,y)|\,\dd\nu(y)\dd\mu(x)
 =\int_N\int_X|K(x,y)|\,\dd\mu(x)\dd\nu(y)=0,
\]
所以修改前后的核积分对几乎每个 $x$ 相同。

另一方面，延拓的算子范数界给
\[
 \sum_n\|Tg_n-Tg\|_2
 \le\sqrt{AB}\sum_n2^{-n}<\infty.
\]
令 $U_N=\sum_{n=1}^N|Tg_n-Tg|$。标量 Minkowski 不等式给出
\[
 \|U_N\|_2\le\sum_{n=1}^N\|Tg_n-Tg\|_2
 \le\sqrt{AB}\sum_{n=1}^{\infty}2^{-n}.
\]
由于 $U_N^2\uparrow(\sum_n|Tg_n-Tg|)^2$，MCT 给出这一无穷和的平方积分有限，
故 $\sum_n|Tg_n(x)-Tg(x)|<\infty$ a.e.；因此 $Tg_n(x)\to Tg(x)$ a.e.。在稠密子空间上
$Tg_n=T_0g_n$，结合 \eqref{eq:schur-pointwise-kernel-limit} 与极限唯一性，得到
\[
 Tg(x)=T_0g(x)=\int_YK(x,y)g(y)\,\dd\nu(y)
\]
几乎处处成立。
\end{proof}

\begin{remark}[一个可回算的 Schur 常数]
在 $X=Y=\R$ 上取 $K(x,y)=e^{-|x-y|}$。平移 $z=x-y$ 给
\[
 \int_\R|K(x,y)|\,\dd y=\int_\R e^{-|z|}\,\dd z=2,
 \qquad
 \int_\R|K(x,y)|\,\dd x=2.
\]
所以 $\|T\|_{2\to2}\le2$。这个算子也可写成与 $e^{-|\cdot|}$ 的卷积；上面的界只用核的行、列
积分即可得到，不需要 Fourier 变换。
\end{remark}

\begin{corollary}[绝对可和的向量级数可逐项积分]\label{corollary:3-4-3-5}
若 $f_n:X\to B$ 强可测且
\[
 \sum_{n=1}^{\infty}\int_X\|f_n(x)\|\,\dd\mu(x)<\infty,
\]
则 $\sum_nf_n(x)$ 在 $B$ 中 a.e. 绝对收敛于一个 Bochner 可积函数 $f$，右侧级数
$\sum_n\int_Xf_n$ 在 $B$ 中绝对收敛，并且
\[
 \int_X\sum_{n=1}^{\infty}f_n(x)\,\dd\mu(x)
 =\sum_{n=1}^{\infty}\int_Xf_n(x)\,\dd\mu(x).
 \tag{3.4.27}\label{eq:bochner-series-integration}
\]
\end{corollary}

\begin{proof}
标量 Tonelli 给
\[
 \int_X\sum_n\|f_n(x)\|\,\dd\mu
 =\sum_n\int_X\|f_n(x)\|\,\dd\mu<\infty,
\]
所以 $\sum_n\|f_n(x)\|<\infty$ a.e.。$B$ 的完备性给出
$f(x)=\sum_nf_n(x)$，并在余下的可测零集上令 $f=0$。所有 $f_n$ 的几乎可分值子空间取可数闭线性包，仍为可分；部分和强可测，
故其点态极限由 Pettis 定理强可测。又
$\|f\|\le\sum_n\|f_n\|\in L^1$，所以 $f$ Bochner 可积。

令 $S_N=\sum_{n\le N}f_n$。则
\[
 \int_X\|f-S_N\|\,\dd\mu
 \le\sum_{n>N}\int_X\|f_n\|\,\dd\mu\longrightarrow0.
\]
Bochner 基本估计给 $\int S_N\to\int f$，而积分线性给
$\int S_N=\sum_{n\le N}\int f_n$。此外
$\sum_n\|\int f_n\|\le\sum_n\int\|f_n\|<\infty$，故右侧在 $B$ 中绝对收敛，并得到
\eqref{eq:bochner-series-integration}。
\end{proof}

取 $f_n(x)=a_n(x)e_n$，\cref{corollary:3-4-3-5} 便回收 Hilbert 值级数逐项积分。这里交换的是
由范数绝对可和保证的级数；右侧在 Banach 空间中绝对收敛，因此任意重排仍有相同的和。一般的
正交平方可和级数若要在积分号下交换，还需另行验证相应
$L^1(H)$ 或 $L^2(H)$ 条件。Schur 条件与 $K\in L^2(X\times Y)$ 的 Hilbert--Schmidt 条件也
不同：前者控制核的行、列绝对积分，后者控制核在乘积空间上的平方积分。

\begin{exercise}[Bochner Fubini 的简单层]\label{exercise:3-4-3-1}
设 $b\in B$、$E\in\mathcal A\otimes\mathcal B$，且 $b=0$ 或
$(\mu\otimes\nu)(E)<\infty$。对 $s=b\1_E$ 直接证明 \eqref{eq:bochner-fubini}，
再推广到可积有限值简单函数；明确指出标量
Cavalieri 公式出现的位置。
\end{exercise}
\begin{exercise}[算子交换]\label{exercise:3-4-3-2}
证明若 $T_n:B\to C$ 在算子范数中趋于 $T$，且 $f$ Bochner 可积，则
$\int T_nf\to\int Tf$，并给出定量误差界。
\end{exercise}
\begin{exercise}[Hilbert 配对]\label{exercise:3-4-3-3}
设 $f\in L^2(X;H)$、$h\in H$ 且 $E\in\mathcal A$、$\mu(E)<\infty$。证明
$\int_Ef$ 存在，并估计 $\|\int_Ef\|$；再用 \eqref{eq:hilbert-inner-product-integral} 计算其与
$h$ 的内积。
\end{exercise}
\begin{exercise}[Schur 测试]\label{exercise:3-4-3-4}
对 $K(x,y)=e^{-a|x-y|}$（$a>0$）计算 Schur 上界。再构造一个满足 Schur 条件但不属于
$L^2(\R^2)$ 的平移不变核，以说明 Schur 条件与 Hilbert--Schmidt 条件不同。
\end{exercise}

至此，本章从简单函数的有限和出发，经过单调极限、乘积积分与 Radon--Nikodym 表示，抵达
Banach 值积分和向量 Fubini。下一章将把积分与平移、缩放以及平均化组合起来，研究卷积、逼近恒等族
与最大函数；那里反复出现的核积分如今都有了严格的存在性与换序依据。

